% -----------------------------------------------------------
\section{1835}
% -----------------------------------------------------------
	\label{EMS-1835}
	
	% ----------------------
	\subsection{Joseph Smith}
	% ----------------------
		\label{EMS-1835-JS}

		% -----
		\subsubsection{Introduction to 1835}
		% -----
		
			sdf
	
		% -----
		\subsubsection{Minutes, 18 January 1835}
		% -----
			\label{EMS-1835-JS-18-JAN-1835}
			% \label{EMS-1830-JS-DAY-3LETTERMONTH-YEAR}
			
			\begin{description}
					\item[Print Sources]
				\end{description}

					\begin{tabular}{p{0.5in}p{1in}p{0.5in}p{1in}}
					JSP	 	& D4:215--218		& JSR 		& --- \\
					DC	 	& --- 				& HC 		& --- \\
					HDDC 	& ---				& TCHC 		& --- \\
					\end{tabular}
					% HDDC = woodford's DC diss; JSR = Marquardt's JS Revelations; TCHC = Vogel HC
				
				\begin{description}
					\item[Online Access] \href{https://www.josephsmithpapers.org/paper-summary/minutes-18-january-1835/1}{Here}
					% \item[Analysis] \hyperref[XX]{Here}
				\end{description}
				
				\begin{description}
					\item[Background]
				\end{description}
				
					% It might also be useful to give a two sentence context of the period: e.g., "This speech was given in Nauvoo to a small congregation one month prior to the destruction of the Nauvoo Expositor. These and other tensions may have been on the prophet's mind when he gave the speech."
				
				\begin{description}
					\item[Original Source]
				\end{description}
				
					% brief description of original source material, with reference to JSP or somewhere else for more info
					% Background material: it might be helpful to have a note that says whether a given document was presented as a speech, was written in a journal, etc.
					% Also a comment here about how reliable the source might be, or what shape it is in, if there are large omissions or parts that are hard to understand, and so on.
				
				\begin{description}
					\item[Text]
				\end{description}
				
					 ...President Rigdon occupieed considerable time. and laid before the council at much length, the importance of assisting the church in this place, to move forward the great work of preparing a place in which the Elders might be endowed, and of printing and sending out the word of the Lord. President Smith took up the subject still further, and occupied a long time in presenting the importance of the same thing and gave decision that Elder [John] Tanner assist with his might, to build up the cause by tarrying in Kirtland...
					 
		% -----
		\subsubsection{Minutes, Discourse, and Blessings, 14--15 February 1835}
		% -----
			\label{EMS-1835-JS-14-15-FEB-1835}
			% \label{EMS-1830-JS-DAY-3LETTERMONTH-YEAR}
			
			\begin{description}
					\item[Print Sources]
				\end{description}

					\begin{tabular}{p{0.5in}p{1in}p{0.5in}p{1in}}
					JSP	 	& D4:219--234		& JSR 		& --- \\
					DC	 	& --- 				& HC 		& 2:180--191 \\
					HDDC 	& ---				& TCHC 		& 2:193--199 \\
					\end{tabular}
					% HDDC = woodford's DC diss; JSR = Marquardt's JS Revelations; TCHC = Vogel HC
				
				\begin{description}
					\item[Online Access] \href{https://www.josephsmithpapers.org/paper-summary/minutes-discourse-and-blessings-14-15-february-1835/1}{Here}
					% \item[Analysis] \hyperref[XX]{Here}
				\end{description}
				
				\begin{description}
					\item[Background]
				\end{description}
				
					% It might also be useful to give a two sentence context of the period: e.g., "This speech was given in Nauvoo to a small congregation one month prior to the destruction of the Nauvoo Expositor. These and other tensions may have been on the prophet's mind when he gave the speech."
				
				\begin{description}
					\item[Original Source]
				\end{description}
				
					% brief description of original source material, with reference to JSP or somewhere else for more info
					% Background material: it might be helpful to have a note that says whether a given document was presented as a speech, was written in a journal, etc.
					% Also a comment here about how reliable the source might be, or what shape it is in, if there are large omissions or parts that are hard to understand, and so on.
				
				\begin{description}
					\item[Text]
				\end{description}
				
					\begin{tightcenter}
							\hypertarget{EMS-1835-JS-14-15-FEB-1835-a}%
						Ordination Blessings
							\marginnote{a}
					\end{tightcenter}

					Kirtland February 14 1835

					This day a meeting was called of those who journeyed to Zion for the <purpose of laying> the foundation of its redemption last season with as many more of the Brethren \& Sisters as felt disposed to attend. President Joseph Smith \supersu{Junr} presided over the meeting. After the Congregation assembled. he arose and requested the attention and read the 15\supersu{th} chapter of John and said, let us endeavor to solemnize our minds that we may receive a blessing by calling on the Lord \& said let us pray

					After an appropriate and affecting prayer was made the Bretheren who went to Zion, were requested to take their seats together in one part of the house by themselves. President Smith arose and stated the reason why this meeting was called. It was this. God had commanded it and it was made known to him by vision and by the Holy Spirit. he then gave a relation of some of the circumstances attending us while journeying to Zion, our trials, sufferings \&c \&c. He said God had not designed all this for nothing, but he had it in remembrance yet, and those who went to Zion, with a determination to lay down their lives, if necessary, it was the Will of God, that they should be ordained to the ministry and go forth to prune the vineyard for the last time, or the coming of the Lord, which was nigh, even fifty six years, should wind up the scene. He said many things such as the weak things, even the smallest and weakest among us shall be powerful and mighty and great things, should be accomplished by you, even from this hour. He said, you shall begin to feel the whisperings of the Spirit of God, and the works of God, shall begin to break forth from this time. you shall be endowed with power from on high.

						\hypertarget{EMS-1835-JS-14-15-FEB-1835-b}%
					President Smith then called
						\marginnote{b}%
					upon all those who went to Zion to know if they were agreed with him in the statements which he had made, he said if you are arise upon your feet, all arose upon their feet. He then called upon the balance of the congregation to know also, if they would sanction the movement. They all raised the right hand. The names of those who went to Zion are as follows...

					A Hymn was then sung for dismission (viz.) Hark listen to the trumpeters they call for volunteers. President Hyrum Smith closed by prayer and dismissed for one hour.

					Met pursuant to adjourment \& president Hyrum Smith opened the meeting by Prayer. President J. Smith \supers{Junr.} \sout{and} arose and said. The first business of the meeting was for the three witnesses of the Book of Mormon, to pray each one and then proceed to choose twelve men from the church as Apostles to go to all nations, kindred toungs and people. The three Witnesses united in prayer (Viz) Oliver Cowdery, David Whitmer \& Martin Harris. These three Witnesses were then blessed by the laying on of the hands of the Presidincy. They then according to a former commandment, proceeded to make choice of the twelve: The names \sout{of} are as follows

					\begin{tabular}{p{5cm}p{1cm}}
					Lyman Johnson & 1 \\
					Brigham Young & 2 \\
					H. C. Kimball & 3 \\
					Orson Hyde & 4 \\
					David W. Patten & 5 \\
					Luke Johnson & 6 \\
					Wm. E. McLel[l]in & 7 \\
					John F. Boynton & 8 \\
					Orson Pratt & 9 \\
					Wm. Smith & 10 \\
					Thomas B. Marsh & 11 \\
					Parley P. Pratt & 12
					\end{tabular}
					
						\hypertarget{EMS-1835-JS-14-15-FEB-1835-c}%
					Lyman Johnson, Brigham Young and
						\marginnote{c}%
					H. C. Kimball came forward and the three witnesses laid their hands upon each ones head and prayed separately. The Blessing of L[yman] Johnson was in the name of Jesus Christ, that he should bear the tidings of salvation to nations, tongues and people, until the utmost corners of the earth shall hear the tidings, and that he shall be a witness of the things of God to nations \& tongues, and that Holy Angels shall administer to him occasionally and that no power of the enemy shall prevent him from going forth and doing the work of the Lord. And that he \sout{shall} <should> live until the gathering was accomplished, according to the Holy Prophets. And that he should be like unto Enoch. And your faith shall be like unto his, and he shall be called great among all the living, and Satan shall tremble before thee, and that he shall see the Saviour come and stand on the Earth with power and great glory.

						\hypertarget{EMS-1835-JS-14-15-FEB-1835-d}%
					The blessing of
						\marginnote{d}%
					Brigham Young was as follows, That he should be strong in Body, that he may go forth and gather the Elect preparatory to the great day of the coming of the Lord, and that he might be strong and mighty declaring the tidings to nations that know not God, that he may add ten talents, that he may come to Zion with many sheaves. He shall go forth from land to land and from sea to sea. and shall behold Heavenly Messengers going forth, \& his life shall be prolonged, and the Holy Priesthood is confirmed upon thee, that he may do wonders in the name of Jesus, that he may cast out Devils, heal the sick, raise the dead, open the eyes of the blind, go forth from land to land and from sea to sea, and that heathen nations shall \sout{not} even call him God himself, if he did not rebuke them.

						\hypertarget{EMS-1835-JS-14-15-FEB-1835-e}%
					Heber C. Kimball, blessing
						\marginnote{e}%
					was in substance as follows. That he shall be made like unto those who have <been> blessed before him and be favored with the same blessing. That he might receive visions, the ministration of angels and hear their voices \& even come into the presence of God. That many millions may be converted by his instrumentality that angels may waft him from place to place and that he may stand unto the coming of our Lord and receive a crown over the kingdom of our Lord, that he be made acquainted with the day when Christ shall come, that he shall be made perfect in faith and that the deaf shall hear, the lame shall walk, the blind shall see, and greater things <than> these shall he do, that he shall have boldness of speech before the nations and great power \&c.

						\hypertarget{EMS-1835-JS-14-15-FEB-1835-f}%
					A hymn was then sung
						\marginnote{f}%
					Glorious things of the[e] are spoken. The congregation was then dismissed by J. Smith Jun\supersu{r}\supers{.}

					Sunday, Feb. 15\supersu{th} at 9 \supers{o} clock A.M. the congregation again assembled \& President [Oliver] Cowdery arose, and made some observations upon the nature of the meeting and then prayed unto the Lord for his assistance \&c \&c.

					After which a number of certificates were read and accepted from brethren, that recently returned from Zion. Then, President Cowdery called forwarward Orson Hyde. David W. Patten and Luke Johnson and proceeded to their ordination \& blessing.

						\hypertarget{EMS-1835-JS-14-15-FEB-1835-g}%
					Orson Hyde, blessing.
						\marginnote{g}%
					Oliver Cowdery proceeded and called upon the Lord to smile upon him and that his faith shall be perfect, and that the blessings pronounced shall be realized. He shall be made mighty and be endued with power from on high, and go forth to the nations of the earth to proclaim the gospel. That he shall escape all the pollutions of the world. The Angels shall uphold him, and that he shall go forth according to the commandment, both to Jew \& Gentile and shall go to all nations, kingdoms and tongues \sout{and shall} All who hear his voice shall acknowledge him to be a servant of God. He shall be equal in holding the keys of the kingdom. He shall stand on the earth and bring souls till Christ comes. We know that he loves thee, and may this thy servant be able to walk through pestilence and not be harmed. The powers of darkness shall have no ascendency over him. He shall have power to smite the earth with pestilence, to divide waters and lead through the Saints. He shall go from land to land and from sea to sea. He shall be like unto one of the three Nephites

						\hypertarget{EMS-1835-JS-14-15-FEB-1835-h}%
					David W. Patten's blessing.
						\marginnote{h}%
					O God, give this, thy servant a knowledge of thy will; May he be like one of old who bore testimony of Jesus. May he be a new man from this day forth. He shall be equal with his brethren, the twelve, and have all the qualifications of the Prophets before him. May his body be strong and never be weary. May he walk and not faint. May he have power over all diseases, and faith according to his desires. May the Heavens be opened upon him speedily, that he may bear testimony from knowledge, that he may go to nations and isles afar off. May he have a knowledge of \sout{thing} the things of the kingdom from the beginning and be able to tear down priest-craft like a Lion May he have power to smite his enemies before him with utter destruction. May he continue till the Lord comes. O Father we seal these blessings upon him, even so Amen.
					
					\begin{tightcenter}
							\hypertarget{EMS-1835-JS-14-15-FEB-1835-i}%
						Luke Johnson'\supers{s} blessing
							\marginnote{i}
					\end{tightcenter}

					Our Father, in Heaven, look down in mercy upon us and upon this thy servant whom we ordain to the ministry of the twelve. He shall be prepared and preserved and become like those we have blessed before him, The nations shall tremble before him. He shall hear the voice of God: he shall comfort the hearts of the saints, always. The angels shall bear him up till he shall finish his ministry. He shall be delivered and come forth with Israel. He shall bear testimony to the kings of the earth, and hold communion with the Father, with the son and with the general assembly and church of the first born. If cast into prison he shall be able to comfort the hearts of his comrades His tongue shall be loosed and he shall have power to lead many to Zion, and set down with them, and the Ancient of Days shall pronounce this blessing. That you have been faithful. He shall have strength wisdom \& power. He shall go among the covenant people and speak all their tongues where he shall go. All these blessings we confirm upon him in the name of Jesus. Amen.
					
					\begin{tightcenter}
							\hypertarget{EMS-1835-JS-14-15-FEB-1835-j}%
						Wm. E. M\supers{c}Lel[l]in's blessing.
							\marginnote{j}
					\end{tightcenter}

					In the name of the Lord, Wisdom \& intelligence shall be poured out upon him, to enable him to perform the great work, that is incumbent upon him. That he may be spared until the saints are gathered, that he may stand before Kings and Rulers to bear testimony and be upheld by holy Angels, and the nations of the earth shall acknowledge that God has sent him. He shall have power to overcome his enemies, and his life shall be spared in the midst of pestilance and destruction, and in the midst of his enemies. He shall be a prince and a saviour to God's people. The Tempter shall not overcome him, nor his enemies prevail against him. The Heavens shall be opened unto him as unto men in days of old. He shall be mighty in the hands of God, and shall convince thousands that God has sent him, and his days may be prolonged until the coming of the son of man. He shall be wafted as on eagles wings, from country to country and from people to people and <be able> to <do> wonders in the midst of this generation, Even so Amen.
					
					\begin{tightcenter}
							\hypertarget{EMS-1835-JS-14-15-FEB-1835-k}%
						John F. Boynton's blessing
							\marginnote{k}
					\end{tightcenter}

					Thou hast prevailed, and thou shalt prevail and thou shalt declare the gospel unto many nations. Thou shalt be made mighty before God. And although thou shalt be cast out from the face of men, yet thou shalt <have> power to prevail, thou shalt lead the Elect triumphantly to the places of refuge. Thou shalt be like thy brethren who have been blessed before thee. Thou shalt stand in that day of calamity when the wicked shall be consumed, and present unto the Father, spotless, the fruit of \sout{his} <thy> labor. Thou shalt overcome all the evils that are in the world. Thou shalt have wisdom to put to silence all the wisdom of the wise, and thou shalt see the face of thy Redeemer in the flesh. These blessings \sout{were} <are> pronounced and sealed upon thee even so Amen.
					
					\begin{tightcenter}
							\hypertarget{EMS-1835-JS-14-15-FEB-1835-l}%
						Williams Smith's blessing
							\marginnote{l}
					\end{tightcenter}

					That he may be purified in heart, that he may have communion with God. That he may be equal with his brethren in holding the keys of this ministry That he may be kept and be instrumental in leading Israel forth, that he may be delivered from the hands of those who seek to destroy him: that he may be enabled to bear testimony to the nations, that Jesus lives. That he may stand in the midst of pestilence and destruction, he shall be mighty in the hands of God, in bringing about the restoration of Israel. The nations shall rejoice at the greatness of the gifts which God has bestowed upon him, That his tongue shall be loosed, he shall have power to do great things in the name of Jesus. He shall be preserved and remain on the earth, until Christ shall come to take vengeance on the wicked. Confirmed
					
		% -----
		\subsubsection{Preface to Doctrine and Covenants, 17 February 1835}
		% -----
			\label{EMS-1835-JS-17-FEB-1835}
			% \label{EMS-1830-JS-DAY-3LETTERMONTH-YEAR}
			
			\begin{description}
					\item[Print Sources]
				\end{description}

					\begin{tabular}{p{0.5in}p{1in}p{0.5in}p{1in}}
					JSP	 	& D4:234--237		& JSR 		& --- \\
					DC	 	& --- 				& HC 		& --- \\
					HDDC 	& ---				& TCHC 		& --- \\
					\end{tabular}
					% HDDC = woodford's DC diss; JSR = Marquardt's JS Revelations; TCHC = Vogel HC
				
				\begin{description}
					\item[Online Access] \href{https://www.josephsmithpapers.org/paper-summary/preface-to-doctrine-and-covenants-17-february-1835/1}{Here}
					% \item[Analysis] \hyperref[XX]{Here}
				\end{description}
				
				\begin{description}
					\item[Background]
				\end{description}
				
					% It might also be useful to give a two sentence context of the period: e.g., "This speech was given in Nauvoo to a small congregation one month prior to the destruction of the Nauvoo Expositor. These and other tensions may have been on the prophet's mind when he gave the speech."
				
				\begin{description}
					\item[Original Source]
				\end{description}
				
					% brief description of original source material, with reference to JSP or somewhere else for more info
					% Background material: it might be helpful to have a note that says whether a given document was presented as a speech, was written in a journal, etc.
					% Also a comment here about how reliable the source might be, or what shape it is in, if there are large omissions or parts that are hard to understand, and so on.
				
				\begin{description}
					\item[Text]
				\end{description}
				
					\begin{tightcenter}
							\hypertarget{EMS-1835-JS-17-FEB-1835-a}%
						PREFACE.
							\marginnote{a}
					\end{tightcenter}

					\textit{To the members of the church of the Latter Day Saints---}

					\textsc{Dear Brethren:}

					We deem it to be unnecessary te entertain you with a lengthy preface to the following volume, but merely to say, that it contains in short, the leading items of the religion which we have professed to believe.

					The first part of the book will be found to contain a series of Lectures as delivered before a Theological class in this place, and in consequence of their embracing the important doctrine of salvation, we have arranged them into the following work.

					The second part contains items or principles for the regulation of the church, as taken from the revelations which have been given since its organization, as well as from former ones.

					There may be an aversion in the minds of some against receiving any thing purporting to be articles of religious faith, in consequence of there being so many now extant; but if men believe a system, and profess that it was given by inspiration, certainly, the more intelligibly they can present it, the better. It does not make a principle untrue to \textit{print} it, neither does it make it true not to print it.

						\hypertarget{EMS-1835-JS-17-FEB-1835-b}%
					The church viewing this
						\marginnote{b}%
					subject to be of importance, appointed, through their servants and delegates the High Council, your servants to select and compile this work. Several reasons might be adduced in favor of this move of the Council, but we only add a few words. They knew that the church was evil spoken of in many places---its faith and belief misrepresented, and the way of truth thus subverted. By some it was represented as disbelieving the bible, by others as being an enemy to all good order and uprightness, and by others as being injurious to the peace of all governments civil and political.

					We have, therefore, endeavored to present, though in few words, \textit{our} belief, and when we say this, humbly trust, the faith and principles of this society as a body.

					We do not present this little volume with any other expectation than that we are to be called to answer to every principle advanced, in that day when the secrets of all hearts will be revealed, and the reward of every man's labor be given him.

					With sentiments of esteem

					and sincere respect, we subscribe

					ourselves your brethren in the bonds of

					the gospel of our Lord Jesus Christ.

					\begin{tightright}
						JOSEPH SMITH jr.

						OLIVER COWDERY.

						SIDNEY RIGDON.

						F[rederick] G. WILLIAMS.
					\end{tightright}

					\textit{Kirtland, Ohio, February} 17, 1835.
					
		% -----
		\subsubsection{Minutes and Blessings, 21 February 1835}
		% -----
			\label{EMS-1835-JS-21-FEB-1835}
			% \label{EMS-1830-JS-DAY-3LETTERMONTH-YEAR}
			
			\begin{description}
					\item[Print Sources]
				\end{description}

					\begin{tabular}{p{0.5in}p{1in}p{0.5in}p{1in}}
					JSP	 	& D4:237--247		& JSR 		& --- \\
					DC	 	& --- 				& HC 		& 2:191--198 \\
					HDDC 	& ---				& TCHC 		& 2:199--205 \\
					\end{tabular}
					% HDDC = woodford's DC diss; JSR = Marquardt's JS Revelations; TCHC = Vogel HC
				
				\begin{description}
					\item[Online Access] \href{https://www.josephsmithpapers.org/paper-summary/minutes-and-blessings-21-february-1835/1}{Here}
					% \item[Analysis] \hyperref[XX]{Here}
				\end{description}
				
				\begin{description}
					\item[Background]
				\end{description}
				
					% It might also be useful to give a two sentence context of the period: e.g., "This speech was given in Nauvoo to a small congregation one month prior to the destruction of the Nauvoo Expositor. These and other tensions may have been on the prophet's mind when he gave the speech."
				
				\begin{description}
					\item[Original Source]
				\end{description}
				
					% brief description of original source material, with reference to JSP or somewhere else for more info
					% Background material: it might be helpful to have a note that says whether a given document was presented as a speech, was written in a journal, etc.
					% Also a comment here about how reliable the source might be, or what shape it is in, if there are large omissions or parts that are hard to understand, and so on.
				
				\begin{description}
					\item[Text]
				\end{description}
				
					\begin{tightcenter}
							\hypertarget{EMS-1835-JS-21-FEB-1835-a}%
						Kirtland February 21\supersu{st} 1835
							\marginnote{a}
					\end{tightcenter}

					Pursuant to adjournment a meeting of the Church of the Latter Day Saints was held and proceeded as follows. (Viz.) Meeting opened by prayer, of President David Whitmer, President O[liver] Cowdery addressed the congregation briefly and then Elder Parley P. Pratt was called to the stand and ordained as one of the Twelve. by J. Smith \supers{Jur.} D. Whitmer \& O. Cowdery

					Parley P. Pratt'\supers{s} blessing, as pronounced upon him by O. Cowdery, O. Lord smile from Heaven upon this thy servant, forgive his sins, sanctify his heart and prepare him to receive the blessing. Increase his love for thee and for thy cause, increase his intelligence, communicate to him all that wisdiom, that prudence, and that understanding which he needs as a minister of Righteousness, and to magnify the Apostleship whereunto he is called. May a double portion of that spirit which was communicated to the Diciples of of our Lord \& Saviour, to lead them into all truth, rest down upon him, go with him where <he> goes. that nothing shall prevail against him. that he may be delivered from prisons, from the power of his enemies, and from the adversary of all righteousness. May he be able to mount up on wings as an Eagle's, to run and not be weary, to walk and not faint.
						\hypertarget{EMS-1835-JS-21-FEB-1835-b}%
					May he have
						\marginnote{b}%
					great wisdom \& intelligence, and be able to lead thine elect through this thorny maze. Let sickness and death have no power over him. Let him be equal with his brethren in bringing many sons and daughters to glory, and many nations to a knowledge of the truth. Great blessings shall rest upon thee. Thy faith shall increase, Thou shalt have great power, to prevail, The viel of the Heavens shall be rolled up, Thou shalt be permitted to gaze within it, and receive instruction from on high. No arm that is formed and lifted against thee, shall prosper, no power shall prevail, for thou shalt have power with God. and shall proclaim his gospel, thou wilt be afflicted, but thou shalt be delivered and conquer all thy foes. Thine office shall never be taken from thee. Thou shalt be called great. Angels shall carry thee from place to place. Thy sins are forgiven, and thy name written in the lambs book of life. Even so Amen.

						\hypertarget{EMS-1835-JS-21-FEB-1835-c}%
					The following charge
						\marginnote{c}%
					was then given Elder P. P. Pratt by President O. Cowdery.

					I am aware, Dear Brother, that the mind naturally claims something new: but the same thing rehearsed, frequently profits us. You will have the same difficulties to encounter in fulfilling this ministry, that the ancient Apostles had. You have enlisted in a cause that requires your whole attention, you ought therefore to count the cost, and become a polished shaft, to become a polished shaft, you must be sensible, requires the labor of years; and your station requires a perfect polish. It is required, not merely to travel a few miles in this country, but in distant countries. You must endure much labor, much toil, and many privations to become perfectly polished. Your calling is not like that of the husbandsman, to cultivate a stinted portion of the planet on which we dwell, and when heaven has given the former and the latter rain, and mellow autumn refined his fruit, gather it in and, congratulates himself for a season in the remission of his toils, while he anticipates his winter evenings of relaxation and fireside enjoyment, But. Dear Brother, it is far otherwise with you. Your labor must be incessant and your toil great. You most go forth and labor till the great work is done. It will require a series of years to accomplish it. But you will have this pleasing consolation that, your Heavenly Father requires it; the field is his, the work is his, and he will not only cheer you, animate you \& buoy you up in your pilgrimage, in yours arduous toils. but, when your work is done, and your labor o[v]er, he will take you to himself.
						\hypertarget{EMS-1835-JS-21-FEB-1835-d}%
					But
						\marginnote{d}%
					before this consummation of your felicity. bring your mind to bear upon what will be imperiously required of you to accomplish the great work that lies before you: Count well the cost. You have read of the trials and persecutions of ancient days. Has not bitter experience taught you that they are the same now? You will be dragged before the Authorities, for the religion you profess, and it were better not to set out than to start, look back, or shrink when dangers thicken upon you, or appalling death stares you in the face. I have spoken these things, Dear Brother, because I have seen them in visions. There are strong dungeons and gloomy prisons for you. These should not appall you. You must be called good or bad man. The ancients passed through the same. They had this testimony, that they had seen the Saviour after he rose from the dead. You must bear some testimony, or your mission, your labor, your toil will be in vain. You must bear the same testimony, that there is but one God and one Mediator; he that has seen him will know, him and testify of him. Beware of Pride, beware of evil; shun the very appearance of it, for the time is coming when if you do not hear to <these> things you will have a fall.
						\hypertarget{EMS-1835-JS-21-FEB-1835-e}%
					Among
						\marginnote{e}%
					your many afflictions you will have many blessings also. but you must pass through many afflictions in order to receive the glory that is in reserve for you. You will see thousands who when they first see you will know nothing about salvation by Jesus Christ, You shall see a nation born in a day. A great work lies before you \& the time is near, when you must bid farewell to your native land, cross the mighty deep and sound the tocsin of alarm to other nations, kindreds tongues and people. Remember that all your hopes of deliverance from danger \& from death, will rest upon your faithfulness to God in his cause. You must necessarily serve him with a perfect heart and a willing mind. Avoid strife and vain glory, think not yourself better than your brethren, but pray for them, as well as for yourself, and if you are faithful, great will be your blessings, but if you are not your stewardship will be taken from you, and another appointed in your stead.

					Elder Pratt gave his hand to President O. Cowdery and said he had received ordination and should fulfil the ministry according to the grace given him. To which the President replied, Go forth and Angels shall bear thee up and thou shalt come forth at the last day bringing many with thee.

						\hypertarget{EMS-1835-JS-21-FEB-1835-f}%
					Thomas
						\marginnote{f}%
					B. Marsh'\supers{s} Blessing by O. Cowdery Dear Brother, You are to be a minister of righteousness and to this ministry and Apostleship you are now to be ordained: and May all temporal and spiritual blessings attend you. Your sins are forgiven you, and you are to go forth and preach the everlasting Gospel. You shall travel from kingdom to kingdom and from nation to nation. Angels shall bear the[e] up, and thou shalt be instrumental in bringing thousands of the redeemed of the Lord to Zion. President David Whitmer sealed the above blessing upon him, even so Amen.
					
					\begin{tightcenter}
							\hypertarget{EMS-1835-JS-21-FEB-1835-g}%
						Orson Pratt's blessing.
							\marginnote{g}
					\end{tightcenter}

					Dear Brother, You are chosen, and set apart to be ordained to this Apostleship and this ministry. You shall go forth and preach the gospel and do a mighty work. You shall be sustained. The Holy spirit shall enlighten thy mind. Thou shalt travel from nation to nation. The Lord God shall preserve thee and return thee safe, with songs of everlasting joy upon thy head. The above was confirmed, by President David Whitmer
					
					\phantom{.}\\

						\hypertarget{EMS-1835-JS-21-FEB-1835-h}%
					The following Charge
						\marginnote{h}%
					was given to the Twelve by President O. Cowdery.

					Dear Brethren, previously to delivering the charge I shall read a part of a revelation. It is known to you that previous to the organizing of this church in 1830 The Lord gave revelations or the Church could not have been organized. The people of this Church were weak in faith compared with the ancients. Those who embarked in this cause, were desirous to know how the work was to be conducted. They read many things in the Book of Mormon, concerning their duty and the way the great work ought to be done; but the minds of men are so constructed, that they will not believe without a testimony of seeing or hearing. The Lord gave us a revelation, that in process of time there should be Twelve chosen to preach his gospel to Jew \& Gentile. Our minds have been on a constant stretch to find who these Twelve were.

						\hypertarget{EMS-1835-JS-21-FEB-1835-i}%
					When the time should come,
						\marginnote{i}%
					we could not tell, but we sought the Lord by fasting and prayer, to have our lives prolonged to see this day, to see you, and to take a retrospect of the difficulties through which we have passed. but, having seen the day, it becomes my duty to deliver to you a charge. and first, a few remarks respecting your ministry. You have many Revelations put into your hands. Revelations to make you acquainted with the nature of your Mission. You will have difficulties by reason of your visiting all the nations of the world. You will need wisdom in a tenfold proportion to what you have ever had. You will have to combat all the prejudices of all nations. He then read Revelation, and proceeded to say, Have you desired this ministry with all your hearts, If you have desired it, you are called of God not of man, to go into all the world. He read again, from the Revelation, what the Lord said to the 12 Brethren, you have your duty presented in this revelation. You have been ordained to the Holy Priesthood, You have received it from those who had their power and authority from an angel.--- You are to preach the gospel to every nation. should you in the least degree, come short of your duty, great will be your condemnation. for the greater the calling, the greater the transgression.
						\hypertarget{EMS-1835-JS-21-FEB-1835-j}%
					I,
						\marginnote{j}%
					therefore, warn you to cultivate great humility, for I know the pride of the human heart. Beware, lest the flatterers of the world lift you up. Beware lest your affections are captivated by worldly objects. Let your ministry be first. Remember the souls of men are committed to your charge, and if you mind your calling you shall always prosper. You have been indebted to other men in the first instance for evidence, on that you have acted. But <it> is necessary that you receive a testimony from Heaven for yourselves, so that you can bear testimony to the truth of the Book of Mormon. And that you have seen the face of God: that is more than the testimony of an Angel. When the proper time arrives, you shall be able to bear this testimony to the world. When you bear testimony that you have seen God. This testimony God will never suffer to fall, but will bear you out. although many will not give heed, yet others will.
						\hypertarget{EMS-1835-JS-21-FEB-1835-k}%
					You will,
						\marginnote{k}%
					therefore, see the necessity of getting this testimony from Heaven. Never cease striving until you have seen God, face to face. Strengthen your faith, cast off your doubts, your sins and all your unbelief and nothing can prevent you from coming to God. your ordination is not full and complete till God has laid his hand upon you. We require as much to qualify us as did those who have gone before us. God is the same. If the Saviour in former days laid his hands on his deciples. Why not in the latter Days. With regard to superiority I must make a few remarks. The ancient Apostles sought to be great. but, brethren, lest the seeds of discord be sown in this matter, understand particularly the voice of the spirit on this occasion. God does not love you better or more than others. You are to contend for the faith once deliverd to the saints. Jacob, you know wrestled till he obtained. It was by fervent prayer and diligent search that you have obtained the testimony you are now able to bear. You are as one, you are equal in bearing the keys of the kingdom to all nations. You are called to preach the gospel of the son of God to the nations of the earth.
							
						\hypertarget{EMS-1835-JS-21-FEB-1835-l}%
					It is the will
						\marginnote{l}%
					of your Heavenly Father that you proclaim his gospel to the ends of the earth. and the Islands of the sea. Be zealous to save souls. The soul of one man is as precious as the soul of another. You are to bear this message to those who consider themselves wise. and such may persecute you, they may seek your life. The adversary has always sought the life of the servants of God. You are, therefore, to be prepared at all times to make a sacrafice of your lives, should God require them in the advancment and building up of his cause. Murmur not at God. Be always prayerful, be always watchful. You will bear with me while I relieve the feelings of my heart. We shall not see another day like this. The time has fully come; The voice of the spirit has come to set these men apart. You will see the time when you will desire to see such a day as this, and you will not see it. Every heart wishes you peace \& prosperity. but the scene, with you, will inevitably change. Let no man take your bishopric, and beware that you lose not your crowns. It will require your whole souls. It will require courage like Enoch's.
						\hypertarget{EMS-1835-JS-21-FEB-1835-m}%
					The time is
						\marginnote{m}%
					near when you will be in the midst of congregations, who will gnash their teeth upon you. This gospel must roll and will roll till it fill the whole Earth. Did I say congregations would gnash upon you, yea I say nations will gnash upon you. You will be considered the worst of Men. Be not discouraged at this, When God pours out his Spirit, the enemy will rage, but, God, remember is on your right hand and on your left. A man, though he may be considered the worst, has joy who is conscious that he pleases God. The lives of those who proclaim the true gospel will be in danger; this has been the case ever since the days of righteous Abel. The same opposition has been manifest whenever men came forward to publish the gospel. The time is coming when you will be considered the worst by many \& by some the best of men. The time is coming when you will be perfectly familiar with the things of God. This testimony will make those who do not believe your testimony, seek your lives. But \sout{their} there are whole nations, who will receive your testimony. They will call you good men. Be not lifted up when you are called good men. Remember you are young men, and you shall be spared, I include the other three. Bear them in mind in your prayers carry their cases to a throne of grace. Although they are not present, yet you and they are equal. This appointment is calculated to create an affection in you, for each other, stronger than death.

						\hypertarget{EMS-1835-JS-21-FEB-1835-n}%
					You will travel
						\marginnote{n}%
					to other Nations, Bear each other in mind. If one or more is cast into prisons, let the others pray for him, and deliver him by their prayers. Your lives shall be in great jeopardy, but the promise of God, is that you shall be delivered. Remember you are not to go to other nations, till you receive your endowment. Tarry at Kirtland until you are endowed with power from on high. You need a fountain of wisdom, knowledge, and intelligence such as you never had. Relative to the endowment, I make a remark or two, that there be no mistake. The world cannot receive the things of God. He can endow you without wor[l]dly pomp or great parade. He can give you that wisdom, that intelligence and that power which characterized the ancient Saints and now \sout{characterized} characterizes the inhabitants of the upper world. The greatness of your commission, consists in this; you are to hold the keys of this ministry. You are to go to the nations afar off; nations that sit in darkness. The [day] is coming when the work of God must be done. Israel shall be gathered, The seed of Jacob shall be gathered from their long dispersion. There will be a feast to Israel, the elect of God. It is a sorrowful tale but, the gospel must be preached and his <God's> ministers be rejected. but where can Israel be found, and receive your testimony, and not rejoice? No where. The Prophecies are full of great things that are to take place in the last days. After the Elect is gathered out, destruction shall come on the inhabitants of the Earth; All nations shall feel the wrath of God, after they have been warned by the saints of the Most High. If you will not warn them others will and you will lose your crowns. You must prepare your minds to bid a long farewell to Kirtland. even till the great day come. You will see what you never expected to see. You will need the mind of Enoch or Elijah \& the faith of the brother of Jared.

						\hypertarget{EMS-1835-JS-21-FEB-1835-o}%
					You must be prepared
						\marginnote{o}%
					to walk by faith, however, appalling the prospect to human view. You, and each of you should feel the force of the imperious mandate, ``Son go labor in my vineyard'' and cheerfully receive what comes, but in the end you will stand while others will fall. You have read in the Revelation concerning ordination. Beware how you ordain, for all nations are not like this nation. They will willingly receive the ordinances at your hand to put you out of the way. There will be times, when nothing but the angels of God can deliver you out of their hand. We appeal to your intelligence, we appeal to your understanding, that we have so far discharged our duty to you. We consider it one of the greatest condescentions of our Heavenly Father in pointing you out to us. You will be stewards over this ministry.

						\hypertarget{EMS-1835-JS-21-FEB-1835-p}%
					We have work to do,
						\marginnote{p}%
					that no other men can do. You must proclaim the Gospel in its simplicity and purity. and we commend you to God and the word of his grace. You have our best wishes, you have our most fervent prayers, that you may be able to bear this testimony, that you have seen the face of God. Therefore, call upon him in faith and mighty prayer, till you prevail, for it <is> your duty and your privelege to bear such testimony for yourselves. We now exhort you to be faithful to fulfil your calling, there must be no lack here. You must fulfil in all things, and permit \sout{is} us to repeat, all nations have a claim on you. You are bound together as the three witnesses were. You, notwithstanding can part \& meet \& meet and part again till your heads are silvered o[v]er with age.

					He then took them separately by the hand and said, Do you with full purpose of heart take part in this ministry, to proclaim the gospel with all diligence with these your brethren, according to the tenor and intent of the charge you have received, each of which answered in the affirmative.

		% -----
		\subsubsection{Minutes and Discourses, 27 February 1835, as Reported by William E. McLellin}
		% -----
			\label{EMS-1835-JS-27-FEB-1835-WM}
			% \label{EMS-1830-JS-DAY-3LETTERMONTH-YEAR}
			
			\begin{description}
				\item[Print Sources]
			\end{description}

				\begin{tabular}{p{0.5in}p{1in}p{0.5in}p{1in}}
				JSP	 	& D4:247--252		& JSR 		& --- \\
				DC	 	& --- 				& HC 		& 2:198--200 \\
				HDDC 	& ---				& TCHC 		& 2:205--207 \\
				\end{tabular}
				% HDDC = woodford's DC diss; JSR = Marquardt's JS Revelations; TCHC = Vogel HC
			
			\begin{description}
				\item[Online Access] \href{https://www.josephsmithpapers.org/paper-summary/minutes-and-discourses-27-february-1835-as-reported-by-william-e-mclellin/1}{Here}
				% \item[Analysis] \hyperref[XX]{Here}
			\end{description}
			
			\begin{description}
				\item[Background]
			\end{description}
			
				% It might also be useful to give a two sentence context of the period: e.g., "This speech was given in Nauvoo to a small congregation one month prior to the destruction of the Nauvoo Expositor. These and other tensions may have been on the prophet's mind when he gave the speech."
			
			\begin{description}
				\item[Original Source]
			\end{description}
			
				% brief description of original source material, with reference to JSP or somewhere else for more info
				% Background material: it might be helpful to have a note that says whether a given document was presented as a speech, was written in a journal, etc.
				% Also a comment here about how reliable the source might be, or what shape it is in, if there are large omissions or parts that are hard to understand, and so on.
			
			\begin{description}
				\item[Text]
			\end{description}
			
					\hypertarget{EMS-1835-JS-27-FEB-1835-WM-a}%
				February 27\supers{th.}
					\marginnote{a}%
				of the same year [1835] the Twelve met in Kirtland by request of President J. Smith Jr. After the council was opened by prayer, he arose and made the following observations, (Viz) ``I have something to lay before this council, an item which they will find to be of great importance to them. I have for myself learned a fact by experience which on reflection gives me deep sorrow. It is a truth that if I now had in my possession every decision which has been given \sout{had} upon important items of doctrine and duties since the rise of this church, they would be of incalculable worth to the saints, but we have neglected to keep record of such things, thinking that prehaps that they would never benefit us afterwards, wh[i]ch had we now, would decide almost any point that might be agitated; and now we cannot bear record to the church nor unto the world of the great and glorious manifestations that have been made to us with that degree of power and authority whi[c]h we otherwise could if we had those decisions to publish abroad.

				Since the twelve are now chosen, I wish to tell them a course which they may pursue and be benefitted hereafter in a point of light of which they, prehaps, are not now aware. At all times when you assemble in the capacity of a council to transact business let the oldest of your number preside, and let one or more be appointed to keep a record of your proceedings and on the decision of every important item, be it what it may, let such decision be noted down, and they will ever after remain upon record as law, covenants and doctrine. \sout{Any} Questions thus decided might at the time appear unimportant, but should they be recorded and one of you lay hands upon them afterward you might find them of infinite worth not only to your brethren but a feast als[o] to your own souls.

					\hypertarget{EMS-1835-JS-27-FEB-1835-WM-b}%
				Should you
					\marginnote{b}%
				assemble from time to time and proceed to discuss important questions and pass decisions upon them and omit to record such decisions, by and by, you will be driven to straits from which you will not be able to extricate yourselves--- not being in a \sout{sufficient} situation to bring your faith to bare with sufficient perfection or power to obtain the desired information. Now in consequence of a neglect to write these things when God reveals them, not esteeming them of sufficient worth the spirit may withdraw and God may be angry, and here is a fountain of intelligence or knowledge of infinite importance which is lost. What was the cause of this? The answer is slothfulness or a neglect to appoint a man to occupy a few moments in writing. Here let me prophecy the time will come when if you neglect to do this, you will fall by the hands of unrighteous men. Were you to be brought before the authorities and accused of any crime or misdemeanor and be as innocent as the angels of God unless you can prove that you were somewhere else, your enemies will prevail against you: but if you can bring twelve men to testify that you were in some other place at that time you will escape their hands. Now if you will be careful to keep minutes of these things as I have said, it will be one of the most important and interesting records ever seen. I have now laid these things before you for your consideration and you are left to act according to your own judgments.''

					\hypertarget{EMS-1835-JS-27-FEB-1835-WM-c}%
				The council
					\marginnote{c}%
				then expressed their approbation of the foregoing remarks and proceeded to nominate and appoint Elders William [E.] McLellin and Orson Hyde to serve as clerks for the `twelve'.

				The following question was then proposed by president J Smith Jun. (viz) What importance is attached to the callings of these twelve apostles differrent from the other callings and offices of the chu[r]ch. After some discussion by Elders [David W.] Patten, [Brigham] Young, M'cLellin and W[illiam] Smith, the following decision was given by President Smith, the Prophet of God.

					\hypertarget{EMS-1835-JS-27-FEB-1835-WM-d}%
				``They are
					\marginnote{d}%
				the twelve apostles who are called to a travelling high council to preside over all the churches of the saints among the gentiles when there is no presidency established. They are to travel and preach among the Gentiles until the Lord shall shall command them to go to the Jews. They are to hold \sout{this} the keys of this ministry--- to unlock the door of the kingdom of heaven unto all nations and preach the Gospel unto every creation. This is the virtue powr and authority of their Apostleship--- Amen. It is all important that the twelve should understand the power and authority of the priesthoods, for without this knowledge they can do nothing to profit. In the first place God manifested himself to me and gave me authority to establish his church, and you have received your authority from God through me; and now it is your duty to go and unlock the kingdom of heaven to foreign nations, for no man can do that thing but yourselves. Neither has any man authority or a right to go to other nations before you; and you, twelve, stand in the same relation to those nations that I stand in to you, that is, as a minister; and you have each the same authority in other nations that I have in this nation.[''] The council was closed by Elder W. E. M'cLellin.

				\begin{tightright}
				William E. M'cLellin\} Clerk
				\end{tightright}

		% -----
		\subsubsection{Minutes and Discourses, 27 February 1835, as Reported by Oliver Cowdery}
		% -----
			\label{EMS-1835-JS-27-FEB-1835-OC}
			% \label{EMS-1830-JS-DAY-3LETTERMONTH-YEAR}
			
			\begin{description}
				\item[Print Sources]
			\end{description}

				\begin{tabular}{p{0.5in}p{1in}p{0.5in}p{1in}}
				JSP	 	& D4:252--254		& JSR 		& --- \\
				DC	 	& --- 				& HC 		& 2:198--200 \\
				HDDC 	& ---				& TCHC 		& 2:205--207 \\
				TPJS 	& 72--74			& 	 		&  \\
				\end{tabular}
				% HDDC = woodford's DC diss; JSR = Marquardt's JS Revelations; TCHC = Vogel HC
			
			\begin{description}
				\item[Online Access] \href{https://www.josephsmithpapers.org/paper-summary/minutes-and-discourses-27-february-1835-as-reported-by-oliver-cowdery/1}{Here}
				% \item[Analysis] \hyperref[XX]{Here}
			\end{description}
			
			\begin{description}
				\item[Background]
			\end{description}
			
				% It might also be useful to give a two sentence context of the period: e.g., "This speech was given in Nauvoo to a small congregation one month prior to the destruction of the Nauvoo Expositor. These and other tensions may have been on the prophet's mind when he gave the speech."
			
			\begin{description}
				\item[Original Source]
			\end{description}
			
				% brief description of original source material, with reference to JSP or somewhere else for more info
				% Background material: it might be helpful to have a note that says whether a given document was presented as a speech, was written in a journal, etc.
				% Also a comment here about how reliable the source might be, or what shape it is in, if there are large omissions or parts that are hard to understand, and so on.
			
			\begin{description}
				\item[Text]
			\end{description}
			
				\begin{tightright}
						\hypertarget{EMS-1835-JS-27-FEB-1835-OC-a}%
					Kirtland Ohio Feb. 27\supers{th} 1835
						\marginnote{a}
				\end{tightright}

				This evening a meeting of nine of the twelve of the Apostles, who had been chosen and ordained was held at the house of President Joseph Smith \supersu{Junr}\supers{.} with Presidents Sidney Rigdon and Frederick G. Williams and certain others (Viz.) Bishop Newel K. Whitney \& Elders John Smith, John Murdock and Even [Evan] M. Greene. It may be well to give, also, the names of those nine who were present. Lyman Johnson, Brigham Young, Heber C. Kimball, Orson Hyde, David W. Patten Luke Johnson, Wm. E. M\supers{c}Lel[l]in, John F. Boynton \& Wm. Smith. Parley P. Pratt had now gone to New-Portage and Orson Pratt \& Thomas B. Marsh had not yet arrived to receive their ordination After the company were assembled, and the meeting opened by prayer of President J. Smith \supersu{Jun}\supers{r.} he (President Smith) arose and said, he had something to lay before the council, and he thought if he were heard patiently, he could lay before the council an item which they would [find] to be of importance. He had for himself, learned a fact, by experience, which on reflection, always gave him a deep sorrow.
					\hypertarget{EMS-1835-JS-27-FEB-1835-OC-b}%
				It is
					\marginnote{b}%
				a fact (said President Smith) that if I now had in my possession, every decision which has been had, upon important items of doctrine and duties, which have been given since the commencement of this work, I would not part with it for any sum of money, but we have neglected to take minutes of such things, thinking, perhaps, that they would never benefit us afterward, which, had we now, would decide almost any point of doctrine, which might be agitated. but, this has been neglected and now we cannot bear record to the church, and to the world, of the great and glorious manifestations which have been made to us, with that degree of power and authority, we otherwise could, if we now had these things to publish abroad. Since the Twelve are now chosen, I wish to tell them a course <which> they may pursue, and be benefitted hereafter, in a point of light of which they are not now aware. If they will, on every time they assemble, appoint a person to preside <over> \sout{of} over them during the meeting and one or more to keep a record of their proceedings, and on the decision of every question or item, let it be what it may, let such decision be noted down, such decision will forever remain upon record, and appear an item of covenant or doctrine. An item thus decided may appear at the time of little or no worth, but should it be published, and one of you lay hands on it after you will find of infinite worth, not only to your brethren, but it will be a feast to your own souls.

					\hypertarget{EMS-1835-JS-27-FEB-1835-OC-c}%
				Here is
					\marginnote{c}%
				another important item. If you assemble from time to time, and proceed to discuss important questions, and pass decisions upon the same, and fail to note them down, by \& by you will be driven to straits, from which, you will not be able to extricate yourselves. because you may be in a situation not to bring your faith to bear with sufficient perfection or power to obtain the desired information, or perhaps, for neglecting to write these things, when God revealed them, not esteeming them of sufficient worth the spirit may withdraw and God may be angry, and here is or was a vast knowledge of infinite importance, which is now lost. What was the cause of this? It came in consequence of Slothfulness, or a neglect to appoint a man to occupy a few moments in writing all these decisions. Here let me prophesy. The time will come, when, if you neglect to do this thing, you will fall by the hands of unrighteous men. Were you to be brought before the authorities, and, be accused of any crime or Misdemeanor, and be as innocent as the Angels of God, unless you can prove yourselves to have been somewhere else, your enemies will prevail over you: but if you can bring twelve men to testify that you were in a certain place at that time, you will escape their hands.
					\hypertarget{EMS-1835-JS-27-FEB-1835-OC-d}%
				Now,
					\marginnote{d}%
				if you will be careful to keep minutes of these things, as I have said, it will be one of the most important records ever seen; for every such decision will, ever after remain as items of doctrine and covenants. I have now placed before you these items, for your consideration, and you are left to act according to your own judgements. The council then expressed their approbation, concerning the foregoing remarks of President Smith, and proceeded to appoint Elders, Orson Hyde \& Wm. E. M\supers{c}Lelin, to serve as clerks for the meeting. After which the following question was proposed by President Smith (viz.) What importance is there attached to the calling of these twelve apostles different from the other callings or officces of the Church? After some discussion by elders, Patten Young, Wm. Smith \& M\supers{c}Lelin, President Smith gave the following decision. They are the Twelve Apostles, who are called to the office of traveling high council, who are to preside over all the churches of the Saints among the Gentiles, where there is no presidency established, and they are to travel and preach among the Gentiles, until the Lord shall command them to go to the Jews. They are to hold the keys of this ministry, to unlock the door of the kingdom of heaven unto all nations, and to preach the Gospel to every creature. This is the power, authority and virtue of their Apostleship. The meeting closed by President J. Smith \supers{Ju}\supersu{n}\supers{\supersu{r}} in Prayer

				\begin{tightright}
				Oliver Cowdery Clerk.
				\end{tightright}

		% -----
		\subsubsection{Minutes, Discourse, and Blessings, 1 March 1835}
		% -----
			\label{EMS-1835-JS-1-MAR-1835}
			% \label{EMS-1830-JS-DAY-3LETTERMONTH-YEAR}
			
			\begin{description}
				\item[Print Sources]
			\end{description}

				\begin{tabular}{p{0.5in}p{1in}p{0.5in}p{1in}}
				JSP	 	& D4:264--279		& JSR 		& --- \\
				DC	 	& --- 				& HC 		& 2:204 \\
				HDDC 	& ---				& TCHC 		& 2:210 \\
				\end{tabular}
				% HDDC = woodford's DC diss; JSR = Marquardt's JS Revelations; TCHC = Vogel HC
			
			\begin{description}
				\item[Online Access] \href{https://www.josephsmithpapers.org/paper-summary/minutes-discourse-and-blessings-1-march-1835/1}{Here}
				% \item[Analysis] \hyperref[XX]{Here}
			\end{description}
			
			\begin{description}
				\item[Background]
			\end{description}
			
				% It might also be useful to give a two sentence context of the period: e.g., "This speech was given in Nauvoo to a small congregation one month prior to the destruction of the Nauvoo Expositor. These and other tensions may have been on the prophet's mind when he gave the speech."
			
			\begin{description}
				\item[Original Source]
			\end{description}
			
				% brief description of original source material, with reference to JSP or somewhere else for more info
				% Background material: it might be helpful to have a note that says whether a given document was presented as a speech, was written in a journal, etc.
				% Also a comment here about how reliable the source might be, or what shape it is in, if there are large omissions or parts that are hard to understand, and so on.
			
			\begin{description}
				\item[Text]
			\end{description}
			
				\begin{tightright}
				March 1\supersu{st} 1835
				\end{tightright}
			
				The council met this morning (Sabbath) to proceed with the ordination of the 70,%
					\footnote{%
						\label{EMS-1835-JS-1-MAR-1835-NOTE-seventy}%
						Note about all the blessings which occur in this meeting and in the previous one as well.
						}
				There being several who had been recently baptized present, and being expedient that they should be confirmed, and also that the sacrament should be administred to the Church, the business of ordination was suspended for a short time. President Joseph Smith Jun\supersu{r} addressed the Church \& the council and the church upon the propriety upon the propriety of attending to this ordinance with pure hearts, and pure desirles. He touched upon the propriety of this institution in the church \& urged the vast importance of doing it with acceptance to the Lord, He asked how long do you suppose a man may partake of this ordinance unworthily and the Lord not withdraw his spirit from him? How long will he thus trifle with sacred things and the Lord not give him over to the buffitings of Satan until the day of Redemption? The church should know if they are unworthy, from time to time, to partake, The servants of the Lord will be forbidden to administer it. Therefore our heart ought \sout{ought} to humble themselves, and we to repent of our sins, and put away evil from among us. some further remarks were made and the sacrament was then administred...
				
		% -----
		\subsubsection{Minutes and Discourses, 7--8 March 1835}
		% -----
			\label{EMS-1835-JS-7-8-MAR-1835}
			% \label{EMS-1830-JS-DAY-3LETTERMONTH-YEAR}
			
			\begin{description}
				\item[Print Sources]
			\end{description}

				\begin{tabular}{p{0.5in}p{1in}p{0.5in}p{1in}}
				JSP	 	& D4:279--287		& JSR 		& --- \\
				DC	 	& --- 				& HC 		& 2:205--208 \\
				HDDC 	& ---				& TCHC 		& 2:211--215 \\
				\end{tabular}
				% HDDC = woodford's DC diss; JSR = Marquardt's JS Revelations; TCHC = Vogel HC
			
			\begin{description}
				\item[Online Access] \href{https://www.josephsmithpapers.org/paper-summary/minutes-and-discourses-7-8-march-1835/1}{Here}
				% \item[Analysis] \hyperref[XX]{Here}
			\end{description}
			
			\begin{description}
				\item[Background]
			\end{description}
			
				% It might also be useful to give a two sentence context of the period: e.g., "This speech was given in Nauvoo to a small congregation one month prior to the destruction of the Nauvoo Expositor. These and other tensions may have been on the prophet's mind when he gave the speech."
			
			\begin{description}
				\item[Original Source]
			\end{description}
			
				% brief description of original source material, with reference to JSP or somewhere else for more info
				% Background material: it might be helpful to have a note that says whether a given document was presented as a speech, was written in a journal, etc.
				% Also a comment here about how reliable the source might be, or what shape it is in, if there are large omissions or parts that are hard to understand, and so on.
			
			\begin{description}
				\item[Text]
			\end{description}
			
				\begin{tightcenter}
				Kirtland March 7\supersu{th} 1835
				\end{tightcenter}
				
				This day a Meeting of the Church of Latter Day Saints was called in this place, for the purpose of blessing in the name of the Lord, those who have heretofore assisted in building, by their labor \& other means, the house of the Lord in this place. The forenoon was occupied by Pr. J. Smith \supersu{Jun}\supers{\supersu{r}} in remarks to the Church, upon the propriety and necessity of purifying itself. In the P.M. the names \sout{of the} several, those who had assisted to build the house were taken and further remarks were made by president J. Smith \supersu{Jun}\supers{\supers{r.}} He said that those who had distinguished themselves Thus far by consecrating to the upbuilding of said house as well as laboring were to be remembered. That those who build it should own it, and have the control of it. After further remarks he proceeded to call a vote of those who had performed this labor, whether they would still go on and perform the remaining part of the same,
				
				Passed by unanimous voice...
				
		% -----
		\subsubsection{Minutes, 12 March 1835}
		% -----
			\label{EMS-1835-JS-12-MAR-1835}
			% \label{EMS-1830-JS-DAY-3LETTERMONTH-YEAR}
			
			\begin{description}
				\item[Print Sources]
			\end{description}

				\begin{tabular}{p{0.5in}p{1in}p{0.5in}p{1in}}
				JSP	 	& D4:287--289		& JSR 		& --- \\
				DC	 	& --- 				& HC 		& 2:209 \\
				HDDC 	& ---				& TCHC 		& 2:217 \\
				\end{tabular}
				% HDDC = woodford's DC diss; JSR = Marquardt's JS Revelations; TCHC = Vogel HC
			
			\begin{description}
				\item[Online Access] \href{https://www.josephsmithpapers.org/paper-summary/minutes-12-march-1835/1}{Here}
				% \item[Analysis] \hyperref[XX]{Here}
			\end{description}
			
			\begin{description}
				\item[Background]
			\end{description}
			
				% It might also be useful to give a two sentence context of the period: e.g., "This speech was given in Nauvoo to a small congregation one month prior to the destruction of the Nauvoo Expositor. These and other tensions may have been on the prophet's mind when he gave the speech."
			
			\begin{description}
				\item[Original Source]
			\end{description}
			
				% brief description of original source material, with reference to JSP or somewhere else for more info
				% Background material: it might be helpful to have a note that says whether a given document was presented as a speech, was written in a journal, etc.
				% Also a comment here about how reliable the source might be, or what shape it is in, if there are large omissions or parts that are hard to understand, and so on.
			
			\begin{description}
				\item[Text]
			\end{description}
			
				Evening of March 12\supers{th.} the twelve assembled and the counsel was opened by president J. Smith Jun. and he proposed that we take our first mission through the eastern States to the Atlantic Ocean and hold conferences in the vicinity of the several branches of the church for the purpose of regulateing all things necessary for their welfare. It was proposed that the twelve should leave Kirtland on the 4\supers{th.} May which was unanimously agreed \sout{too} to. It was then proposed that during their present mission, Elder B[righam] Young should open a door to the remnants of Joseph who dwelt among the Gentiles which was carrid...
				
		% -----
		\subsubsection{Minutes, 16 March 1835}
		% -----
			\label{EMS-1835-JS-16-MAR-1835}
			% \label{EMS-1830-JS-DAY-3LETTERMONTH-YEAR}
			
			\begin{description}
				\item[Print Sources]
			\end{description}

				\begin{tabular}{p{0.5in}p{1in}p{0.5in}p{1in}}
				JSP	 	& D4:289--293		& JSR 		& --- \\
				DC	 	& --- 				& HC 		& --- \\
				HDDC 	& ---				& TCHC 		& --- \\
				\end{tabular}
				% HDDC = woodford's DC diss; JSR = Marquardt's JS Revelations; TCHC = Vogel HC
			
			\begin{description}
				\item[Online Access] \href{https://www.josephsmithpapers.org/paper-summary/minutes-16-march-1835/1}{Here}
				% \item[Analysis] \hyperref[XX]{Here}
			\end{description}
			
			\begin{description}
				\item[Background]
			\end{description}
			
				% It might also be useful to give a two sentence context of the period: e.g., "This speech was given in Nauvoo to a small congregation one month prior to the destruction of the Nauvoo Expositor. These and other tensions may have been on the prophet's mind when he gave the speech."
			
			\begin{description}
				\item[Original Source]
			\end{description}
			
				% brief description of original source material, with reference to JSP or somewhere else for more info
				% Background material: it might be helpful to have a note that says whether a given document was presented as a speech, was written in a journal, etc.
				% Also a comment here about how reliable the source might be, or what shape it is in, if there are large omissions or parts that are hard to understand, and so on.
			
			\begin{description}
				\item[Text]
			\end{description}
			
				...A. High Council was called for church business.
				
				...John Smith opened the case in few words stating that the matter being so intricate that he could not bring it as a charge against any one of these brethren particularly, and the brethren severally were called upon to make their own remarks. After this the counsellors severally spoke, and the complainant and the accused made their remarks, when the case was submitted, President Williams being short, President Rigdon spoke one hour When President J. Smith Jun\supersu{r} proceeded to give decision After speaking to considerable extent said: The decision is, that brother Lyman as well as brother Cutler be rebuked and stand reproved, until they openly acknowledge before the church, that they have injured Elder Cahoon by speaking evil of his character, and conduct without any just ground or occasion, and that as Elder Cutler has heretofore declared before the Church, that he has not murmured nor complained, because he has not received pay for working on the house, to his satisfaction, and it is known to the contrary, that he be as frank in acknowledging this fact, as he has been in making the declaration. The vote was then taken and passed by unanimous voice, and the council closed by prayer of President S. Rigdon...
				
		% -----
		\subsubsection{Minutes, 26 April 1835}
		% -----
			\label{EMS-1835-JS-26-APR-1835}
			% \label{EMS-1830-JS-DAY-3LETTERMONTH-YEAR}
			
			\begin{description}
				\item[Print Sources]
			\end{description}

				\begin{tabular}{p{0.5in}p{1in}p{0.5in}p{1in}}
				JSP	 	& D4:293--295		& JSR 		& --- \\
				DC	 	& --- 				& HC 		& 2:218--219 \\
				HDDC 	& ---				& TCHC 		& 2:225 \\
				\end{tabular}
				% HDDC = woodford's DC diss; JSR = Marquardt's JS Revelations; TCHC = Vogel HC
			
			\begin{description}
				\item[Online Access] \href{https://www.josephsmithpapers.org/paper-summary/minutes-26-april-1835/1}{Here}
				% \item[Analysis] \hyperref[XX]{Here}
			\end{description}
			
			\begin{description}
				\item[Background]
			\end{description}
			
				% It might also be useful to give a two sentence context of the period: e.g., "This speech was given in Nauvoo to a small congregation one month prior to the destruction of the Nauvoo Expositor. These and other tensions may have been on the prophet's mind when he gave the speech."
			
			\begin{description}
				\item[Original Source]
			\end{description}
			
				% brief description of original source material, with reference to JSP or somewhere else for more info
				% Background material: it might be helpful to have a note that says whether a given document was presented as a speech, was written in a journal, etc.
				% Also a comment here about how reliable the source might be, or what shape it is in, if there are large omissions or parts that are hard to understand, and so on.
			
			\begin{description}
				\item[Text]
			\end{description}
				
				\begin{tightright}
				Kirtland April 26, 1835.
				\end{tightright}
				
				This day, pursuant to previous appointment, the Twelve Apostles and the Seventy (a part of whom had already been chosen,) assembled in the temple (altho unfinished.) with a numerous concourse of people in order to receive our charge and instructions from President Joseph Smith Jun relative to our mission and duties...
				
		% -----
		\subsubsection{Blessing to Thomas B. Marsh, 26 April 1835}
		% -----
			\label{EMS-1835-JS-26-APR-1835-TM}
			% \label{EMS-1830-JS-DAY-3LETTERMONTH-YEAR}
			
			\begin{description}
				\item[Print Sources]
			\end{description}

				\begin{tabular}{p{0.5in}p{1in}p{0.5in}p{1in}}
				JSP	 	& ---\footnote{\label{EMS-1835-JS-26-APR-1835-TM-jsp}	Explain its absence, and the absence of the other blessings below.} & JSR 		& --- \\
				DC	 	& --- 		& HC 		& --- \\
				HDDC 	& ---		& TCHC 		& --- \\
				\end{tabular}
				% HDDC = woodford's DC diss; JSR = Marquardt's JS Revelations; TCHC = Vogel HC
			
			\begin{description}
				\item[Online Access] \href{https://www.josephsmithpapers.org/paper-summary/blessing-to-thomas-b-marsh-26-april-1835/1}{Here}
				% \item[Analysis] \hyperref[XX]{Here}
			\end{description}
			
			\begin{description}
				\item[Background]
			\end{description}
			
				% It might also be useful to give a two sentence context of the period: e.g., "This speech was given in Nauvoo to a small congregation one month prior to the destruction of the Nauvoo Expositor. These and other tensions may have been on the prophet's mind when he gave the speech."
			
			\begin{description}
				\item[Original Source]
			\end{description}
			
				% brief description of original source material, with reference to JSP or somewhere else for more info
				% Background material: it might be helpful to have a note that says whether a given document was presented as a speech, was written in a journal, etc.
				% Also a comment here about how reliable the source might be, or what shape it is in, if there are large omissions or parts that are hard to understand, and so on.
			
			\begin{description}
				\item[Text]
			\end{description}
			
				Thomas B. Marsh'\supers{s} Blessing by O. Cowdery Dear Brother, You are to be a minister of righteousness and to this ministry and Apostleship you are now to be ordained: and May all temporal and spiritual blessings attend you. Your sins are forgiven you, and you are to go forth and preach the everlasting Gospel. You shall travel from kingdom to kingdom and from nation to nation. Angels shall bear the[e] up, and thou shalt be instrumental in bringing thousands of the redeemed of the Lord to Zion. President David Whitmer sealed the above blessing upon him, even so Amen.%
					\footnote{%
						\label{EMS-1835-JS-26-APR-1835-TM-NOTE-blessings}%
						Connected this blessing and the next one to OP back to the others that happened earlier in the year.
						}
				
		% -----
		\subsubsection{Blessing to Orson Pratt, 26 April 1835}
		% -----
			\label{EMS-1835-JS-26-APR-1835-OP}
			% \label{EMS-1830-JS-DAY-3LETTERMONTH-YEAR}
			
			\begin{description}
				\item[Print Sources]
			\end{description}

				\begin{tabular}{p{0.5in}p{1in}p{0.5in}p{1in}}
				JSP	 	& ---		& JSR 		& --- \\
				DC	 	& --- 		& HC 		& --- \\
				HDDC 	& ---		& TCHC 		& --- \\
				\end{tabular}
				% HDDC = woodford's DC diss; JSR = Marquardt's JS Revelations; TCHC = Vogel HC
			
			\begin{description}
				\item[Online Access] \href{https://www.josephsmithpapers.org/paper-summary/blessing-to-orson-pratt-26-april-1835/1}{Here}
				% \item[Analysis] \hyperref[XX]{Here}
			\end{description}
			
			\begin{description}
				\item[Background]
			\end{description}
			
				% It might also be useful to give a two sentence context of the period: e.g., "This speech was given in Nauvoo to a small congregation one month prior to the destruction of the Nauvoo Expositor. These and other tensions may have been on the prophet's mind when he gave the speech."
			
			\begin{description}
				\item[Original Source]
			\end{description}
			
				% brief description of original source material, with reference to JSP or somewhere else for more info
				% Background material: it might be helpful to have a note that says whether a given document was presented as a speech, was written in a journal, etc.
				% Also a comment here about how reliable the source might be, or what shape it is in, if there are large omissions or parts that are hard to understand, and so on.
			
			\begin{description}
				\item[Text]
			\end{description}
			
				\begin{tightcenter}
				Orson Pratt's blessing.
				\end{tightcenter}
				
				Dear Brother, You are chosen, and set apart to be ordained to this Apostleship and this ministry. You shall go forth and preach the gospel and do a mighty work. You shall be sustained. The Holy spirit shall enlighten thy mind. Thou shalt travel from nation to nation. The Lord God shall preserve thee and return thee safe, with songs of everlasting joy upon thy head. The above was confirmed, by President David Whitmer
				
		% -----
		\subsubsection{Minutes and Discourse, 2 May 1835, as Reported by William E. McLellin--A}
		% -----
			\label{EMS-1835-JS-2-MAY-1835-A}
			% \label{EMS-1830-JS-DAY-3LETTERMONTH-YEAR}
			
			\begin{description}
				\item[Print Sources]
			\end{description}

				\begin{tabular}{p{0.5in}p{1in}p{0.5in}p{1in}}
				JSP	 	& D4:299--306		& JSR 		& --- \\
				DC	 	& --- 				& HC 		& 2:219--222 \\
				HDDC 	& ---				& TCHC 		& 2:225--228 \\
				TPJS 	& 74--76			& 	 		&  \\
				\end{tabular}
				% HDDC = woodford's DC diss; JSR = Marquardt's JS Revelations; TCHC = Vogel HC
			
			\begin{description}
				\item[Online Access] \href{https://www.josephsmithpapers.org/paper-summary/minutes-and-discourse-2-may-1835-as-reported-by-william-e-mclellin-a/1}{Here}
				% \item[Analysis] \hyperref[XX]{Here}
			\end{description}
			
			\begin{description}
				\item[Background]
			\end{description}
			
				% It might also be useful to give a two sentence context of the period: e.g., "This speech was given in Nauvoo to a small congregation one month prior to the destruction of the Nauvoo Expositor. These and other tensions may have been on the prophet's mind when he gave the speech."
			
			\begin{description}
				\item[Original Source]
			\end{description}
			
				% brief description of original source material, with reference to JSP or somewhere else for more info
				% Background material: it might be helpful to have a note that says whether a given document was presented as a speech, was written in a journal, etc.
				% Also a comment here about how reliable the source might be, or what shape it is in, if there are large omissions or parts that are hard to understand, and so on.
			
			\begin{description}
				\item[Text]
			\end{description}
				
				\begin{tightcenter}
						\hypertarget{EMS-1835-JS-2-MAY-1835-A-a}%
					Kirtland May 2\supers{d---} 1835
						\marginnote{a}
				\end{tightcenter}
				
				According to appointment the Presidency, the Twelve a part of the Seventy, and some other Elders of the church met in conference this morning in order to consult the affairs of the church. \&c. Conference was opened by Elder Brigham Young by prayer. President J. Smith Jun\supersu{r} presiding. After conference was opened and the Twelve took their seats. he stated that it would be the duty of the twelve to appoint the oldest one of their number to preside in their councils, beginning at the oldest and so on until the youngest has presided and then beginning at the oldest again. \&c. The Twelve took their seats regularly according to their ages as follows. T[homas] B. Marsh David W. Patten. Brigham Young, Heber C. Kimball Orson Hyde, Wm. E. M\supers{c}Lel[l]in Parley P. Pratt \sout{J[ohn] F. Boynton} Luke Johnson William Smith Orson Pratt John F. Boynton \& Lyman Johnson. The president then stated that the Twelve will have no right to go into Zion or any of its stakes and there undertake to regulate the affairs thereof where there is a standing High Council. But it is their duty to go abroad and regulate all matters relative to the different branches of the Church. When the Twelve are together, or a quorum of them in any Church, they will have authority to act independently and make decisions, and those decisions, and those decissions are valid. But where there is not a quorum they will have to do business by the voice of the Church.
					\hypertarget{EMS-1835-JS-2-MAY-1835-A-b}%
				No
					\marginnote{b}%
				standing high council has authority \sout{his} to go into the churches abroad and regulate the matters thereof, for this belongs to the Twelve. No High Council will ever be established only in Zion or one of its Stakes. When the twelve pass a decision it is in the name of the church, Therefore, it is valid. No individual has a right to go into any church and ordain any minister for the Church, unless it is by the Voice of the Church. No Elder has a right to go into any branch of the church and appoint meetings or regulate the Church without the consent or advic[e] of the presiding Elder of said branch. If the first Seventy are all occupied, and there is a call for more laborers it will be the duty of the seven presidents of the first seventy to call and ordain other Seventy and send them forth to labor, in the vineyard until if need be they set a part apart seven times Seventy, even until there shall be one hundred \& forty and four thousand. The Seventy are not to attend the conferences of the Twelve unless they are called upon or requested to by the Twelve. The twelve and the Seventy have particularly to depend upon their Ministry for their support and that of their families, and they have a right by virtue of their offices to call upon the Church to assist them...
				
					\hypertarget{EMS-1835-JS-2-MAY-1835-A-c}%
				...The Elders
					\marginnote{c}%
				of Kirtland and its vicinity were next called upon, or their circumstances were considered. Their names being enrolled. President J. Smith \supersu{Jun}\supers{\supers{r}} arose with the list in his hands and made some very appropriate remarks, relative to the deliverance of Zion and so much of the Authority being present, he moved that we never give up the struggle for Zion, even until Death. or until Zion is Redeemed. The vote was unanimous and with apparent deep feeling. Voted that all the Elders of the Church are bound to travel in the World to preach the gospel with all their might mind \& Strength when their circumstances will admit of it, and that the door is now opened. Motioned, seconded \& voted that Elder Brigham Youngs, John P. Green[e] \& Amos Orton are appointed to go and preach the gospel to the remnants of Joseph. <the door to be opened by> \sout{Motioned, seconded \& voted} Elder B. Young and this will open a door to all the house of Joseph...
				
		% -----
		\subsubsection{Minutes and Discourse, 2 May 1835, as Reported by William E. McLellin--B}
		% -----
			\label{EMS-1835-JS-2-MAY-1835-B}
			% \label{EMS-1830-JS-DAY-3LETTERMONTH-YEAR}
			
			\begin{description}
				\item[Print Sources]
			\end{description}

				\begin{tabular}{p{0.5in}p{1in}p{0.5in}p{1in}}
				JSP	 	& D4:306--308		& JSR 		& --- \\
				DC	 	& --- 				& HC 		& 2:219--222 \\
				HDDC 	& ---				& TCHC 		& 2:225--228 \\
				TPJS 	& 74--76			& 	 		&  \\
				\end{tabular}
				% HDDC = woodford's DC diss; JSR = Marquardt's JS Revelations; TCHC = Vogel HC
			
			\begin{description}
				\item[Online Access] \href{https://www.josephsmithpapers.org/paper-summary/minutes-and-discourse-2-may-1835-as-reported-by-william-e-mclellin-b/1}{Here}
				% \item[Analysis] \hyperref[XX]{Here}
			\end{description}
			
			\begin{description}
				\item[Background]
			\end{description}
			
				% It might also be useful to give a two sentence context of the period: e.g., "This speech was given in Nauvoo to a small congregation one month prior to the destruction of the Nauvoo Expositor. These and other tensions may have been on the prophet's mind when he gave the speech."
			
			\begin{description}
				\item[Original Source]
			\end{description}
			
				% brief description of original source material, with reference to JSP or somewhere else for more info
				% Background material: it might be helpful to have a note that says whether a given document was presented as a speech, was written in a journal, etc.
				% Also a comment here about how reliable the source might be, or what shape it is in, if there are large omissions or parts that are hard to understand, and so on.
			
			\begin{description}
				\item[Text]
			\end{description}
					
					\hypertarget{EMS-1835-JS-2-MAY-1835-B-a}%
				...In the
					\marginnote{a}%
				midst of this grand assemblage, President J Smith Jun arose and made many remarks, among which were the following. ``It will be the duty of the twelve when in council to take their seats together according to their ages. The oldest to be seated at the head, and preside in the first council, the next oldest in the second; and so on until the youngest has presided''.

				The twelve Apostles have no right to go into Zion or any of its stakes where there is a regular high council established, to regulate any matters pertaining thereunto: But it is their duty to go abroad and regulate and set in order all matters relative to the different branches of this church of the Latter Day Saints.

				When the twelve are all together or a quorum of them in any church, they have authority to act independently of the church and form decisions and those decisions will be valid; but when there is not a quorum of them together, they must transact business by the common consent of the church.

					\hypertarget{EMS-1835-JS-2-MAY-1835-B-b}%
				No standing hig[h]
					\marginnote{b}%
				council has authority to go into the churches abroad and regulate the matters thereof, for this belongs to the Twelve.

				No standing high council will ever be established except in Zion or one of its stakes.

				When the Twelve pass a decision, it is in the name of the church, therefore, it is valid.

				No official member of the church has authority to go into any branch thereof and ordain any minister for the church unless it is by the voice of that branch. No elder has authority to go into any branch of the church and appoint meetings or attempt to regulate the affairs of the church without the advice and consent of the presiding Elder of that branch.

				If the first seventy are all employed and there is a call for more labourers in the vineyard it will be the the duty of the seven presidents of the first seventy to call and ordain other seventy and send them forth to labour until, if need be, they set apart seven times seventy, and even until there are one hundred and forty \& four thousand thus set apart to the ministry. The seventy are not to attend the conferences of the Twelve, unless they are called upon or requested so to do by the twelve.

				The twelve and the seventy have particula[r]ly to depend upon their ministry for their support and that of their families, and they have a right by virtue of their offices to call upon the churches to assist them.

				Resolved in this grand Council; That we never give up the struggle for Zion until it is redeemed altho'. we should die in the contest. The vote was unanimous of all that were in the house.

				\begin{tightright}
				W. E. Mc.Lellin\} Clerk
				\end{tightright}
				
		% -----
		\subsubsection{Instruction on Priesthood, between circa 1 March and circa 4 May 1835 [D\&C 107]}
		% -----
			\label{EMS-1835-JS-CIR-1-MAR-4-MAY-1835}
			% \label{EMS-1830-JS-DAY-3LETTERMONTH-YEAR}
			
			\begin{description}
				\item[Print Sources]
			\end{description}

				\begin{tabular}{p{0.5in}p{1in}p{0.5in}p{1in}}
				JSP	 	& D4:308--321		& JSR 		& 267--273 \\
				DC	 	& Section 107 		& HC 		& 2:210--217 \\
				HDDC 	& 1398--1432		& TCHC 		& 2:218--222 \\
				\end{tabular}
				% HDDC = woodford's DC diss; JSR = Marquardt's JS Revelations; TCHC = Vogel HC
			
			\begin{description}
				\item[Online Access] \href{https://www.josephsmithpapers.org/paper-summary/instruction-on-priesthood-between-circa-1-march-and-circa-4-may-1835-dc-107/1}{Here}
				% \item[Analysis] \hyperref[DC107]{Here}
			\end{description}
			
			\begin{description}
				\item[Background]
			\end{description}
			
				% It might also be useful to give a two sentence context of the period: e.g., "This speech was given in Nauvoo to a small congregation one month prior to the destruction of the Nauvoo Expositor. These and other tensions may have been on the prophet's mind when he gave the speech."
			
			\begin{description}
				\item[Original Source]
			\end{description}
			
				% brief description of original source material, with reference to JSP or somewhere else for more info
				% Background material: it might be helpful to have a note that says whether a given document was presented as a speech, was written in a journal, etc.
				% Also a comment here about how reliable the source might be, or what shape it is in, if there are large omissions or parts that are hard to understand, and so on.
			
			\begin{description}
				\item[Text]
			\end{description}
			
				\begin{tightcenter}
					SECTION III.

					ON PRIESTHOOD.
				\end{tightcenter}

				1 There are, in the church, two priesthoods, namely: the Melchizedek, and the Aaronic, including the Levitical priesthood. Why the first is called the Melchizedek priesthood, is because Melchizedek was such a great high priest: before his day it was called \textit{the holy priesthood, after the order of the Son of God;} but out of respect or reverence to the name of the Supreme Being, to avoid the too frequent repetition of his name, they, the church, in ancient days, called that priesthood after Melchizedek, or the Melchizedek priesthood.

				2 All other authorities, or offices in the church are appendages to this priesthood; but there are two divisions, or grand heads---one is the Melchizedek priesthood, and the other is the Aaronic, or Levitical priesthood.

				3 The office of an elder comes under the priesthood of Melchizedek. The Melchisedek priesthood holds the right of presidency, and has power and authority over all the offices in the church, in all ages of the world, to administer in spiritual things.

				4 The presidency of the high priesthood, after the order of Melchizedek, have a right to officiate in all the offices in the church.

				5 High priests, after the order of the Melchizedek priesthood, have a right to officiate in their own standing, under the direction of the presidency, in administering spiritual things, and also in the office of an elder, priest, (of the Levitical order,) teacher, deacon and member.

				6 An elder has a right to officiate in his stead when the high priest is not present.

				7 The high priest, and elder, are to administer in spiritual things, agreeably to the covenants and commandments of the church; and they have a right to officiate in all these offices of the church when there are no higher authorities present.

				8 The second priesthood is called the priesthood of Aaron, because it was conferred upon Aaron and his seed, throughout all their generations. Why it is called the lesser priesthood, is because it is an appendage to the greater, or the Melchizedek priesthood, and has power in administering outward ordinances. The bishopric is the presidency of this priesthood and holds the keys, or authority of the same. No man has a legal right to this office, to hold the keys of this priesthood, except he be a litteral descendant of Aaron. But as a high priest, of the Melchizedek priesthood, has authority to officiate in all the lesser offices, he may officiate in the office of bishop when no literal descendant of Aaron can be found; provided he is called and set apart and ordained unto this power by the hands of the presidency of the Melchizedek priesthood.

				9 The power and authority of the higher or Melchizedek priesthood, is to hold the keys of all the spiritual blessings of the church---to have the privilege of receiving the mysteries of the kingdom of heaven---to have the heavens opened unto them---to commune with the general assembly and church of the first born, and to enjoy the communion and presence of God the Father, and Jesus the Mediator of the new covenant.

				10 The power and authority of the lesser, or Aaronic priesthood, is, to hold the keys of the ministring of angels, and to administer in outward ordinances---the letter of the gospel---the baptism of repentance for the remission of sins, agreeably to the covenants and commandments.

				11 Of necessity there are presidents, or presiding offices growing out of, or appointed of, or from among those who are ordained to the several offices in these two priesthoods. Of the Melchizedek priesthood, three presiding high priests, chosen by the body, appointed and ordained to that office, and upheld by the confidence, faith and prayer of the church, form a quorum of the presidency of the church. The twelve travelling counsellors are called to be the twelve apostles, or special witnesses of the name of Christ, in all the world: thus differing from other officers in the church in the duties of their calling. And they form a quorum equal in authority and power to the three presidents, previously mentioned. The seventy are also called to preach the gospel, and to be especial witnesses unto the Gentiles and in all the world. Thus differing from other officers in the church in the duties of their calling: and they form a quorum equal in authority to that of the twelve especial witnesses or apostles, just named. And every descision made by either of these quorums, must be by the unanimous voice of the same; that is, every member in each quorum must be agreed to its decisions in order to make their decisions of the same power or validity one with the other. -[A majority may form a quorum when circumstances render it impossible to be otherwise.]- Unless this is the case, their decisions are not entitled to the same blessings which the decisions of a quorum of three presidents were anciently, who were ordained after the order of Melchizedek, and were righteous and holy men. The decisions of these quorums, or either of them are to be made in all righteousness; in holiness and lowliness of heart; meekness and long suffering; and in faith and virtue and knowledge; temperance, patience, godliness brotherly kindness and charity, because the promise is, if these things abound in them, they shall not be unfruitful in the knowledge of the Lord. And in case that any decision, of these quorums, is made in unrighteousness, it may be brought before a general assembly of the several quorums which constitute the spiritual authorities of the church, otherwise there can be no appeal from their decision.

				12 The twelve are a travelling, presiding high council, to officiate in the name of the Lord, under the direction of the presidency of the church, agreeably to the institution of heaven; to build up the church, and regulate all the affairs of the same, in all nations: first unto the Gentiles, and secondly unto the Jews.

				13 The seventy are to act in the name of the Lord, under the direction of the twelve, or the travelling high council, in building up the church and regulating all the affairs of the same, in all nations: first unto the Gentiles and then to the Jews:---the twelve being sent out, holding the keys, to open the door by the proclamation of the gospel of Jesus Christ; and first unto the Gentiles and then unto the Jews.

				14 The standing high councils, at the stakes of Zion, form a quorum equal in authority, in the affairs of the church, in all their decisions, to the quorum of the presidency, or to the travelling high council.

				15 The high council in Zion, forms a quorum equal in authority, in the affairs of the church, in all their decisions, to the councils of the twelve at the stakes of Zion.

				16 It is the duty of the travelling high council to call upon the seventy, when they need assistance, to fill the several calls for preaching and administering the gospel, in stead of any others.

				17 It is the duty of the twelve in all large branches of the church, to ordain evangelical ministers, as they shall be designated unto them by revelation.

				18 The order of this priesthood was confirmed to be handed down from father to son, and rightly belongs to the literal descendants of the chosen seed, to whom the promises were made. This order was instituted in the days of Adam, and came down by lineage in the following manner:

				19 From Adam to Seth, who was ordained by Adam at the age of 69 years, and was blessed by him three years previous to his (Adam's) death, and received the promise of God by his father, that his posterity should be the chosen of the Lord, and that they should be preserved unto the end of the earth, because he -[Seth]- was a perfect man, and his likeness was the express likeness of his father's, insomuch that he seemed to be like unto his father in all things; and could be distinguished from him only by his age.

				20 Enos was ordained at the age of 134 years, and four months, by the hand of Adam.

				21 God called upon Cainan in the wilderness, in the fortieth year of his age, and he met Adam in journeying to the place Shedolamak: he was eighty seven years old when he received his ordination.

				22 Mahalaeel was 496 years and seven days old when he was ordained by the hand of Adam, who also blessed him.

				23 Jared was 200 years old when he was ordained under the hand of Adam, who also blessed him.

				24 Enoch, was 25 years old when he was ordained under the hand of Adam, and he was 65 and Adam blessed him---and he saw the Lord: and he walked with him, and was before his face continually: and he walked with God 365 years: making him 430 years old when he was translated.

				25 Methuselah was 100 years old when he was ordained under the hand of Adam.

				26 Lamech was 32 years old when he was ordained under the hand of Seth.

				27 Noah was 10 years old when he was ordained under the hand of Methuselah.

				28 Three years previous to the death of Adam, he called Seth, Enos, Cainan, Mahalaleel, Jared, Enoch and Methuselah, who were all high priests, with the residue of his posterity, who were righteous, into the valley of Adam-ondi-ahman, and there bestowed upon them his last blessing. And the Lord appeared unto them, and they rose up and blessed Adam, and called him Michael, the Prince, the Archangel. And the Lord administered comfort unto Adam, and said unto him, I have set thee to be at the head: a multitude of nations shall come of thee; and thou art a prince over them for ever.

				29 And Adam stood up in the midst of the congregation, and notwithstanding he was bowed down with age, being full of the Holy Ghost, predicted whatsoever should befall his posterity unto the latest generation. These things were all written in the book of Enoch, and are to be testified of in due time.

				30 It is the duty of the twelve, also, to ordain and set in order all the other officers of the church, agreeably to the revelation which says:

				31 To the church of Christ in the land of Zion, in addition to the church laws, respecting church business: Verily, I say unto you, says the Lord of hosts, There must needs be presiding elders, to preside over those who are of the office of an elder; and also priests, to preside over those who are of the office of a priest; and also teachers to preside over those who are of the office of a teacher, in like manner; and also the deacons: wherefore, from deacon to teacher, and from teacher to priest, and from priest to elder, severally as they are appointed, according to the covenants and commandments of the church; then comes the high priesthood, which is the greatest of all. Wherefore, it must needs be that one be appointed, of the high priesthood, to preside over the priesthood; and he shall be called president of the high priesthood of the church, or, in other words, the presiding high priest over the high priesthood of the church. From the same comes the administering of ordinances and blessings upon the church, by the laying on of the hands.

				32 Wherefore the office of a bishop is not equal unto it; for the office of a bishop is in administering all temporal things: nevertheless, a bishop must be chosen from the high priest hood, unless he is a literal descendant of Aaron; for unless he is a literal descendant of Aaron he cannot hold the keys of that priesthood. Nevertheless, a high priest, that is after the order of Melchizedek, may be set apart unto the ministering of temporal things, having a knowledge of them by the Spirit of truth, and also to be a judge in Israel, to do the business of the church to sit in judgment upon transgressors, upon testimony, as it shall be laid before him, according to the laws, by the assistance of his counsellors, whom he has chosen, or will chose among the elders of the church. This is the duty of a bishop who is not a literal descendant of Aaron, but has been ordained to the high priesthood after the order of Melchizedek.

				33 Thus shall he be a judge, even a common judge among the inhabitants of Zion, or in a stake of Zion, or in any branch of the church where he shall be set apart unto this ministry, until the borders of Zion are enlarged, and it becomes necessary to have other bishops, or judges in Zion, or elsewhere: and inasmuch as there are other bishops appointed they shall act in the same office.

				34 But a literal descendant of Aaron has a legal right to the presidency of this priesthood, to the keys of this ministry, to act in the office of bishop independently, without counsellors, except in a case where a president of the high priesthood, after the order of Melchizedek, is tried; to sit as a judge in Israel.---And the decision of either of these councils, agreeably to the commandment which says;

				35 Again, verily, I say unto you: The most important business of the church, and the most difficult cases of the church, inasmuch as there is not satisfaction upon the decision of the bishop, or judges, it shall be handed over and carried up unto the council of the church, before the presidency of the high priesthood; and the presidency of the council of the high priesthood shall have power to call other high priests, even twelve, to assist as counsellors; and thus the presidency of the high priesthood, and its counsellors shall have power to decide upon testimony according to the laws of the church. And after this decision it shall be had in remembrance no more before the Lord; for this is the highest council of the church of God, and a final decision upon controversies, in spiritual matters.

				36 There is not any person belonging to the church, who is exempt from this council of the church.

				37 And inasmuch as a president of the high priesthood shall transgress, he shall be had in remembrance before the common council of the church, who shall be assisted by twelve counsellors of the high priesthood; and their decision upon his head shall be an end of controversy concerning him. Thus, none shall be exempted from the justice and the laws of God; that all things may be done in order and in solemnity, before him, according to truth and righteousness.

				38 And again, verily I say unto you, the duty of a president over the office of a deacon, is to preside over twelve deacons, to sit in council with them, and to teach them their duty---edifying one another, as it is given according to the covenants.

				39 And also the duty of the president over the office of the teachers, is to preside over twenty four of the teachers, and to sit in council with them---teaching them the duties of their office, as given in the covenants.

				40 Also the duty of the president over the priesthood of Aaron, is to preside over forty eight priests, and sit in council with them, to teach them the duties of their office, as is given in the covenants. This president is to be a bishop; for this is one of the duties of this priesthood.

				41 Again, the duty of the president over the office of elders is to preside over ninety six elders, and to sit in council with them, and to teach them according to the covenants. This presidency is a distinct one from that of the seventy, and is designed for those who do not travel into all the world.

				42 And again, the duty of the president of the office of the high priesthood is to preside over the whole church, and to be like unto Moses. Behold, here is wisdom---yea, to be a seer, a revelator, a translator and a prophet---having all the gifts of God which he bestows upon the head of the church.

				43 And it is according to the vision, showing the order of the seventy, that they should have seven presidents to preside over them, chosen out of the number of the seventy, and the seventh president of these presidents is to preside over the six; and these seven presidents are to choose other seventy besides the first seventy, to whom they belong, and are to preside over them; and also other seventy until seven times seventy, if the labor in the vineyard of necessity requires it. And these seventy are to be travelling ministers unto the Gentiles, first, and also unto the Jews, whereas other offices of the church who belong not unto the twelve neither to the seventy, are not under the responsibility to travel among all nations, but are to travel as their circumstances shall allow, notwithstanding they may hold as high and responsible offices in the church.

				44 Wherefore, now let every man learn his duty, and to act in the office in which he is appointed, in all diligence. He that is slothful shall not be counted worthy to stand, and he that learns not his duty and shows himself not approved, shall not be counted worthy to stand; even so. Amen.

		% -----
		\subsubsection{Letters to John Burk, Sally Waterman Phelps, and Almira Mack Scobey, 1--2 June 1835}
		% -----
			\label{EMS-1835-JS-1-2-JUN-1835}
			% \label{EMS-1830-JS-DAY-3LETTERMONTH-YEAR}
			
			\begin{description}
				\item[Print Sources]
			\end{description}

				\begin{tabular}{p{0.5in}p{1in}p{0.5in}p{1in}}
				JSP	 	& D4:326--339		& JSR 		& --- \\
				DC	 	& --- 				& HC 		& 2:228--231 \\
				HDDC 	& ---				& TCHC 		& 2:234 \\
				TPJS 	& 76--78			& 	 		&  \\
				\end{tabular}
				% HDDC = woodford's DC diss; JSR = Marquardt's JS Revelations; TCHC = Vogel HC
			
			\begin{description}
				\item[Online Access] \href{https://www.josephsmithpapers.org/paper-summary/letters-to-john-burk-sally-waterman-phelps-and-almira-mack-scobey-1-2-june-1835/1}{Here}
				% \item[Analysis] \hyperref[XX]{Here}
			\end{description}
			
			\begin{description}
				\item[Background]
			\end{description}
			
				% It might also be useful to give a two sentence context of the period: e.g., "This speech was given in Nauvoo to a small congregation one month prior to the destruction of the Nauvoo Expositor. These and other tensions may have been on the prophet's mind when he gave the speech."
			
			\begin{description}
				\item[Original Source]
			\end{description}
			
				% brief description of original source material, with reference to JSP or somewhere else for more info
				% Background material: it might be helpful to have a note that says whether a given document was presented as a speech, was written in a journal, etc.
				% Also a comment here about how reliable the source might be, or what shape it is in, if there are large omissions or parts that are hard to understand, and so on.
			
			\begin{description}
				\item[Text]
			\end{description}
			
					\hypertarget{EMS-1835-JS-1-2-JUN-1835-a}%
				According to
					\marginnote{a}%
				the order of the kingdom begun in the last days, to prepare men for the rest of the Lord, the elders in Zion or <in> her immediate region have no authority, nor right to medelle [meddle] with her affairs to regulate, or \uline{even} \uline{hold} \uline{any} \uline{courts}. The high council has been organized expressly <to> administer \sout{over} <in> all her spiritual affairs; And the bishop and his council are set over her temporal matters: so thus the elders' acts are null and void: Now the Lord wants the tares and wheat to grow to gether <while \sout{and} in an organized states,>; Zion must be redeemed with Judgments and her Converts by righteousness: Every elder that can, after he has provided for his family (if any he has) and paid his debts must go forth and clear his skirts from the blood of this generation: while they are there, instead of holding courts to stop communion, or any thing <else,> let every one laber to bind up the broken hearted; reclaim the wanderer, and persuade back into the Kingdom, such as have been cut off. by encouraging <them> to lay to and work rightousness, and prepare with one heart and one mind to redeem Zion, that goodly land of promise, where the willing and the obedient shall be blessed. Souls are as precious in the sight of God as they ever were, and the elders were never called to drive any down to hell, but \sout{entice} <invite> and persuade men every where to repent--- It is the acceptable year of the Lord. The priests too, should not be <idle> their duties are plain and unless <they> do them, they can not expect to be approved. Righteousness must govern the saints in all things, and <when> the new covenants come forth the priests will learn that great <things> may be expected <at> their hands.
				
					\hypertarget{EMS-1835-JS-1-2-JUN-1835-b}%
				The Teachers
					\marginnote{b}%
				and Deacons are the standing ministers of the church, and in the absence of other officers, they will have great things required at their hands--- They must strengthen the members;--- persuade such as are out of the way to repent, and meekly urge and persuade every one to forgive each other all their trespasses, offences, and sins; Bear and forbear one with another, brethren, for so God does with us. Cease to find fault and learn to do well: Pray <for> your enemies in the church, and curse not <your> foes <with>out: for vengence belongs to God. Know you not, that ``there is more Joy in heaven over \uline{One} sinner that repents, than \sout{over} there is oever <\uline{ninety}> \uline{nine} Just persons, that need no repentance!'' Strive not about the mysteries of the kingdom; to one is given the word of \uline{wisdom}; to another the power of healing difficulties: To every officer we say be \uline{merciful} and you shall find mercy. Your Brethren, who leave their families and go forth to warn the present generation of the great things to <come>, expect great things of those who enjoy the society of the sain[ts] and their families--- They pray that they may be very prayerful and very humble, \sout{will} work\sout{ing} diligently \sout{for} spiritually <\&> temporally for the redemption of Zion, when \uline{all} the pure in heart can return with songs, of everlasting Joy, to enjoy the good of the land of Zion: Brethren, in the name of Jesus we entreat you to live worthy of the blessing which shall <be> heired [inherited] by the faithful after Zion is redeemed.

					\hypertarget{EMS-1835-JS-1-2-JUN-1835-c}%
				To Elder Lyman Wight:
					\marginnote{c}%
				The Lord is not pleased with him because he did not come to Kirtland as is manifest to the Council: The Council fear, that in consequence of self sufficiency and self knowledge, which hertofore at times headed our worthy brother aside from the strict principles of <the gospel of> our Lord Jesus Christ, unless he goes forth into the vineyard, and magnifies his calling according to the commandment, satan will sift him as chaff. Let him be reminded that it is necessary for him to repent, and humble himself before God and go about his master's business; and let those who have been appointed, whose appointments have been sanctioned by the voice of the church regulate the affairs of Zion, Lest he that puts his hand forth to steady the ark, be made an ensample, as was Uzzah in the \sout{David} Days of David...
				
				\footnote{%
					% \label{EMS-1835-JS-1-2-JUN-1835-phelps}%
					In the Phelps part of this letter, note that he says, ``President Smith preached last Sabbath, and I gave him the text: `This is my beloved Son; hear ye him!' He preached one of the greatest sermons I ever heard--- it was about 3$\frac{1}{2}$ hours long--- and unfolded more mysteries than I can write at this time.''
					}%
				...Brother W. W phelps has left a little space for me to occupy and I gladly improve it, I would be glad to see the children of Zion and del[i]ver the <word> of Eternal <life> to them from my own mouth but cannot this year nevertheless the day will come that I shall enjoy this privileege I trust. and we all shall receive an inheritance in the land of refuge which is so much to be desired seeing it is under the direction of the Allmighty therefore let us live faithful before the Lord and it shall be well with us I feel for all the Chilldren of Zion and pray for them in all my prayrs peace be multiplyed unto their redeemtion and favor from God Amen

				\begin{tightright}
				Joseph Smith Jr%
					\footnote{%
						\label{EMS-1835-JS-1-2-JUN-1835-JS}%
						These bits written in JS's hand.
						}
				\end{tightright}
				
		% -----
		\subsubsection{Letter to Church Brethren, 15 June 1835}
		% -----
			\label{EMS-1835-JS-15-JUN-1835}
			% \label{EMS-1830-JS-DAY-3LETTERMONTH-YEAR}
			
			\begin{description}
				\item[Print Sources]
			\end{description}

				\begin{tabular}{p{0.5in}p{1in}p{0.5in}p{1in}}
				JSP	 	& D4:345--347		& JSR 		& --- \\
				DC	 	& --- 				& HC 		& --- \\
				HDDC 	& ---				& TCHC 		& --- \\
				\end{tabular}
				% HDDC = woodford's DC diss; JSR = Marquardt's JS Revelations; TCHC = Vogel HC
			
			\begin{description}
				\item[Online Access] \href{https://www.josephsmithpapers.org/paper-summary/letter-to-church-brethren-15-june-1835/1}{Here}
				% \item[Analysis] \hyperref[XX]{Here}
			\end{description}
			
			\begin{description}
				\item[Background]
			\end{description}
			
				% It might also be useful to give a two sentence context of the period: e.g., "This speech was given in Nauvoo to a small congregation one month prior to the destruction of the Nauvoo Expositor. These and other tensions may have been on the prophet's mind when he gave the speech."
			
			\begin{description}
				\item[Original Source]
			\end{description}
			
				% brief description of original source material, with reference to JSP or somewhere else for more info
				% Background material: it might be helpful to have a note that says whether a given document was presented as a speech, was written in a journal, etc.
				% Also a comment here about how reliable the source might be, or what shape it is in, if there are large omissions or parts that are hard to understand, and so on.
			
			\begin{description}
				\item[Text]
			\end{description}
			
				\begin{tightright}
				Kirtland June 15, 1835
				\end{tightright}

				Dear brethren in the Lord,

				I send you my love and warmest wishes for your prosperity in the great cause of our Redeemer.

				We are now commencing to prepare and print the New Translation, together with all the revelations which God has been pleased to give us in these last days, and as we are in want of funds to go on with so great and glorious a work, brethren <we> want you should donate and loan us all the means or money you can that we may be enable[d] to accomplish the work as a great means towards the salvation of Men.

				My love to my relatives \&c

				\begin{tightright}
				your brother in the bonds of the New Covenant.

				Joseph Smith Jr
				\end{tightright}
				
		% -----
		\subsubsection{Ordination and Blessing of Cyrus Smalling, 30 June 1835}
		% -----
			\label{EMS-1835-JS-30-JUN-1835}
			% \label{EMS-1830-JS-DAY-3LETTERMONTH-YEAR}
			
			\begin{description}
				\item[Print Sources]
			\end{description}

				\begin{tabular}{p{0.5in}p{1in}p{0.5in}p{1in}}
				JSP	 	& D4:349--353		& JSR 		& --- \\
				DC	 	& --- 		& HC 		& --- \\
				HDDC 	& ---		& TCHC 		& --- \\
				\end{tabular}
				% HDDC = woodford's DC diss; JSR = Marquardt's JS Revelations; TCHC = Vogel HC
			
			\begin{description}
				\item[Online Access] \href{https://www.josephsmithpapers.org/paper-summary/ordination-and-blessing-of-cyrus-smalling-30-june-1835/1}{Here}
				% \item[Analysis] \hyperref[XX]{Here}
			\end{description}
			
			\begin{description}
				\item[Background]
			\end{description}
			
				% It might also be useful to give a two sentence context of the period: e.g., "This speech was given in Nauvoo to a small congregation one month prior to the destruction of the Nauvoo Expositor. These and other tensions may have been on the prophet's mind when he gave the speech."
			
			\begin{description}
				\item[Original Source]
			\end{description}
			
				% brief description of original source material, with reference to JSP or somewhere else for more info
				% Background material: it might be helpful to have a note that says whether a given document was presented as a speech, was written in a journal, etc.
				% Also a comment here about how reliable the source might be, or what shape it is in, if there are large omissions or parts that are hard to understand, and so on.
			
			\begin{description}
				\item[Text]
			\end{description}
			
					\hypertarget{EMS-1835-JS-30-JUN-1835-a}%
				[Cyrus Smalling] ordination blessing which he received under the hands of presidents Joseph Smith, jr. Oliver Cowdery and Sidney Rigdon, in Kirtland, Geauga County; Ohio, June 30, 1835, when he was set apart by prophecy to the apostleship of the first Seventy, in the place of Ezra Thayer, who had left the Church and the apostleship.

				Brother [Smalling] in the name of the Lord Jesus we lay our hands upon you, and ordain you to be one of the first Seventy, and to stand in the office of elder Ezra Thayer, who by transgression fell, for thou, by the prayer of faith, has obtained his mission and taken his bishopric. We therefore set thee apart for this mission and ministry, to preach the gospel and testify to this generation of the things which thou hast seen and heard. In the name of the Lord, therefore, thou Shalt go forth, and his angel shall go with thee whithersoever thou goest; for thou shalt go and nothing shall stay thee. Thou shalt preach to nations that thou knowest not of at this time: thou shalt go to the islands of the seas. When thy locks shall be silvered over with age, thou shalt be renewed that thou shalt receive again the vigor of thy youth and yet stand before the Kings of the earth, if thou desirest it.
					\hypertarget{EMS-1835-JS-30-JUN-1835-b}%
				Thou shalt
					\marginnote{b}%
				be a witness for the Lord of the things that thou Knowest until the Lord comes in his glory, and thou shalt see him with thine eyes in the flesh;--- for the desire of thy heart, in righteousness, shall be granted. Thou shalt go to Europe and preach the gospel to the Kings of the earth. Thou shalt stand before bishops, archbishops and lord bishops, and even the pope himself. Thou shalt push together thy thousands: thou shalt proclaim to Catholics, to Presbyterians and all other sects and denominations, and teach bishops, archbichops and popes. Thou shalt be changed in the twinkling of an eye,--- and many other blessings that we cannot describe, which the Lord will show thee. Thou art of the seed of Abraham, of the tribe of Benjamin. In the name of Jesus we seal these blessings upon thy head through thy diligence, faithfulness, prayer of faith and holy desires; even so. Amen.

				\begin{tabular}{p{7cm}p{3cm}p{1cm}}
				 & Wm. W. Phelps,) & \\
				 &  & Clerks. \\
				 & John Whitmer,) & 
				\end{tabular}

				\{Recorded in this book Dec. 21,1835, Oliver Cowdery Recorder.\}
				
		% -----
		\subsubsection{Revelation, circa June 1835 [D\&C 68]}
		% -----
			\label{EMS-1835-JS-CIR-JUN-1835}
			% \label{EMS-1830-JS-DAY-3LETTERMONTH-YEAR}
			
			\begin{description}
				\item[Print Sources]
			\end{description}

				\begin{tabular}{p{0.5in}p{1in}p{0.5in}p{1in}}
				JSP	 	& D4:354--359		& JSR 		& --- \\
				DC	 	& Section 68 		& HC 		& 1:227--229 \\
				HDDC 	& 853--865			& TCHC 		& 1:158--160 \\
				\end{tabular}
				% HDDC = woodford's DC diss; JSR = Marquardt's JS Revelations; TCHC = Vogel HC
			
			\begin{description}
				\item[Online Access] \href{https://www.josephsmithpapers.org/paper-summary/revelation-circa-june-1835-dc-68/1}{Here}
				% \item[Analysis] \hyperref[DC68]{Here}
			\end{description}
			
			\begin{description}
				\item[Background]
			\end{description}
			
				% It might also be useful to give a two sentence context of the period: e.g., "This speech was given in Nauvoo to a small congregation one month prior to the destruction of the Nauvoo Expositor. These and other tensions may have been on the prophet's mind when he gave the speech."
			
			\begin{description}
				\item[Original Source]
			\end{description}
			
				% brief description of original source material, with reference to JSP or somewhere else for more info
				% Background material: it might be helpful to have a note that says whether a given document was presented as a speech, was written in a journal, etc.
				% Also a comment here about how reliable the source might be, or what shape it is in, if there are large omissions or parts that are hard to understand, and so on.
			
			\begin{description}
				\item[Text]
			\end{description}
			
				\begin{tightcenter}
				REVELATIONS.
				\end{tightcenter}

				\textit{A Revelation, given November,} 1831\textit{, to Orson Hyde, Luke Johnson, Lyman Johnson and William E. McLel[l]in. The mind and will of the Lord, as made known by the voice of the Spirit to a conference concerning certain elders: and also certain items, as made known, in addition to the covenants and commandments:---}

				My servant, Orson Hyde, was called, by his ordinance, to proclaim the everlasting gospel, by the spirit of the living God, from people to people, and from land to land, in the congregations of the wicked, in their synagogues, reasoning with and expounding all scriptures unto them: And behold and lo, this is an ensample unto all those who were ordained unto this priesthood, whose mission is appointed unto them to go forth: And this is the ensample unto them, that they shall speak as they are moved upon by the Holy Ghost; and whatsoever they shall speak, when moved upon by the Holy Ghost, shall be scripture; shall be the will of the Lord; shall be the mind of the Lord; shall be the word of the Lord; shall be the voice of the Lord, and the power of God unto salvation; Behold this is the promise of the Lord unto you, O ye my servants: wherefore, be of good cheer, and do not fear, for I the Lord am with you, and will stand by you; and ye shall bear record of me even Jesus Christ, that I am the Son of the living God; that I was; that I am; and that I am to come. This is the word of the Lord unto you my servant, Orson Hyde; and also unto my servant, Luke Johnson, and unto my servant, Lyman Johnson, and unto my servant William E. McLelin; and unto all the faithful elders of my church: Go ye into all the world; preach the gospel to every creature; acting in the authority which I have given you; baptizing in the name of the Father, and of the Son, and of the Holy Ghost; and he that believeth, and is baptized, shall be saved, and he that believeth not shall be damned; and he that believeth shall be blessed with signs following, even as it is written: And unto you it shall be given to know the signs of the times, and the signs of the coming of the Son of man; and of as many as the Father shall bear record, to you it shall be given power to seal them up unto eternal life: Amen.

				And now concerning the items in addition to the covenants and commandments, they are these; there remaineth hereafter, in the due time of the Lord, other bishops to be set apart unto the church to minister even according to the first: Wherefore they shall be high priests who are worthy, and they shall be appointed by the first presidency of the Melchisedek priesthood, except they be literal descendants of Aaron; and if they be literal descendants of Aaron, they have a legal right to the bishopric, if they are the first born among the sons of Aaron: for the first born holds the right of presidency over this priesthood, and the keys or authority of the same. No man has a legal right to this office, to hold the keys of this priesthood, except he be a literal descendant and the first born of Aaron: but as a high priest of the Melchizedek priesthood, has authority to officiate in all the lesser offices, he may officiate in the office of bishop when no literal descendant of Aaron can be found; provided he is called and set apart, and ordained unto this power under the hands of the first presidency of the Melchizedek priesthood. And a literal descendant of Aaron, also, must be designated by this presidency, and found worthy, and annointed, and ordained under the hands of this presidency, otherwise they are not legally authorized to officiate in their priesthood: but by virtue of the decree concerning their right of the priesthood descending from father to son, they may claim their annointing, if at any time they can prove their lineage, or do ascertain it by revelation from the Lord under the hands of the above named presidency.

				And again, no bishop or high priest, who shall be set apart for this ministry, shall be tried or condemned for any crime save it be before the first presidency of the church; and inasmuch as he is found guilty before this presidency, by testimony that cannot be impeached, he shall be condemned, and if he repents he shall be forgiven, according to the covenants and commandments of the church.

				And again, inasmuch as parents have children in Zion, or in any of her stakes which are organized, that teach them not to understand the doctrine of repentance; faith in Christ the Son of the living God; and of baptism and the gift of the Holy Ghost by the laying on of the hands, when eight years old: the sin be upon the head of the parents, for this shall be a law unto the inhabitants of Zion, or in any of her stakes which are organized, and their children shall be baptized for the remission of their sins when eight years old, and receive the laying on of the hands: and they shall also teach their children to pray, and to walk uprightly before the Lord. And the inhabitants of Zion shall also observe the Sabbath day to keep it holy. And the inhabitants of Zion, also, shall remember their labors, inasmuch as they are appointed to labor, in all faithfulness, for the idler shall be had in remembrance before the Lord. Now I the Lord am not well pleased with the inhabitants of Zion, for there are idlers among them; and their children are also growing up in wickedness: They also seek not earnestly the riches of eternity, but their eyes are full of greediness. These things ought not to be, and must be done away from among them: wherefore let my servant Oliver Cowdery, carry these sayings unto the land of Zion. And a commandment I give unto them, that he that observeth not his prayers before the Lord in the season thereof, let him be had in remembrance before the judge of my people. These sayings are true and faithful: wherefore transgress them not, neither take therefrom. Behold I am Alpha and Omega, and I come quickly: Amen.
				
		% -----
		\subsubsection{Minutes, 14 July 1835}
		% -----
			\label{EMS-1835-JS-14-JUL-1835}
			% \label{EMS-1830-JS-DAY-3LETTERMONTH-YEAR}
			
			\begin{description}
				\item[Print Sources]
			\end{description}

				\begin{tabular}{p{0.5in}p{1in}p{0.5in}p{1in}}
				JSP	 	& D4:365--368		& JSR 		& --- \\
				DC	 	& --- 				& HC 		& 2:236--237 \\
				HDDC 	& ---				& TCHC 		& 2:240--241 \\
				\end{tabular}
				% HDDC = woodford's DC diss; JSR = Marquardt's JS Revelations; TCHC = Vogel HC
			
			\begin{description}
				\item[Online Access] \href{https://www.josephsmithpapers.org/paper-summary/minutes-14-july-1835/1}{Here}
				% \item[Analysis] \hyperref[XX]{Here}
			\end{description}
			
			\begin{description}
				\item[Background]
			\end{description}
			
				% It might also be useful to give a two sentence context of the period: e.g., "This speech was given in Nauvoo to a small congregation one month prior to the destruction of the Nauvoo Expositor. These and other tensions may have been on the prophet's mind when he gave the speech."
			
			\begin{description}
				\item[Original Source]
			\end{description}
			
				% brief description of original source material, with reference to JSP or somewhere else for more info
				% Background material: it might be helpful to have a note that says whether a given document was presented as a speech, was written in a journal, etc.
				% Also a comment here about how reliable the source might be, or what shape it is in, if there are large omissions or parts that are hard to understand, and so on.
			
			\begin{description}
				\item[Text]
			\end{description}
			
				...President J. Smith Jun\supers{r.} addressed the council in points of duty such as observing covenants \&c. President Smith testified to the truth of the above covenant. President O. Cowdery testified that he himself framed the above covenant...
				
		% -----
		\subsubsection{Letter to Sally Waterman Phelps, 20 July 1835}
		% -----
			\label{EMS-1835-JS-20-JUL-1835}
			% \label{EMS-1830-JS-DAY-3LETTERMONTH-YEAR}
			
			\begin{description}
				\item[Print Sources]
			\end{description}

				\begin{tabular}{p{0.5in}p{1in}p{0.5in}p{1in}}
				JSP	 	& D4:368--371		& JSR 		& --- \\
				DC	 	& --- 				& HC 		& --- \\
				HDDC 	& ---				& TCHC 		& --- \\
				\end{tabular}
				% HDDC = woodford's DC diss; JSR = Marquardt's JS Revelations; TCHC = Vogel HC
			
			\begin{description}
				\item[Online Access] \href{https://www.josephsmithpapers.org/paper-summary/letter-to-sally-waterman-phelps-20-july-1835/1}{Here}
				% \item[Analysis] \hyperref[XX]{Here}
			\end{description}
			
			\begin{description}
				\item[Background]
			\end{description}
			
				% It might also be useful to give a two sentence context of the period: e.g., "This speech was given in Nauvoo to a small congregation one month prior to the destruction of the Nauvoo Expositor. These and other tensions may have been on the prophet's mind when he gave the speech."
			
			\begin{description}
				\item[Original Source]
			\end{description}
			
				% brief description of original source material, with reference to JSP or somewhere else for more info
				% Background material: it might be helpful to have a note that says whether a given document was presented as a speech, was written in a journal, etc.
				% Also a comment here about how reliable the source might be, or what shape it is in, if there are large omissions or parts that are hard to understand, and so on.
			
			\begin{description}
				\item[Text]
			\end{description}
			
					\hypertarget{EMS-1835-JS-20-JUL-1835-a}%
				Sister [Sally Waterman] phelps,
					\marginnote{a}%
				as Brother [William W.] phelps has given me an opertunity in his \sout{paper} letter, I gladly improve it, how far it may be to your edifycation I know not, but trust you will reciev[e] it in a maner that shall answer my intention in so doing--- for my intention is to give you a word of consolation to streng[th]en you in the absence of your most worthy Companion; and--- husband whose Merits and experiance and acquirements, but few can compete with in this generation and fewer I fear will ever appretiate the worth of such men; men upon whom God, in his wi[s]dom hath bestowed gifts, that duly qualify them to lead men in the way of life and salvation.
					\hypertarget{EMS-1835-JS-20-JUL-1835-b}%
				I
					\marginnote{b}%
				consider in some degree how great a trial you must have in this seperation, but I think I may safely say, that you may rest with a firm reliance that God will so order it that you may not be seperated only but for a short season, and then <your joy will be full,> and if faithful <he will> return and teach you things that have been hid from the wise and prudent, hiden things of old times, as Moses said in Deut. 33\supers{d} <chap> 19\supers{th} verse for they shall suck of the abundance of [the] [se]as and of the treasures hid in the sand. Some of the[se] things have begun [to] come forth therefore Lift up your heart and be glad; be comforted be faithful: pray in faith; and in fine say to all the saints in Zion as the Lord God liveth the redemtion of Zion is nigh at hand, and we shall live to see it
				
				\begin{tightright}
				J. Sm[i]th Jr...
				\end{tightright}
				
		% -----
		\subsubsection{Letter to Quorum of the Twelve, 4 August 1835}
		% -----
			\label{EMS-1835-JS-4-AUG-1835}
			% \label{EMS-1830-JS-DAY-3LETTERMONTH-YEAR}
			
			\begin{description}
				\item[Print Sources]
			\end{description}

				\begin{tabular}{p{0.5in}p{1in}p{0.5in}p{1in}}
				JSP	 	& D4:371--378		& JSR 		& --- \\
				DC	 	& --- 				& HC 		& 2:240--241 \\
				HDDC 	& ---				& TCHC 		& 2:245--246 \\
				\end{tabular}
				% HDDC = woodford's DC diss; JSR = Marquardt's JS Revelations; TCHC = Vogel HC
			
			\begin{description}
				\item[Online Access] \href{https://www.josephsmithpapers.org/paper-summary/letter-to-quorum-of-the-twelve-4-august-1835/1}{Here}
				% \item[Analysis] \hyperref[XX]{Here}
			\end{description}
			
			\begin{description}
				\item[Background]
			\end{description}
			
				% It might also be useful to give a two sentence context of the period: e.g., "This speech was given in Nauvoo to a small congregation one month prior to the destruction of the Nauvoo Expositor. These and other tensions may have been on the prophet's mind when he gave the speech."
			
			\begin{description}
				\item[Original Source]
			\end{description}
			
				% brief description of original source material, with reference to JSP or somewhere else for more info
				% Background material: it might be helpful to have a note that says whether a given document was presented as a speech, was written in a journal, etc.
				% Also a comment here about how reliable the source might be, or what shape it is in, if there are large omissions or parts that are hard to understand, and so on.
			
			\begin{description}
				\item[Text]
			\end{description}
			
					\hypertarget{EMS-1835-JS-4-AUG-1835-a}%
				...We further
					\marginnote{a}%
				inform the twelve, that as far as we can learn from the churches through which you have traveled, that you have set yourselves up as an independant counsel subject to no authority of the church--- a kind of out laws. This impression is wrong, and will if presisted in, bring down the wrath and indignation upon your heads. The other ten are directed to procceed on and finish the conferences, and the two may act their own Judgment, whether to proceed or return...
				
				...We wish you to understand, that your duty requires you to seek first the kingdom of Heaven and its righteousness, that is--- attend to the first things \uline{first}, and then all things will be added, and that complaint about your families will be less frequent---22 Dont preach yourselves crucified for your wives sake, but remembr that Christ was crucified, and you are sent out to be special witnesses of this thing. Men do not wish to hear these \uline{little} things, for there is no salvation in them, but there is in the other.
				
					\hypertarget{EMS-1835-JS-4-AUG-1835-b}%
				Let the
					\marginnote{b}%
				hands of the ten be strengthened, and let them go forth in the name of the Lord, in the power of their mission, giving diligent heed to the direction of the Holy Spirit--- we say, be strong in the Lord and in the power of his might, for great things await you, and great blessings are in store for you. Let the power of the two be upon the seventy until the two make full satisfaction; for the seventy shall be blessed and are blessed. That man who presumes to speak evil of the dignities which God has set in his church, to his family, or to any body else, shall be cursed in his generation. Remember the 109 Psalm His bishopric shall be taken from him unless he speedily repents. Be it known that God is God and when he speaks let all the congregation say: Amen...
				
		% -----
		\subsubsection{Letter to Church Officers in Clay County, Missouri, 31 August 1835}
		% -----
			\label{EMS-1835-JS-31-AUG-1835}
			% \label{EMS-1830-JS-DAY-3LETTERMONTH-YEAR}
			
			\begin{description}
				\item[Print Sources]
			\end{description}

				\begin{tabular}{p{0.5in}p{1in}p{0.5in}p{1in}}
				JSP	 	& D4:399--407		& JSR 		& --- \\
				DC	 	& --- 				& HC 		& --- \\
				HDDC 	& ---				& TCHC 		& --- \\
				TPJS 	& 78--79			& 	 		&  \\
				\end{tabular}
				% HDDC = woodford's DC diss; JSR = Marquardt's JS Revelations; TCHC = Vogel HC
			
			\begin{description}
				\item[Online Access] \href{https://www.josephsmithpapers.org/paper-summary/letter-to-church-officers-in-clay-county-missouri-31-august-1835/1}{Here}
				% \item[Analysis] \hyperref[XX]{Here}
			\end{description}
			
			\begin{description}
				\item[Background]
			\end{description}
			
				% It might also be useful to give a two sentence context of the period: e.g., "This speech was given in Nauvoo to a small congregation one month prior to the destruction of the Nauvoo Expositor. These and other tensions may have been on the prophet's mind when he gave the speech."
			
			\begin{description}
				\item[Original Source]
			\end{description}
			
				% brief description of original source material, with reference to JSP or somewhere else for more info
				% Background material: it might be helpful to have a note that says whether a given document was presented as a speech, was written in a journal, etc.
				% Also a comment here about how reliable the source might be, or what shape it is in, if there are large omissions or parts that are hard to understand, and so on.
			
			\begin{description}
				\item[Text]
			\end{description}
			
					\hypertarget{EMS-1835-JS-31-AUG-1835-a}%
				Kirtland Aug. 31, 1835.
					\marginnote{a}

				The Presidency of Kirtland and Zion say that the Lord has manifested by revelation of his spirit that the high preeist, teachers, Priests and deacons, or in other words all the officers in the land of Clay Co. Mo. belonging to the church are more or less in transgression, because they have not enjoyed the Spirit of God sufficiently to be able to comprehend their duties respecting themselves, and the welfare of Zion. Thereby having been left to act, in a manner that is detrimental to the interest, and also a hindrance, to the redemption of Zion.

					\hypertarget{EMS-1835-JS-31-AUG-1835-b}%
				Now if
					\marginnote{b}%
				they will be wise, they will humble themselves in a peculiar manner that God may open the eyes of their understanding, It will be clearly manifest that the design and purposes of the Almighty; are with regard to them and the children of Zion: that they should let the high counsel which is appointed of God and ordained for that purpose, <make and> regulate all the affairs of Zion: and that it is the will of God, that her children should stand still, and see the salvation of her redemption; and the officers of the church should go forth, inasmuch as they can leave their families in comfortable circumstances; and gather up the saints, even the strength of the Lords house. And those who cannot go forth consistently with the will of God their circumstances prev[e]nting them; remain in deep humility: and in as much, [as] th[e]y do any thing confine themselves to teaching the first principles of the Gospel: not endeavoring to institute regulations or laws for Zion, without having been appointed of God.

					\hypertarget{EMS-1835-JS-31-AUG-1835-c}%
				Now we
					\marginnote{c}%
				say there is no need of ordaining in Zion, or appointing any more officers: but let all those that are ordained magnify themselves before the Lord: by going into the vineyard and cleansing their garments from the blood of this generation. It is one thing to be ordained to preach the gospel, and to push the people together to Zion, and it is another thing to be annointed to lay the foundation and build up the City of Zion, and execute her laws. Therefore it is certain that many of the Elders have come under great condemnation, in endeavoring to steady the ark of God, in a place where th[e]y have not been sent.

				The high counsel and bishops court have been established to do the business of Zion, and her children are not bound, to accknowledge any of those who feel disposed to run to Zion and set themselves to be their rulers. Let not her children be duped in this way, but let them prove those who say they are apostles, and are not.

					\hypertarget{EMS-1835-JS-31-AUG-1835-d}%
				The Elders
					\marginnote{d}%
				have no right to regulate Zion, but they have a right to preach the gospel. They will all do well to repent and humble themselves, and all the church, and also we, ourselves receive the admonition and so now endeavor and pray to this end. When the children of Zion are \sout{in} stranger[s] in a strange land their harps must be hu[n]g upon the willows: and they cannot sing the songs of Zion: but should mourn and not dance. Therefore brethren, it remains for all such to be exercised with prayer, and continual suplication, until Zion is redeemed. We realize the situation that all the brethren and sisters must be in, being deprived of their spiritual privileges, which are enjoyed by those who sit in heavenly places in Christ Jesus; where there are no mobs to rise up and bind their consciences. Nevertheless, it is wisdom that the church should make but little or no stir in that region, and cause as little excitement as posible and endure their afflictions patiently until the time appointed--- and the Governor of Mo. fulfils his promise in setting the church over upon their own lands. We would suggest an idea that it would be wisdom for all the members of the church on the return of the Bishop [Edward Partridge], to make known to him their names places of residince \&c. that it may be known where they all are when the Governor shall give directions for you to be set over on your lands

					\hypertarget{EMS-1835-JS-31-AUG-1835-e}%
				Again
					\marginnote{e}%
				it is the will of the Lord, that the church should attend to their communion on the sabbath day, and let them remember the commandment which says ``talk not of Judgment'' we are commanded not to give the childrens bread unto the dogs; neither cast our pearl[s] before swine, least they trample them under their feet, and turn again and rend you. Therefore--- let us be wise in all things, and keep all the commandments of God, that our salvation may be sure: having our armour ready and prepared against the time appointed; and having on the whole armour of righteousness, we may be able to stand in that trying day. We say also that if there are any doors open for the Elders to preach the first principles of the gospel: let them not keep silence: rail not against the sects, neither talk against their tenets. But preach Christ and him crucified. love to God, and love to man, observing always to make mention of our republican principles, thereby if posible, we may allay the prejudice of the people, be meek and lowly of heart, and the Lord God of our fathers shall be with you for evermore

				Amen.

				Sanctioned and signed by the Presidents

				\begin{tightright}
				Joseph Smith Jr.

				Oliver Cowdery

				Sidney Rigdon

				F[rederick] G. Williams

				W[illiam] W. Phelps

				John Whitmer.
				\end{tightright}

					\hypertarget{EMS-1835-JS-31-AUG-1835-f}%
				P.S. Br Hesekiah [Hezekiah] Peck
					\marginnote{f}

				We remember your familly, with all the fi[r]st families of the church, who first embraced the truth, we remember your losses and sorrows our first ties are not broken, we participate with you in the evil as well as the good, in the sorrows as well as the Joys, our union we trust is stronger than death, and shall never be severed. Remember us unto all who believe in the fulness of the gospel of our Lord and Saviour Jesus Christ. We hereby authorize you Hezekiah Peck, our beloved brother to read this epistle and \sout{communication} communicate it unto all the brotherhood in all that region of Country, Dictated by me your unworthy brother, and fellow laborer in the testimony of the book of Mormon. Signed by my own hand in the token of the everlasting covenant.

				\begin{tightcenter}
				Joseph Smith Jr.
				\end{tightcenter}
				
		% -----
		\subsubsection{Declaration on Government and Law, circa August 1835 [D\&C 134]}
		% -----
			\label{EMS-1835-JS-CIR-AUG-1835-GOV}
			% \label{EMS-1830-JS-DAY-3LETTERMONTH-YEAR}
			
			\begin{description}
				\item[Print Sources]
			\end{description}

				\begin{tabular}{p{0.5in}p{1in}p{0.5in}p{1in}}
				JSP	 	& D4:408--412		& JSR 		& --- \\
				DC	 	& Section 27 		& HC 		& 1:106--108 \\
				HDDC 	& 393--403			& TCHC 		& 1:75--76 \\
				\end{tabular}
				% HDDC = woodford's DC diss; JSR = Marquardt's JS Revelations; TCHC = Vogel HC
			
			\begin{description}
				\item[Online Access] \href{https://www.josephsmithpapers.org/paper-summary/appendix-4-declaration-on-government-and-law-circa-august-1835-dc-134/1}{Here}
				% \item[Analysis] \hyperref[DC27]{Here}
			\end{description}
			
			\begin{description}
				\item[Background]
			\end{description}
			
				JS's involvement with the production of this document isn't clear, but he did later endorse it.
			
				% It might also be useful to give a two sentence context of the period: e.g., "This speech was given in Nauvoo to a small congregation one month prior to the destruction of the Nauvoo Expositor. These and other tensions may have been on the prophet's mind when he gave the speech."
			
			\begin{description}
				\item[Original Source]
			\end{description}
			
				% brief description of original source material, with reference to JSP or somewhere else for more info
				% Background material: it might be helpful to have a note that says whether a given document was presented as a speech, was written in a journal, etc.
				% Also a comment here about how reliable the source might be, or what shape it is in, if there are large omissions or parts that are hard to understand, and so on.
			
			\begin{description}
				\item[Text]
			\end{description}
			
				\begin{tightcenter}
					SECTION CII.
				
					Of Governments and Laws in General.
				\end{tightcenter}
				
				\textit{That our belief, with regard to earthly governments and laws in general, may not be misinterpreted nor misunderstood, we have thought proper to present, at the close of this volume, our opinion concerning the same.}
				
				1 We believe that Governments were instituted of God for the benefit of man, and that he holds men accountable for their acts in relation to them, either in making laws or administering them, for the good and safety of society.
				
				2 We believe that no Government can exist, in peace, except such laws are framed and held inviolate as will secure to each individual the free exercise of con[s]cience, the right and control of property and the protection of life.

				3 We believe that all Governments necessarily require civil officers and magistrates to enforce the laws of the same, and that such as will administer the law in equity and justice should be sought for and upheld by the voice of the people, (if a Republic,) or the will of the Sovereign.

				4 We believe that religion is instituted of God, and that men are amenable to him and to him only for the exercise of it, unless their religious opinion prompts them to infringe upon the rights and liberties of others; but we do not believe that human law has a right to interfere in prescribing rules of worship to bind the consciences of men, nor dictate forms for public or private devotion; that the civil magistrate should restrain crime, but never control conscience; should punish guilt, but never surpress the freedom of the soul.

				5 We believe that all men are bound to sustain and uphold the respective Governments in which they reside, while protected in their inherent and unalienable rights by the laws of such Governments, and that sedition and rebellion are unbecoming every citizen thus protected, and should be punished accordingly; and that all Governments have a right to enact such laws as in their own judgments are best calculated to secure the public interest, at the same time, however, holding sacred the freedom of conscience.

				6 We believe that every man should be honored in his station: rulers and magistrates as such—being placed for the protection of the innocent and the punishment of the guilty; and that to the laws all men owe respect and deference, as without them peace and harmony would be supplanted by anarchy and terror: human laws being instituted for the express purpose of regulating our interests as individuals and nations, between man and man, and divine laws, given of heaven, prescribing rules on spiritual concerns, for faith and worship, both to be answered by man to his Maker.

				7 We believe that Rulers, States and Governments have a right, and are bound to enact laws for the protection of all citizens in the free exercise of their religious belief; but we do not believe that they have a right, in justice, to deprive citizens of this privilege, or proscribe them in their opinions, so long as a regard and reverence is shown to the laws, and such religious opinions do not justify sedition nor conspiracy.

				8 We believe that the commission of crime should be punished according to the nature of the offence: that murder, treason, robbery, theft and the breach of the general peace, in all respects, should be punished according to their criminality and their tendency to evil among men, by the laws of that Government in which the offence is committed: and for the public peace and tranquility, all men should step forward and use their ability in bringing offenders, against good laws, to punishment.

				9 We do not believe it just to mingle religious influence with civil Government, whereby one religious society is fostered and another proscribed in its spiritual privileges, and the individual rights of its members, as citizens, denied.

				10 We believe that all religious societies have a right to deal with their members for disorderly conduct according to the rules and regulations of such societies, provided that such dealing be for fellowship and good standing; but we do not believe that any religious society has authority to try men on the right of property or life, to take from them this world's goods, or put them in jeopardy either life or limb, neither to inflict any physical punishment upon them,—they can only excommunicate them from their society and withdraw from their fellowship.

				11 We believe that men should appeal to the civil law for redress of all wrongs and grievances, where personal abuse is inflicted, or the right of property or character infringed, where such laws exist as will protect the same; but we believe that all men are justified in defending themselves, their friends and property, and the Government, from the unlawful assaults and encroachments of all persons, in times of exigencies, where immediate appeal cannot be made to the laws, and relief afforded.

				12 We believe it just to preach the gospel to the nations of the earth, and warn the righteous to save themselves from the corruption of the world; but we do not believe it right to interfere with bond-servants, neither preach the gospel to, nor baptize them, contrary to the will and wish of their masters, nor to meddle with, or influence them in the least to cause them to be dissatisfied with their situations in this life, thereby jeopardizing the lives of men: such interference we believe to be unlawful and unjust, and dangerous to the peace of every Government allowing human beings to be held in servitude. [\textit{remainder of page blank}]

		% -----
		\subsubsection{Revelation, circa August 1835 [D\&C 27]}
		% -----
			\label{EMS-1835-JS-CIR-AUG-1835}
			% \label{EMS-1830-JS-DAY-3LETTERMONTH-YEAR}
			
			\begin{description}
				\item[Print Sources]
			\end{description}

				\begin{tabular}{p{0.5in}p{1in}p{0.5in}p{1in}}
				JSP	 	& D4:408--412		& JSR 		& --- \\
				DC	 	& Section 27 		& HC 		& 1:106--108 \\
				HDDC 	& 393--403			& TCHC 		& 1:75--76 \\
				\end{tabular}
				% HDDC = woodford's DC diss; JSR = Marquardt's JS Revelations; TCHC = Vogel HC
			
			\begin{description}
				\item[Online Access] \href{https://www.josephsmithpapers.org/paper-summary/revelation-circa-august-1835-dc-27/1}{Here}
				% \item[Analysis] \hyperref[DC27]{Here}
			\end{description}
			
			\begin{description}
				\item[Background]
			\end{description}
			
				% It might also be useful to give a two sentence context of the period: e.g., "This speech was given in Nauvoo to a small congregation one month prior to the destruction of the Nauvoo Expositor. These and other tensions may have been on the prophet's mind when he gave the speech."
			
			\begin{description}
				\item[Original Source]
			\end{description}
			
				% brief description of original source material, with reference to JSP or somewhere else for more info
				% Background material: it might be helpful to have a note that says whether a given document was presented as a speech, was written in a journal, etc.
				% Also a comment here about how reliable the source might be, or what shape it is in, if there are large omissions or parts that are hard to understand, and so on.
			
			\begin{description}
				\item[Text]
			\end{description}
			
				\begin{tightcenter}
				SECTION L.

				\textit{Revelation given September}, 1830.
				\end{tightcenter}

				1 Listen to the voice of Jesus Christ, your Lord, your God and your Redeemer, whose word is quick and powerful. For behold I say unto you, that it mattereth not what ye shall eat, or what ye shall drink, when ye partake of the sacrament, if it so be that ye do it with an eye single to my glory; remembering unto the Father my body which was laid down for you, and my blood which was shed for the remission of your sins: wherefore a commandment I give unto you, that you shall not purchase wine, neither strong drink of your enemies: wherefore you shall partake of none, except it is made new among you, yea, in this my Father's kingdom which shall be built up on the earth.

				2 Behold this is wisdom in me: wherefore marvel not for the hour cometh that I will drink of the fruit of the vine with you on the earth, and with Moroni, whom I have sent unto you to reveal the book of Mormon, containing the fulness of my everlasting gospel; to whom I have committed the keys of the record of the stick of Ephraim; and also with Elias, to whom I have committed the keys of bringing to pass the restoration of all things, or the restorer of all things spoken by the mouth of all the holy prophets since the world began, concerning the last days: and also John the son of Zacharias, which Zacharias he (Elias) visited and gave promise that he should have a son, and his name should be John, and he should be filled with the spirit of Elias; which John I have sent unto you, my servants, Joseph Smith, jr. and Oliver Cowdery, to ordain you unto this first priesthood which you have received, that you might be called and ordained even as Aaron: and also Elijah, unto whom I have committed the keys of the power of turning the hearts of the fathers to the children and the hearts of the children to the fathers, that the whole earth may not be smitten with a curse: and also, with Joseph, and Jacob, and Isaac, and Abraham your fathers; by whom the promises remain; and also with Michael, or Adam, the father of all, the prince of all, the ancient of days:

				3 And also with Peter, and James, and John, whom I have sent unto you, by whom I have ordained you and confirmed you to be apostles and especial witnesses of my name, and bear the keys of your ministry: and of the same things which I revealed unto them: unto whom I have committed the keys of my kingdom, and a dispensation of the gospel for the last times; and for the fulness of times, in the which I will gather together in one all things both which are in heaven and which are on earth: and also with all those whom my Father hath given me out of the world: wherefore lift up your hearts and rejoice, and gird up your loins, and take upon you my whole armor, that ye may be able to withstand the evil day, having done all ye may be able to stand. Stand, therefore, having your loins girt about with truth; having on the breastplate of righteousness; and your feet shod with the preparation of the gospel of peace which I have sent mine angels to commit unto you, taking the shield of faith wherewith ye shall be able to quench all the fiery darts of the wicked; and take the helmet of salvation, and the sword of my Spirit, which I will pour out upon you, and my word which I reveal unto you, and be agreed as touching all things whatsoever ye ask of me, and be faithful until I come, and ye shall be caught up that where I am ye shall be also. Amen.
				
		% -----
		\subsubsection{Minutes, 19 September 1835}
		% -----
			\label{EMS-1835-JS-19-SEP-1835}
			% \label{EMS-1830-JS-DAY-3LETTERMONTH-YEAR}
			
			\begin{description}
				\item[Print Sources]
			\end{description}

				\begin{tabular}{p{0.5in}p{1in}p{0.5in}p{1in}}
				JSP	 	& D4:422--427		& JSR 		& --- \\
				DC	 	& --- 				& HC 		& 2:277--280 \\
				HDDC 	& ---				& TCHC 		& 2:267--270 \\
				\end{tabular}
				% HDDC = woodford's DC diss; JSR = Marquardt's JS Revelations; TCHC = Vogel HC
			
			\begin{description}
				\item[Online Access] \href{https://www.josephsmithpapers.org/paper-summary/minutes-19-september-1835/1}{Here}
				% \item[Analysis] \hyperref[XX]{Here}
			\end{description}
			
			\begin{description}
				\item[Background]
			\end{description}
			
				% It might also be useful to give a two sentence context of the period: e.g., "This speech was given in Nauvoo to a small congregation one month prior to the destruction of the Nauvoo Expositor. These and other tensions may have been on the prophet's mind when he gave the speech."
			
			\begin{description}
				\item[Original Source]
			\end{description}
			
				% brief description of original source material, with reference to JSP or somewhere else for more info
				% Background material: it might be helpful to have a note that says whether a given document was presented as a speech, was written in a journal, etc.
				% Also a comment here about how reliable the source might be, or what shape it is in, if there are large omissions or parts that are hard to understand, and so on.
			
			\begin{description}
				\item[Text]
			\end{description}
			
					\hypertarget{EMS-1835-JS-19-SEP-1835-a}%
				...Object
					\marginnote{a}%
				of the council stated by President J. Smith Ju\supersu{r} as follows: Some weeks since Elder Jared Carter preached on the sabbath in the church, and some of the brethren found fault with his teachings, and this council was called by him (J. S. Jun\supers{r}) to decide this matter and to see who was in fault...
				
				...President Joseph Smith \supers{Jun\supers{r.}} then arose and said that the decision of his mind is that brother Jared erred in judgement in not understanding what the brethren desired of him when they labored with him. and he erred in spirit when he taught in the church, the things testified of here. And that the hand of the destroyer was laid upon him because he had a rebelious spirit. from the beginning, and the word of the Lord, had been spoken by <my> mouth--- that it should upon him and this council should see it and now that he has been seized by the destroyer comes in fulfilment of his word. and God required him to bear testimony of it before the church and warn them to be careful \& not do as he had done.
					\hypertarget{EMS-1835-JS-19-SEP-1835-b}%
				But
					\marginnote{b}%
				instead of doing this, he said he would prove the Book of Mormon, and one thing or another not being sufficiently humble to deliver just the message that was required, and so he stumbled and could not get the spirit. and the brethren were not edified, and he did not do the thing that God required. but erred in choosing words to communicate his thoughts. Such as commanding the prayers of the church instead of soliciting them: and also of making himself an example for the Church, when it, was only the things which he suffered which were to be as a check upon transgressors. His rebeling against the advice and counsel of the Presidents was the cause of his falling to the hands of the Destroyer, again as he had done before when he rebelled against the counsel that was given him by the Authorities of the church. And that in all this Elder Carter has not designed to do wickedly, but he erred in judgement and deserves reproof. And the decision is that he shall acknowledge his errors on the morrow before the congregation, and say, brethren, I am fully convinced, that I have erred in spirit in my remarks before you, when I spoke here a few sabbaths since, \& now I ask your forgiveness, and if he do this in full faith and is truly humble before God. Then God will bless him abundintly as he hath not been wont to do...\footnote{%
					\label{EMS-1835-JS-19-SEP-1835-background}%
					Note about the circumstances of these remarks.
					}
					
		% -----
		\subsubsection{Record, 22 September 1835}
		% -----
			\label{EMS-1835-JS-22-SEP-1835}
			% \label{EMS-1830-JS-DAY-3LETTERMONTH-YEAR}
			
			\begin{description}
				\item[Print Sources]
			\end{description}

				\begin{tabular}{p{0.5in}p{1in}p{0.5in}p{1in}}
				JSP	 	& J1:61--62		& JSR 		& --- \\
				DC	 	& --- 				& HC 		& --- \\
				HDDC 	& ---				& TCHC 		& --- \\
				PJS 	& 1:98			& 	 		&  \\
				\end{tabular}
				% HDDC = woodford's DC diss; JSR = Marquardt's JS Revelations; TCHC = Vogel HC
			
			\begin{description}
				\item[Online Access] \href{https://www.josephsmithpapers.org/paper-summary/journal-1835-1836/2}{Here}
				% \item[Analysis] \hyperref[XX]{Here}
			\end{description}
			
			\begin{description}
				\item[Background]
			\end{description}
			
				% It might also be useful to give a two sentence context of the period: e.g., "This speech was given in Nauvoo to a small congregation one month prior to the destruction of the Nauvoo Expositor. These and other tensions may have been on the prophet's mind when he gave the speech."
			
			\begin{description}
				\item[Original Source]
			\end{description}
			
				% brief description of original source material, with reference to JSP or somewhere else for more info
				% Background material: it might be helpful to have a note that says whether a given document was presented as a speech, was written in a journal, etc.
				% Also a comment here about how reliable the source might be, or what shape it is in, if there are large omissions or parts that are hard to understand, and so on.
			
			\begin{description}
				\item[Text]
			\end{description}
			
				September 22, 1835. This day Joseph Smith, jr. labored with Oliver Cowdery, in obtaining and writing blessings. We were thronged a part of the time with company, so that our labor, in this thing, was hindered; but we obtained many precious things, and our souls were blessed. O Lord, may thy Holy Spirit be with thy servants forever. Amen.
				
				\textbf{September 23.\supers{th} [22\supers{nd}] This day Joseph Smith, Jr. was at home writing blessings for my most beloved Brotheren <I>, have been hindered by a multitude of visitors but the Lord has blessed our Souls this day. May God grant <to> continue his mercies unto my house, this <night,> \sout{day} For Christ sake. This day my Soul has desired the salvation of Brother Ezra, Thay[e]r. Also Brother Noah, Packard. Came to my house and let the Chappel Committee have one thousand dollers, by loan, for the building the house of the Lord; Oh may God bless him with an hundred fold! even of the <things of> Earth, for this ritious [righteous] act. My heart is full of desire to day, to <be> blessed of the God, of Abraham; with prosperity, untill I will be able to pay all my depts; for it is \sout{my} <the> delight of my soul to <be> honest. Oh Lord that thou knowes right well! help me and I will give to the poor.}
					
		% -----
		\subsubsection{Blessing to David Whitmer, 22 September 1835}
		% -----
			\label{EMS-1835-JS-22-SEP-1835-DW}
			% \label{EMS-1830-JS-DAY-3LETTERMONTH-YEAR}
			
			\begin{description}
				\item[Print Sources]
			\end{description}

				\begin{tabular}{p{0.5in}p{1in}p{0.5in}p{1in}}
				JSP	 	& D4:428--430		& JSR 		& --- \\
				DC	 	& --- 				& HC 		& 2:281 \\
				HDDC 	& ---				& TCHC 		& 2:270 \\
				\end{tabular}
				% HDDC = woodford's DC diss; JSR = Marquardt's JS Revelations; TCHC = Vogel HC
			
			\begin{description}
				\item[Online Access] \href{https://www.josephsmithpapers.org/paper-summary/blessing-to-david-whitmer-22-september-1835/1}{Here}
				% \item[Analysis] \hyperref[XX]{Here}
			\end{description}
			
			\begin{description}
				\item[Background]
			\end{description}
			
				% It might also be useful to give a two sentence context of the period: e.g., "This speech was given in Nauvoo to a small congregation one month prior to the destruction of the Nauvoo Expositor. These and other tensions may have been on the prophet's mind when he gave the speech."
			
			\begin{description}
				\item[Original Source]
			\end{description}
			
				% brief description of original source material, with reference to JSP or somewhere else for more info
				% Background material: it might be helpful to have a note that says whether a given document was presented as a speech, was written in a journal, etc.
				% Also a comment here about how reliable the source might be, or what shape it is in, if there are large omissions or parts that are hard to understand, and so on.
			
			\begin{description}
				\item[Text]
			\end{description}
			
					\hypertarget{EMS-1835-JS-22-SEP-1835-DW-a}%
				Blessed of
					\marginnote{a}%
				the Lord is brother David [Whitmer], for he truly is a faithful friend to mankind, and he should be beloved by all, because of the integrity of his heart. All his words are as steadfast as the pillars of heaven, because truth is his only meditation, and he delighteth in it, and shall rejoice in it forever. The Lord God of Abraham, of Isaac and of Jacob, shall be on his right hand and on his left, and shall go before his face, and shall be his rear ward, and his enemies shall become an easy pray unto him; for behold, he it is, whom the Lord hath appointed to be captain of his host, under the guidance and direction of him who is appointed to say unto the strength of the Lord's house, Go forth and build up the waste places of Zion. A mighty shaft shall he be in the quiver of the Almighty in bringing about the redemption of Zion, and in avenging the wrongs of the innocent.
					\hypertarget{EMS-1835-JS-22-SEP-1835-DW-b}%
				He
					\marginnote{b}%
				shall yet stand upon the land of Zion, from whence he has been driven, and shall find inheritance there; and shall be a ruler in Zion until he is old and well stricken in years. And shall enjoy an abundance of the precious things of the lasting mountains, and shall have part with his brethren in all the good things of this earth, and shall never want a friend. He shall bring down his adversaries under his feet, and shall walk upon their ashes when their names are blotted out. His name shall be a blessing among all nations, and his testimony shall shine as fair as the sun, and as a diamond it shall remain untarnished. There shall not be a spot upon his character while he lives, neither upon his seed after him unto the latest posterity: he shall not be forsaken nor his seed found begging bread. Amen.
				
				\begin{tightright}
				Oliver Cowdery, \uline{Clerk and Recorder}.
				\end{tightright}
				
				Given like the foregoing blessings, by vision, to Joseph Smith, jr. the Seer, September 22, 1835, and recorded Oct. 2, 1835.
				
		% -----
		\subsubsection{Blessing to John Whitmer, 22 September 1835}
		% -----
			\label{EMS-1835-JS-22-SEP-1835-JW}
			% \label{EMS-1830-JS-DAY-3LETTERMONTH-YEAR}
			
			\begin{description}
				\item[Print Sources]
			\end{description}

				\begin{tabular}{p{0.5in}p{1in}p{0.5in}p{1in}}
				JSP	 	& D4:430--433		& JSR 		& --- \\
				DC	 	& --- 		& HC 		& --- \\
				HDDC 	& ---		& TCHC 		& --- \\
				\end{tabular}
				% HDDC = woodford's DC diss; JSR = Marquardt's JS Revelations; TCHC = Vogel HC
			
			\begin{description}
				\item[Online Access] \href{https://www.josephsmithpapers.org/paper-summary/blessing-to-john-whitmer-22-september-1835/1}{Here}
				% \item[Analysis] \hyperref[XX]{Here}
			\end{description}
			
			\begin{description}
				\item[Background]
			\end{description}
			
				% It might also be useful to give a two sentence context of the period: e.g., "This speech was given in Nauvoo to a small congregation one month prior to the destruction of the Nauvoo Expositor. These and other tensions may have been on the prophet's mind when he gave the speech."
			
			\begin{description}
				\item[Original Source]
			\end{description}
			
				% brief description of original source material, with reference to JSP or somewhere else for more info
				% Background material: it might be helpful to have a note that says whether a given document was presented as a speech, was written in a journal, etc.
				% Also a comment here about how reliable the source might be, or what shape it is in, if there are large omissions or parts that are hard to understand, and so on.
			
			\begin{description}
				\item[Text]
			\end{description}
			
				Blessed of the Lord is brother John [Whitmer]: he shall live to a good old age, and his generations after him shall be blessed. He shall live to a good old age, and wax more and more steadfast in the work of the Lord, and he shall make a choice record of Israel unto the memory of his name. He shall truly be chastened wherein he steps aside, yet he shall escape the hand of the Destroyer, and shall never fall by the hands of his enemies: and he shall be avenged of the wrongs he has sustained by them. He shall be blessed, and a means of bringing many blessings upon his father's family. He shall be blessed with abundance of the good things of \sout{the} this earth, with houses and lands, in the fruit of the field, with horses and with asses and with she asses, with gold and with silver, and with precious stones, even the things of the lasting mountains. He shall be made mighty in the hands of his God in bringing to pass the redemption of Zion. The Lord shall be on his right hand and on his left; shall go before his face and shall be his rear ward, and shall turn aside all the shafts of \sout{the} <his> enemies. And he shall go down \sout{to his} <in> honor and with gray hairs, to the grave, and shall be caught up in the cloud to meet the Lord, and shall ever be with him; even so. Amen.
				
				\begin{tightright}
				Oliver Cowdery, \uline{Clerk and Recorder}.
				\end{tightright}
				
				Given September 22, 1835, and recorded October 2, 1835.
				
		% -----
		\subsubsection{Blessing to John Corrill, 22 September 1835}
		% -----
			\label{EMS-1835-JS-22-SEP-1835-JC}
			% \label{EMS-1830-JS-DAY-3LETTERMONTH-YEAR}
			
			\begin{description}
				\item[Print Sources]
			\end{description}

				\begin{tabular}{p{0.5in}p{1in}p{0.5in}p{1in}}
				JSP	 	& D4:433--434		& JSR 		& --- \\
				DC	 	& --- 		& HC 		& --- \\
				HDDC 	& ---		& TCHC 		& --- \\
				\end{tabular}
				% HDDC = woodford's DC diss; JSR = Marquardt's JS Revelations; TCHC = Vogel HC
			
			\begin{description}
				\item[Online Access] \href{https://www.josephsmithpapers.org/paper-summary/blessing-to-john-corrill-22-september-1835/}{Here}
				% \item[Analysis] \hyperref[XX]{Here}
			\end{description}
			
			\begin{description}
				\item[Background]
			\end{description}
			
				% It might also be useful to give a two sentence context of the period: e.g., "This speech was given in Nauvoo to a small congregation one month prior to the destruction of the Nauvoo Expositor. These and other tensions may have been on the prophet's mind when he gave the speech."
			
			\begin{description}
				\item[Original Source]
			\end{description}
			
				% brief description of original source material, with reference to JSP or somewhere else for more info
				% Background material: it might be helpful to have a note that says whether a given document was presented as a speech, was written in a journal, etc.
				% Also a comment here about how reliable the source might be, or what shape it is in, if there are large omissions or parts that are hard to understand, and so on.
			
			\begin{description}
				\item[Text]
			\end{description}
			
					\hypertarget{EMS-1835-JS-22-SEP-1835-JC-a}%
				Blessed of
					\marginnote{a}%
				the Lord is brother Corrill, John Corrill for the hand of the Lord shall be over him, and there are none that surpass him in understanding pertaining to architecture: unto this end hath the Lord raised him up, and he shall yet stand upon the land of Zion and be greatly multiplied in the midst of her: for he shall build the house of the Lord in Zion, and shall do a work that none other man understandeth to do, in building the houses of the Lord. It shall be a gift unto him: he shall be enlightened by visions and dreams, and his soul shall be satisfied with the abundance with which the Lord his God shall bless him: and his seed after him shall be blessed from generation to generation. And his name shall be had in remembrance and be known among all nations; and it shall be said of him among governors, rulers and kings: This man reared the temple of God in Zion.
					\hypertarget{EMS-1835-JS-22-SEP-1835-JC-b}%
				And he
					\marginnote{b}%
				shall be reverenced because of his great sagacity. His life shall be prolonged until he is very old. He shall be accounted among the wise counsellors of Israel, and he shall hear the voice of the Lord his God, and receive counsel by the ministering of angels; and shall be very terrible unto his enemies, and his soul shall be satisfied in seeing the wrath of God poured out upon them to the uttermost. He shall have communion with Abraham, Isaac and Jacob, and with the spirits of just men made perfect, and shall be caught up, whether in the body or out of the body, I cannot tell, God knoweth, he shall be caught up to the third heaven and behold unuterable things. And thus shall the hand of God be upon him all the days of his life, unto his joy and the satisfaction of his soul: and he shall have a crown of eternal life. And these are the blessings that I want should come upon my brother Corrill. Amen.
				
				\begin{tightright}
				Oliver Cowdery, \uline{[C]lerk and Recorder}.
				\end{tightright}
				
				Given September 22, 1835, and recorded October 3, 1835.
				
		% -----
		\subsubsection{Blessing to William W. Phelps, 22 September 1835}
		% -----
			\label{EMS-1835-JS-22-SEP-1835-WP}
			% \label{EMS-1830-JS-DAY-3LETTERMONTH-YEAR}
			
			\begin{description}
				\item[Print Sources]
			\end{description}

				\begin{tabular}{p{0.5in}p{1in}p{0.5in}p{1in}}
				JSP	 	& D4:435--436		& JSR 		& --- \\
				DC	 	& --- 		& HC 		& --- \\
				HDDC 	& ---		& TCHC 		& --- \\
				\end{tabular}
				% HDDC = woodford's DC diss; JSR = Marquardt's JS Revelations; TCHC = Vogel HC
			
			\begin{description}
				\item[Online Access] \href{https://www.josephsmithpapers.org/paper-summary/blessing-to-william-w-phelps-22-september-1835/1}{Here}
				% \item[Analysis] \hyperref[XX]{Here}
			\end{description}
			
			\begin{description}
				\item[Background]
			\end{description}
			
				% It might also be useful to give a two sentence context of the period: e.g., "This speech was given in Nauvoo to a small congregation one month prior to the destruction of the Nauvoo Expositor. These and other tensions may have been on the prophet's mind when he gave the speech."
			
			\begin{description}
				\item[Original Source]
			\end{description}
			
				% brief description of original source material, with reference to JSP or somewhere else for more info
				% Background material: it might be helpful to have a note that says whether a given document was presented as a speech, was written in a journal, etc.
				% Also a comment here about how reliable the source might be, or what shape it is in, if there are large omissions or parts that are hard to understand, and so on.
			
			\begin{description}
				\item[Text]
			\end{description}
			
					\hypertarget{EMS-1835-JS-22-SEP-1835-WP-a}%
				Blessed of
					\marginnote{a}%
				the Lord is brother Phelps, William W. Phelps, for he shall have the desires of his heart in the gift that pertaineth to writing the law of God, and in being an instrument in assisting to lift up an ensign to the nations. And according to the greatness of the desires of his soul, for blessings upon his friends, so shall the blessings of the Lord come upon him to the uttermost. The Lord will chasten him because he taketh honor to himself, and when his soul is greatly humbled he will forsake the evil; then shall the light of the Lord break upon him as at noon-day, and in him shall be no darkness so great is the glory that shall come upon him. And blessed is his name among all nations. He shall have part in that that coucheth beneath: and it shall be revealed unto him things by the hand of the Lord's annointed that have been kept secret from the foundation of the world, concerning the last days. And he shall be a blessing unto his posterity from generation to generation, and shall be satisfied in beholding his enemies cut off out of the earth.
					\hypertarget{EMS-1835-JS-22-SEP-1835-WP-b}%
				He
					\marginnote{b}%
				shall be filled with a fulness of the good things of the earth: with houses and with lands, with the fruit of the vine and with the fat of the olive, and he shall feed on the finest of the wheat; and because of his liberal soul the Lord will make him rich; even with treachers of gold, silver, precious stones, and with all precious metals. He shall be a wise lawyer in Israel; for he shall understand the law of the Lord perfectly; and the renowned among men shall acknowledge his superior wisdom pertaining to the laws of nations and of kingdoms. Behold, he shall have understanding in all sciences and languages, and with his brother <Oliver [Cowdery]> shall write and arrange many books for the good of the church, that the young may grow up in wisdom: these shall remain for a memorial unto their names, and the generations to come shall call them blessed. His days shall be prolonged upon the land of Zion; and when he is old and bowed down with many years, behold, he shall lift up his eyes and the heavens shall be opened upon him. He shall be lifted up at the last day, and his rest shall be glorious. Amen.
				
				\begin{tightright}
				Oliver Cowdery, \uline{clerk and Recorder}.
				\end{tightright}
				
				Given Sept. 22, 1835, and recorded Oct. 3, 1835.
				
		% -----
		\subsubsection{Blessing from Oliver Cowdery, 22 September 1835}
		% -----
			\label{EMS-1835-JS-22-SEP-1835-OC}
			% \label{EMS-1830-JS-DAY-3LETTERMONTH-YEAR}
			
			\begin{description}
				\item[Print Sources]
			\end{description}

				\begin{tabular}{p{0.5in}p{1in}p{0.5in}p{1in}}
				JSP	 	& D4:436--441		& JSR 		& --- \\
				DC	 	& --- 		& HC 		& --- \\
				HDDC 	& ---		& TCHC 		& --- \\
				\end{tabular}
				% HDDC = woodford's DC diss; JSR = Marquardt's JS Revelations; TCHC = Vogel HC
			
			\begin{description}
				\item[Online Access] \href{https://www.josephsmithpapers.org/paper-summary/blessing-from-oliver-cowdery-22-september-1835/1}{Here}
				% \item[Analysis] \hyperref[XX]{Here}
			\end{description}
			
			\begin{description}
				\item[Background]
			\end{description}
			
				% It might also be useful to give a two sentence context of the period: e.g., "This speech was given in Nauvoo to a small congregation one month prior to the destruction of the Nauvoo Expositor. These and other tensions may have been on the prophet's mind when he gave the speech."
			
			\begin{description}
				\item[Original Source]
			\end{description}
			
				% brief description of original source material, with reference to JSP or somewhere else for more info
				% Background material: it might be helpful to have a note that says whether a given document was presented as a speech, was written in a journal, etc.
				% Also a comment here about how reliable the source might be, or what shape it is in, if there are large omissions or parts that are hard to understand, and so on.
			
			\begin{description}
				\item[Text]
			\end{description}
			
					\hypertarget{EMS-1835-JS-22-SEP-1835-OC-a}%
				Blessed of
					\marginnote{a}%
				the Lord is my brother, for the integrity of his heart and the steadfastness of his soul. Upheld by the arm of the Almighty he shall never fall, but shall be strengthened by his right hand till he overcomes. Like Jacob of old he shall wrestle with the angel, and as a prince shall he have power with God, and shall prevail. Ever faithful to <his> friends and true to his word, the goodness of the Most High shall sustain him, and thousands shall stand up to defend him from the hand of his enemies, and put forth the hand and ward off the blow were it needful: but ere his foes are aware he shall be hid under the pavillion of the Lord Jehovah; for, with the voice of his thunder shall he strike terror to their hearts, and as with the wings of an eagle shall my brother be carried beyond all harm, by the power of the Annointed. From amid the burning bush, like Moses of old, shall he hear the voice, saying, I am the God of thy fathers, Abraham, Isaac and Jacob, I have seen, I have seen, the affliction of my people and am come down to deliver them: go, thou, and say to the strength of my house, To your tents, O Israel: build up the wastes and raise up the foundation of desolation that this generation has made.
					\hypertarget{EMS-1835-JS-22-SEP-1835-OC-b}%
				Thus
					\marginnote{b}%
				shall he be honored of the Lord, and thus shall it be recorded of him, that the generations to come may bless his name, in Israel, saying, The Lord make thee as Joseph the Seer, who was of the house of Ephraim, the brother of Manasseh: the Lord do thee good and bring peace and blessings opun [upon] thy house as he brought them upon the house of Joseph the Seer, who was raised up of a choice vine from the stem of Jacob through the root of Joseph, even that Joseph who was separated from his brethren--- For, like Joseph of old shall he be: he shall save the just from desolation, by the wise counsel of the Almighty; for by his direction shall they gather into store-houses and barns, till they overflow with the richness of the fruit of harvest: and by this means shall the just be saved from famine, while the nations of the wicked are distressed and faint. In due time shall he go forth toward the north, and by the power of his word shall the deep begin to give way and the ice melt before the sun. By the keys of the Kingdom shall he lead Israel into the land of Zion while the house of Jacob shouts in the danse and in the song--- Joy, O my soul, in that day, for thou shalt be with him and bear thy part in the keys which are confirmed <upon> thee for an everlasting priesthood, forever and ever--- Joy, O my heart, in that day, with thanks giving and with praise, for thou shalt stand with him before the hosts of Israel--- the lame shall leap as a hart, the old shall renew his strength, and the virgin of Israel, with the youth, shall exalt the name of our God upon harps and instruments of tens[e] strings.
					\hypertarget{EMS-1835-JS-22-SEP-1835-OC-c}%
				He
					\marginnote{c}%
				shall be a shure arrow in the bow of his God, for he shall be hid under the shadow of his wing. His loins shall be like iron, girded by the hand of the Lord, and his feet shall be swift to execute the commandment of the Most High when he shall say, Destroy. His name shall be had in everlasting remembrance, and his \sout{name} <seed> after him, for they shall be saved to the uttermost. His fame shall be sounded in foreign lands, even to the ends of the earth, as well as nigh at home: for in this the times shall change--- a prophet shall have honor in his own country. His learning and wisdom shall astonish the great, for they shall acknowledge that by his intelligence he has far surpassed their \sout{learng} <learning> and their science. In palaces of governors, rulers and kings shall he be honored, even in his person, for God shall give him power to prevail. He shall be a lawgiver to Israel and shall teach the house of Jacob the statutes of the Most High. His testimony shall shine like the sun, and the weight of his influence shall be like the great river that rises on the east of the lasting hills, and flows into the great deep--- so shall his righteousness ever abound. He shall partake of the blessings of Abraham, Isaac and Jacob: the chief things of the ancient mountains, the precious things that couch beneath, and of the treasures hid in the sand. The records of past ages and generations, and the histories of ancient days shall he bring forth: even the record of the Nephites shall he again obtain, with all those hid up by Mormon, and others who were righteous, and many others, till he is overwhelmed with knowledge. No precious thing shall slumber from his possession, for he shall be covered with the most choice of all ages, till his soul shall be satisfied and his heart shall say, Enough, Enough!
					\hypertarget{EMS-1835-JS-22-SEP-1835-OC-d}%
				In
					\marginnote{d}%
				his hands shall the Urim and Thummim remain and the holy ministry, and the keys of the evangelical priesthood, also, for an everlasting priesthood forever, even the patriarchal; for, behold, he is the first patriarch in the last days. He shall sit in the great assembly and general council of patriarchs, and execute the will and commandment of God under the direction of the Ancient of Days; for he shall have his place and act in his station. Behold, my brother Joseph is blessed: blessed are all who bless him, and blessed are all those whom he blessess. Multitudes, multitudes, shall come to a knowledge of the truth through his ministry, and he shall be welcomed into the presence of kings and the great ones of the earth; for he shall claim his place among the nobles of the earth and shall be reverenced by them. He shall also be filled withe abundance of the fat of the earth: his flocks shall bring forth thousands and tens of thousands: his fats [vats] shall overflow with wine and oil: his cattle shall increase to a multitude: he shall have horses and mules, asses, she asses and dromedaries, camels and elephants, and all swift beasts, and when he goes forth in haste his chariots shall roar like the approach of an army: he shall have gold and silver, precious stones, diamonds, pearls, and the pure platina [platinum], with the antiquities of every kind.
					\hypertarget{EMS-1835-JS-22-SEP-1835-OC-e}%
				Thus
					\marginnote{e}%
				shall God bless, and thus shall he be prospered: and he shall have peace after a little; for his enemies shall be consumed, many of them, and many shall turn and be his friends in very deed: he shall remain to a good old age, even till his head is like the pure wool. Behold, there is no end to the vision, of the multiplicity of blessings and glories which shall come upon my brother Joseph. He shall possess a mansion on high and have an inheritance in that city which is like pure gold, even like transparent glass. His rest shall be glorious and his name remain forever. Thus closes the vision, and thus it shall be; even so. Amen. Given in the evening of September 22, 1835, and recorded October 3, 1835. Oliver Cowdery.
				
		% -----
		\subsubsection{Minutes, 28--29 September 1835}
		% -----
			\label{EMS-1835-JS-28-29-SEP-1835}
			% \label{EMS-1830-JS-DAY-3LETTERMONTH-YEAR}
			
			\begin{description}
				\item[Print Sources]
			\end{description}

				\begin{tabular}{p{0.5in}p{1in}p{0.5in}p{1in}}
				JSP	 	& D4:446--455		& JSR 		& --- \\
				DC	 	& --- 				& HC 		& 2:284--286 \\
				HDDC 	& ---				& TCHC 		& 2:274--277 \\
				PJS 	& 1:101--102			& 	 		&  \\
				\end{tabular}
				% HDDC = woodford's DC diss; JSR = Marquardt's JS Revelations; TCHC = Vogel HC
			
			\begin{description}
				\item[Online Access] \href{https://www.josephsmithpapers.org/paper-summary/minutes-28-29-september-1835/1}{Here}
				% \item[Analysis] \hyperref[XX]{Here}
			\end{description}
			
			\begin{description}
				\item[Background]
			\end{description}
			
				% It might also be useful to give a two sentence context of the period: e.g., "This speech was given in Nauvoo to a small congregation one month prior to the destruction of the Nauvoo Expositor. These and other tensions may have been on the prophet's mind when he gave the speech."
			
			\begin{description}
				\item[Original Source]
			\end{description}
			
				% brief description of original source material, with reference to JSP or somewhere else for more info
				% Background material: it might be helpful to have a note that says whether a given document was presented as a speech, was written in a journal, etc.
				% Also a comment here about how reliable the source might be, or what shape it is in, if there are large omissions or parts that are hard to understand, and so on.
			
			\begin{description}
				\item[Text]
			\end{description}
			
				Charge preferred against Lorenzo Young, which is this that poor men ought not to raise up seed or children.
				
				Preferred by W[illiam] W. Phelps,
				
				President O. Cowdery testified that the above declaration fell from the lips of said Young at the home of President J. Smith \supers{Jun\supersu{r}} also that he believed it was right to have sexual intercourse notwithstanding. President J. Smith \supers{Ju}\supersu{n}\supers{\supers{r.}} corroborates the above testimony, also that he did not intend to have any more children. After hearing. The case was fairly laid open by the counsellors on both sides. The accuser made his remarks. The accused then arose and made an humble acknowledgement to the satisfaction of the Presidency \& counsellors who retain him in full fellowship as an Elder in the church of the Latter. Day. Saints.
				
		% -----
		\subsubsection{Letter to the Elders of the Church, 2 October 1835}
		% -----
			\label{EMS-1835-JS-2-OCT-1835}
			% \label{EMS-1830-JS-DAY-3LETTERMONTH-YEAR}
			
			\begin{description}
				\item[Print Sources]
			\end{description}

				\begin{tabular}{p{0.5in}p{1in}p{0.5in}p{1in}}
				JSP	 	& D5:6--15		& JSR 		& --- \\
				DC	 	& --- 			& HC 		& 2:253--259 \\
				HDDC 	& ---			& TCHC 		& 2:257--262 \\
				TPJS 	& 79--83			& 	 		&  \\
				PJS 	& 1:102			& 	 		&  \\
				\end{tabular}
				% HDDC = woodford's DC diss; JSR = Marquardt's JS Revelations; TCHC = Vogel HC
			
			\begin{description}
				\item[Online Access] \href{https://www.josephsmithpapers.org/paper-summary/letter-to-the-elders-of-the-church-2-october-1835/1}{Here}
				% \item[Analysis] \hyperref[XX]{Here}
			\end{description}
			
			\begin{description}
				\item[Background]
			\end{description}
			
				% It might also be useful to give a two sentence context of the period: e.g., "This speech was given in Nauvoo to a small congregation one month prior to the destruction of the Nauvoo Expositor. These and other tensions may have been on the prophet's mind when he gave the speech."
			
			\begin{description}
				\item[Original Source]
			\end{description}
			
				% brief description of original source material, with reference to JSP or somewhere else for more info
				% Background material: it might be helpful to have a note that says whether a given document was presented as a speech, was written in a journal, etc.
				% Also a comment here about how reliable the source might be, or what shape it is in, if there are large omissions or parts that are hard to understand, and so on.
				
				Note the spaces in between lines and block quotations here.
			
			\begin{description}
				\item[Text]
			\end{description}
			
					\hypertarget{EMS-1835-JS-2-OCT-1835-a}%
				\phantom{ }\textit{To the elders of the church of Latter Day Saints.}
					\marginnote{a}

				After so long a time, and after so many things having been said, I feel it my duty to drop a few hints, that, perhaps, the elders, traveling through the world to warn the inhabitants of the earth to flee the wrath to come, and save themselves from this untoward generation, may be aided in a measure, in doctrine, and in the way of their duty. I have been laboring in this cause for eight years, during which time I have traveled much, and have had much experience. I removed from Seneca county, N. Y. to Geauga county, Ohio, in February, 1831.

				Having received, by an heavenly vision, a commandment, in June following, to take my journey to the western boundaries of the State of Missouri, and there designate the very spot, which was to be the central spot, for the commencement of the gathering together of those who embrace the fulness of the everlasting gospel---I accordingly undertook the journey with certain ones of my brethren, and, after a long and tedious journey, suffering many privations and hardships, I arrived in Jackson county Missouri; and, after viewing the country, seeking diligently at the hand of God, he manifested himself unto me, and designated to me and others, the very spot upon which he designed to commence the work of the gathering, and the upbuilding of an holy city, which should be called Zion:---Zion because it is to be a place of righteousness, and all who build thereon, are to worship the true and living God---and all believe in one doctrine even the doctrine of our Lord and Savior Jesus Christ.

					\hypertarget{EMS-1835-JS-2-OCT-1835-b}%
				``Thy
					\marginnote{b}%
				watchmen shall lift up the voice; with the voice together shall they sing: for they shall see eye to eye, when the Lord shall bring again Zion.''---Isaiah 52:8.

				Here we pause for a moment, to make a few remarks upon the idea of gathering to this place. It is well known that there were lands belonging to the government, to be sold to individuals; and it was understood by all, at least we believed so, that we lived in a free country, a land of liberty and of laws, guaranteeing to every man, or any company of men, the right of purchasing lands, and settling, and living upon them: therefore we thought no harm in advising the Latter Day Saints, or Mormons, as they are reproachfully called, to gather to this place, inasmuch as it was their duty, (and it was well understood so to be,) to purchase, \textit{with money}, lands, and live upon them---not infringing upon the civil rights of any individual, or community of people: always keeping in view the saying, ``Do unto others as you would wish to have others do unto you.'' Following also the good injunction: ``Deal justly, love mercy, and walk humbly with thy God.''

					\hypertarget{EMS-1835-JS-2-OCT-1835-c}%
				These
					\marginnote{c}%
				were our motives in teaching the people, or Latter Day Saints, to gather together, beginning at this place. And inasmuch as there are those who have had different views from this, we feel, that it is a cause of deep regret: For, be it known unto all men, that our principles concerning this thing, have not been such as have been represented by those who, we have every reason to believe, are designing and wicked men, that have said that this was our doctrine:---to infringe upon the rights of a people who inhabit our civil and free country: such as to drive the inhabitants of Jackson county from their lands, and take possession thereof unlawfully. Far, yea, far be such a principle from our hearts: it never entered into our mind, and we only say, that God shall reward such in that day when he shall come to make up his jewels.

				But to return to my subject: after having ascertained the very spot, and having the happiness of seeing quite a number of the families of my brethren, comfortably situated upon the land, I took leave of them, and journeyed back to Ohio, and used every influence and argument, that lay in my power, to get those who believe in the everlasting covenant, whose circumstances would admit, and whose families were willing to remove to the place which I now designated to be the land of Zion: And thus the sound of the gathering, and of the doctrine, went abroad into the world; and many we have reason to fear, having a zeal not according to knowledge, not understanding the pure principles of the doctrine of the church, have no doubt, in the heat of enthusiasm, taught and said many things which are derogatory to the genuine character and principles of the church, and for these things we are heartily sorry, and would apologize if an apology would do any good.

					\hypertarget{EMS-1835-JS-2-OCT-1835-d}%
				But we
					\marginnote{d}%
				pause here and offer a remark upon the saying which we learn has gone abroad, and has been handled in a manner detrimental to the cause of truth, by saying, ``that in preaching the doctrine of gathering, we break up families, and give license for men to leave their families; women their husbands; children their parents, and slaves their masters, thereby deranging the order, and breaking up the harmony and peace of society.'' We shall here show our faith, and thereby, as we humbly trust, put an end to these faults, and wicked misrepresentations, which have caused, we have every reason to believe, thousands to think they were doing God's service, when they were persecuting the children of God: whereas, if they could have enjoyed the true light, and had a just understanding of our principles, they would have embraced them with all their hearts, and been rejoicing in the love of the truth.

				And now to show our doctrine on this subject, we shall commence with the first principles of the gospel, which are repentance, and baptism for the remission of sins, and the gift of the Holy Ghost by the laying on of the hands. This we believe to be our duty, to teach to all mankind the doctrine of repentance, which we shall endeavor to show from the following quotations:

					\hypertarget{EMS-1835-JS-2-OCT-1835-e}%
				``Then
					\marginnote{e}%
				opened he their understanding, that they might understand the scriptures, and said unto them, thus it is written, and thus it behoved Christ to suffer, and to rise from the dead, the third day; and that repentance and remission of sins should be preached in his name among all nations, beginning at Jerusalem.''---Luke 24:45, 46, 47.

				By this we learn, that it behoved Christ to suffer, and to be crucified, and rise again on the third day, for the express purpose that repentance and remission of sins should be preached unto all nations.

				``Then Peter said unto them, repent, and be baptized every one of you, in the name of Jesus Christ, for the remission of sins, and ye shall receive the gift of the Holy Ghost. For the promise is unto you, and to your children, and to all that are afar off, even as many as the Lord our God shall call.''---Acts 2:38, 39.

					\hypertarget{EMS-1835-JS-2-OCT-1835-f}%
				By this
					\marginnote{f}%
				we learn, that the promise of the Holy Ghost, is unto as many as the doctrine of repentance was to be preached, which was unto all nations. And we discover also, that the promise was to extend by lineage: for Peter says, ``not only unto you, but unto your children, and unto all that are afar off.'' From this we infer that it was to continue unto their children's children, and even unto as many generations as should come after, even as many as the Lord their God should call.--- We discover here that we are blending two principles together, in these quotations. The first is the principle of repentance, and the second is the principle of remission of sins. And we learn from Peter, that remission of sins is obtained by baptism in the name of the Lord Jesus Christ; and the gift of the Holy Ghost follows inevitably: for, says Peter, ``you shall receive the gift of the Holy Ghost.'' Therefore we believe in preaching the doctrine of repentance in all the world, both to old and young, rich and poor, bond and free, as we shall endeavor to show hereafter---how and in what manner, and how far it is binding upon the consciences of mankind, making proper distinctions between old and young men, women and children, and servants.

				But we discover, in order to be benefitted by the doctrine of repentance, we must believe in obtaining the remission of sins. And in order to obtain the remission of sins, we must believe in the doctrine of baptism, in the name of the Lord Jesus Christ. And if we believe in baptism for the remission of sins, we may expect a fulfilment of the promise of the Holy Ghost: for the promise extends to all whom the Lord our God shall call. And hath he not surely said, as you will find in the last chapter of Revelations:

					\hypertarget{EMS-1835-JS-2-OCT-1835-g}%
				``And
					\marginnote{g}%
				the Spirit and the bride say, Come. And let him that heareth, say, Come. And let him that is athirst, come. And whosoever will, let him take the water of life freely.'' Rev. 22:17.

				Again the Savior says:

				``Come unto me, all ye that labor, and are heavy laden, and I will give you rest. Take my yoke upon you, and learn of me; for I am meek and lowly in heart; and ye shall find rest unto your souls. For my yoke is easy, and my burden is light.''---Math. 11:28, 29, 30.

				Again Isaiah says:

				``Look unto me, and be ye saved, all the ends of the earth: for I am God, and there is none else. I have sworn by myself, the word is gone out of my mouth in righteousness, and shall not return, that unto me every knee shall bow, every tongue shall swear. Surely, shall one say, in the Lord have I righteousness and strength: even to him shall men come; and all that are incensed against him shall be ashamed.''---Isaiah 45:22, 23, 24.

				And to show further connections in proof of the doctrine above named, we quote the following scriptures:

				``Him hath God exalted with his right hand, to be a Prince and a Savior, for to give repentance to Israel, and forgiveness of sins. And we are his witnesses of these things; and so is also the Holy Ghost, whom God hath given to them that obey him.''---Acts 5:31, 32.

					\hypertarget{EMS-1835-JS-2-OCT-1835-h}%
				``But
					\marginnote{h}%
				when they believed Philip, preaching the things concerning the kingdom of God, and the name of Jesus Christ, they were baptized, both men and women. Then Simon himself believed also; and when he was baptized, he continued with Philip, and wondered, beholding the miracles and signs which were done. Now when the apostles, which were at Jerusalem, heard that Samaria had received the word of God, they sent unto them Peter and John; who, when they were come down, prayed for them, that they might receive the Holy Ghost. (For as yet he was fallen upon none of them: only they were baptized in the name of the Lord Jesus.)--- Then laid they their hands on them, and they received the Holy Ghost. . . . And as they went on their way, they came unto a certain water; and the eunuch said, See, here is water; what doth hinder me to be baptized?---And Philip said, If thou believest with all thine heart thou mayest. And he answered and said, I believe that Jesus Christ is the Son of God. And he commanded the chariot to stand still: and they went down both into the water, both Philip and the eunuch; and he baptized him. And, when they were come up out of the water, the Spirit of the Lord caught away Philip, that the eunuch saw him no more: and he went on his way rejoicing. But Philip was found at Azotus: and, passing through, he preached in all the cities, till he came to Cesarea.''---Acts 8:12, 13, 14, 15, 16, 17,------36, to the end.

					\hypertarget{EMS-1835-JS-2-OCT-1835-i}%
				``While
					\marginnote{i}%
				Peter yet spake these words, the Holy Ghost fell on all them which heard the word. And they of the circumcision, which believed, were astonished, as many as came with Peter, because that on the Gentiles also was poured out the gift of the Holy Ghost: for they heard them speak with tongues, and magnify God. Then answered Peter, Can any man forbid water, that these should not be baptized, which have received the Holy Ghost as well as we? And he commanded them to be baptized in the name of the Lord. Then prayed they him to tarry certain days.''---Acts 10:44, 45, 46, 47, 48.

				``And on the Sabbath, we went out of the city, by a river side, where prayer was wont to be made; and we sat down, and spake unto the women that resorted thither. And a certain woman, named Lydia, a seller of purple, of the city of Thyatira, which worshipped God, heard us: whose heart the Lord opened, that she attended unto the things which were spoken of Paul. And when she was baptized, and her household, she besought us, saying, If ye have judged me to be faithful to the Lord, come into my house, and abide there. And she constrained us. . . . . And at midnight Paul and Silas prayed, and sang praises unto God: and the prisoners heard them. And suddenly there was a great earthquake, so that the foundations of the prison were shaken; and immediately all the doors were opened, and every one's bands were loosed. And the keeper of the prison awaking out of his sleep, and seeing the prison doors open, he drew out his sword, and would have killed himself, supposing that the prisoners had been fled. But Paul cried with a loud voice, saying, Do thyself no harm; for we are all here. Then he called for a light, and sprang in, and came trembling, and fell down before Paul and Silas; and brought them out, and said, Sirs, what must I do to be saved? And they said believe on the Lord Jesus Christ, and thou shalt be saved and thy house. And they spake unto him the word of the Lord, and to all that were in his house. And he took them the same hour of the night, and washed their stripes, and was baptized, he and all his, staightway. And when he had brought them into his house, he set meat before them, and rejoiced, believing in God with all his house.''---Acts 16:13, 14, 15.------25, to 35.

					\hypertarget{EMS-1835-JS-2-OCT-1835-j}%
				``And it
					\marginnote{j}%
				came to pass, that, while Apollos was at Corinth, Paul, having passed through the upper coasts, came to Ephesus; and finding certain disciples, he said unto them, Have ye received the Holy Ghost since ye believed? And they said unto him, We have not so much as heard whether there be any Holy Ghost. And he said unto them, Unto what then were ye baptized? And they said, Unto John's baptism. Then said Paul, John verily baptized with the baptism of repentance, saying unto the people, that they should believe on him which should come after him, that is, on Christ Jesus. When they heard this, they were baptized in the name of the Lord Jesus. And, when Paul had laid his hands upon them, the Holy Ghost came on them; and they spake with tongues, and prophesied.''---Acts 19:1, 2, 3, 4, 5, 6.

				[``]And one Ananias, a devout man, according to the law, having a good report of all the Jews which dwelt there, Came unto me, and stood, and said unto me, Brother Saul, receive thy sight. And the same hour I looked up upon him. And he said, the God of our fathers hath chosen thee, that thou shouldst know his will, and see that Just One, and shouldst hear the voice of his mouth. For thou shalt be his witness unto all men, of what thou hast seen and heard. And now why tarriest thou? arise, and be baptized, and wash away thy sins, calling on the name of the Lord.''---Acts 22:12, 13, 14, 15, 16.

					\hypertarget{EMS-1835-JS-2-OCT-1835-k}%
				``For,
					\marginnote{k}%
				when for the time ye ought to be teachers, ye have need that one teach you again which be the first principles of the oracles of God; and are become such as have need of milk, and not of strong meat. For every one that useth milk, is unskilful in the word of righteousness; for he is a babe. But strong meat belongeth to them that are of full age, even those who by reason of use, have their senses exercised to discern both good and evil.''---Heb. 5:12, 13, 14.

				``Therefore, leaving the principles of the doctrine of Christ, let us go on unto perfection; not laying again the foundation of repentance from dead works, and of faith towards God, of the doctrine of baptisms, and of laying on of hands, and of resurrection of the dead, and of eternal judgment. And this will we do, if God permit. For it is impossible for those who were once enlightened, and have tasted of the heavenly gift, and were made partakers of the Holy Ghost, and have tasted the good word of God, and the powers of the world to come, if they shall fall away, to renew them again unto repentance; seeing they crucify to themselves the Son of God afresh, and put him to an open shame.['']---Heb. 6:1, 2, 3, 4, 5, 6.

				These quotations are so plain, in proving the doctrine of repentance and baptism for the remission of sins, I deem it unnecessary to enlarge this letter with comments upon them---but I shall continue the subject in my next.

				In the bonds of the new and everlasting covenant,
				
				\begin{tightright}
				JOSEPH SMITH, jr.
				\end{tightright}

				\textsc{John Whitmer}, Esq.
				
		% -----
		\subsubsection{Record, 5 October 1835}
		% -----
			\label{EMS-1835-JS-5-OCT-1835}
			% \label{EMS-1830-JS-DAY-3LETTERMONTH-YEAR}
			
			\begin{description}
				\item[Print Sources]
			\end{description}

				\begin{tabular}{p{0.5in}p{1in}p{0.5in}p{1in}}
				JSP	 	& J1:68		& JSR 		& --- \\
				DC	 	& --- 				& HC 		& --- \\
				HDDC 	& ---				& TCHC 		& --- \\
				PJS 	& 1:104			& 	 		&  \\
				\end{tabular}
				% HDDC = woodford's DC diss; JSR = Marquardt's JS Revelations; TCHC = Vogel HC
			
			\begin{description}
				\item[Online Access] \href{https://www.josephsmithpapers.org/paper-summary/journal-1835-1836/5}{Here}
				% \item[Analysis] \hyperref[XX]{Here}
			\end{description}
			
			\begin{description}
				\item[Background]
			\end{description}
			
				% It might also be useful to give a two sentence context of the period: e.g., "This speech was given in Nauvoo to a small congregation one month prior to the destruction of the Nauvoo Expositor. These and other tensions may have been on the prophet's mind when he gave the speech."
			
			\begin{description}
				\item[Original Source]
			\end{description}
			
				% brief description of original source material, with reference to JSP or somewhere else for more info
				% Background material: it might be helpful to have a note that says whether a given document was presented as a speech, was written in a journal, etc.
				% Also a comment here about how reliable the source might be, or what shape it is in, if there are large omissions or parts that are hard to understand, and so on.
			
			\begin{description}
				\item[Text]
			\end{description}
			
				Monday 5th returned home being much fatiegued riding in the rain spent the remainder of the day in reading and meditation \&c and in the evening attend[ed] a high councel of the twelve apostles, had a glorious time and gave them many instruction concerning their duties for time to come, told them that it was the will of God they should take their families to Missouri next season, also attend this fall the solemn assembly of the first Elders for the organization of the school of the prophets, and attend to the ordinence of the washing of feet and to prepare their hearts in all humility for an endowment \sout{from} with power from on high to which they all agreed with one accord, and seamed to be greatly rejoiced may God spare the lives of the twelve with one accord to a good old age for christ the redeemers sake amen
				
		% -----
		\subsubsection{Blessing to Newel K. Whitney, 7 October 1835}
		% -----
			\label{EMS-1835-JS-7-OCT-1835}
			% \label{EMS-1830-JS-DAY-3LETTERMONTH-YEAR}
			
			\begin{description}
				\item[Print Sources]
			\end{description}

				\begin{tabular}{p{0.5in}p{1in}p{0.5in}p{1in}}
				JSP	 	& D5:18--21; J1:69--71		& JSR 		& --- \\
				DC	 	& --- 			& HC 		& 2:288--289 \\
				HDDC 	& ---			& TCHC 		& 2:283--284 \\
				PJS 	& 1:106--107			& 	 		&  \\
				\end{tabular}
				% HDDC = woodford's DC diss; JSR = Marquardt's JS Revelations; TCHC = Vogel HC
			
			\begin{description}
				\item[Online Access] \href{https://www.josephsmithpapers.org/paper-summary/blessing-to-newel-k-whitney-7-october-1835/1}{Here}
				% \item[Analysis] \hyperref[XX]{Here}
			\end{description}
			
			\begin{description}
				\item[Background]
			\end{description}
			
				% It might also be useful to give a two sentence context of the period: e.g., "This speech was given in Nauvoo to a small congregation one month prior to the destruction of the Nauvoo Expositor. These and other tensions may have been on the prophet's mind when he gave the speech."
			
			\begin{description}
				\item[Original Source]
			\end{description}
			
				% brief description of original source material, with reference to JSP or somewhere else for more info
				% Background material: it might be helpful to have a note that says whether a given document was presented as a speech, was written in a journal, etc.
				% Also a comment here about how reliable the source might be, or what shape it is in, if there are large omissions or parts that are hard to understand, and so on.
				
				There is a period at the end of what I have here, and not at the end of the J1 account; check them all.
			
			\begin{description}
				\item[Text]
			\end{description}
			
					\hypertarget{EMS-1835-JS-7-OCT-1835-a}%
				Blessed of
					\marginnote{a}%
				the lord is bro [Newel K.] Whitney even the bishop of the church of the latter day saints, for the bishoprick shall never be taken away from him while he liveth and the time cometh that he shall overcome all the narrow mindedness of his heart and all his covetous desires that so easily besetteth him and <he> shall \sout{deliver} deal with a liberal hand to the poor and the needy the sick and the afflicted the widow and the fatherless and marviously [marvelously] and miraculously shall the Lord his God provid for him. even that he shall be blessed with \sout{a} <all the \sout{the}> fullness of the good things of this earth and his seed after him from generation to generation and it shall come to pass that according to to the measure that he meeteth out with a liberal hand unto the poor so shall it be measured to him again by the hand of his God even an hundred fold Angels shall guard <his> house and shall guard the lives of his posterity, and they shall become very great and very numerous on the earth,
					\hypertarget{EMS-1835-JS-7-OCT-1835-b}%
				whomsoever
					\marginnote{b}%
				he blesseth they shall be blessed. whomsoever he curseth they shall be cursed. and when his enemies seek him unto his hurt and distruction let him rise up and curse and the hand of God shall be upon his enemies in Judgment they shall be utterly confounded and brought to dessolation, therefor he shall be preserved unto the utmost and his <\textbf{life}> \sout{day} shall be precious in the sight of the Lord. he shall rise up and shake himself as a lion riseth out of his nest and roareth untill he shaketh the hills and as a lion goeth forth among the Lesser beasts, so shall the goings forth of him <\textbf{be}> whom the Lord hath anointed to exalt the poor and to humble the rich, therefor his name shall be on high and his rest among the sanctified.
				
		% -----
		\subsubsection{Revelation, 18 October 1835}
		% -----
			\label{EMS-1835-JS-18-OCT-1835}
			% \label{EMS-1830-JS-DAY-3LETTERMONTH-YEAR}
			
			\begin{description}
				\item[Print Sources]
			\end{description}

				\begin{tabular}{p{0.5in}p{1in}p{0.5in}p{1in}}
				JSP	 	& D5:21--23		& JSR 		& --- \\
				DC	 	& --- 			& HC 		& --- \\
				HDDC 	& ---			& TCHC 		& --- \\
				\end{tabular}
				% HDDC = woodford's DC diss; JSR = Marquardt's JS Revelations; TCHC = Vogel HC
			
			\begin{description}
				\item[Online Access] \href{https://www.josephsmithpapers.org/paper-summary/revelation-18-october-1835/1}{Here}
				% \item[Analysis] \hyperref[XX]{Here}
			\end{description}
			
			\begin{description}
				\item[Background]
			\end{description}
			
				% It might also be useful to give a two sentence context of the period: e.g., "This speech was given in Nauvoo to a small congregation one month prior to the destruction of the Nauvoo Expositor. These and other tensions may have been on the prophet's mind when he gave the speech."
			
			\begin{description}
				\item[Original Source]
			\end{description}
			
				% brief description of original source material, with reference to JSP or somewhere else for more info
				% Background material: it might be helpful to have a note that says whether a given document was presented as a speech, was written in a journal, etc.
				% Also a comment here about how reliable the source might be, or what shape it is in, if there are large omissions or parts that are hard to understand, and so on.
				
				Note the spaces in between lines and block quotations here.
			
			\begin{description}
				\item[Text]
			\end{description}
			
				Oct 18, 1835. Sabbath
				
				This day assembled in the house of the Lord as usual and the Spirit of the Lord decended upon J. Smith Jr---the seer and he propheced: saying the L[or]d has showd to me this day by the Spirit of Revelation that the distress, and sickness that has heretofore prevailed among the children of Zion will be mitigated from this time forth.
				
		% -----
		\subsubsection{Prayer, 23 October 1835}
		% -----
			\label{EMS-1835-JS-23-OCT-1835}
			% \label{EMS-1830-JS-DAY-3LETTERMONTH-YEAR}
			
			\begin{description}
				\item[Print Sources]
			\end{description}

				\begin{tabular}{p{0.5in}p{1in}p{0.5in}p{1in}}
				JSP	 	& D5:23--25; J1:111--112	& JSR 		& --- \\
				DC	 	& --- 			& HC 		& 2:291 \\
				HDDC 	& ---			& TCHC 		& 2:285--286 \\
				PJS 	& 1:147--148			& 	 		&  \\
				\end{tabular}
				% HDDC = woodford's DC diss; JSR = Marquardt's JS Revelations; TCHC = Vogel HC
			
			\begin{description}
				\item[Online Access] \href{https://www.josephsmithpapers.org/paper-summary/prayer-23-october-1835/1}{Here}
				% \item[Analysis] \hyperref[XX]{Here}
			\end{description}
			
			\begin{description}
				\item[Background]
			\end{description}
			
				% It might also be useful to give a two sentence context of the period: e.g., "This speech was given in Nauvoo to a small congregation one month prior to the destruction of the Nauvoo Expositor. These and other tensions may have been on the prophet's mind when he gave the speech."
			
			\begin{description}
				\item[Original Source]
			\end{description}
			
				% brief description of original source material, with reference to JSP or somewhere else for more info
				% Background material: it might be helpful to have a note that says whether a given document was presented as a speech, was written in a journal, etc.
				% Also a comment here about how reliable the source might be, or what shape it is in, if there are large omissions or parts that are hard to understand, and so on.
				
				Note the spaces in between lines and block quotations here.
			
			\begin{description}
				\item[Text]
			\end{description}
			
				Copy of a prayer offered up. on the 23\supers{d} day of Oct 1835, by the following individuals, at 4 oclock P.M. viz. Joseph Smith jn, Oliver Cowdery; David Whitmer, Hirum [Hyrum] Smith John Whitmer, Sidn[e]y Rigdon, Samuel H. Smith, Frederick G. Williams, and W\supersu{m}\supers{.} W. Phelps, assembled and united in prayer, with one voice before the Lord, for the following blessings:
				
				That the Lord will give us means sufficient to deliver us from all our afflictions and difficulties, wherein we are placed by means of our debts; that he will open the way and deliver Zion in the appointed time and that without the shedding of blood; that he will hold our lives precious, and grant that we may live to the common age of man, and never fall into the hands nor power of the mob in Missourie nor in any other place; that he will also preserve our posterity, that none of them fall even to the end of time; that he will give us the blessings of the earth sufficient to carry us to Zion, and that we may purchase inheritances in that land, even enough to carry on <and accomplish> the work unto which he has appointed us; and also that he will assist all others who desire, according\sout{ly} to his commandments, to go up and purchase inheritances; and all this easily and without perplexity, and trouble; and finally, that in the end he will save us in his Celestial Kingdom. Amen.
				
				\begin{tightright}
				Oliver Cowdery \uline{Clerk}
				\end{tightright}
				
		% -----
		\subsubsection{Revelation, 27 October 1835}
		% -----
			\label{EMS-1835-JS-27-OCT-1835}
			% \label{EMS-1830-JS-DAY-3LETTERMONTH-YEAR}
			
			\begin{description}
				\item[Print Sources]
			\end{description}

				\begin{tabular}{p{0.5in}p{1in}p{0.5in}p{1in}}
				JSP	 	& D5:25--26; J1:75--76	& JSR 		& 273 \\
				DC	 	& --- 			& HC 		& 2:292--293 \\
				HDDC 	& ---			& TCHC 		& 2:286--287 \\
				\end{tabular}
				% HDDC = woodford's DC diss; JSR = Marquardt's JS Revelations; TCHC = Vogel HC
			
			\begin{description}
				\item[Online Access] \href{https://www.josephsmithpapers.org/paper-summary/revelation-27-october-1835/1}{Here}
				% \item[Analysis] \hyperref[XX]{Here}
			\end{description}
			
			\begin{description}
				\item[Background]
			\end{description}
			
				% It might also be useful to give a two sentence context of the period: e.g., "This speech was given in Nauvoo to a small congregation one month prior to the destruction of the Nauvoo Expositor. These and other tensions may have been on the prophet's mind when he gave the speech."
			
			\begin{description}
				\item[Original Source]
			\end{description}
			
				% brief description of original source material, with reference to JSP or somewhere else for more info
				% Background material: it might be helpful to have a note that says whether a given document was presented as a speech, was written in a journal, etc.
				% Also a comment here about how reliable the source might be, or what shape it is in, if there are large omissions or parts that are hard to understand, and so on.
				
				Footnote to additional info in the J1 account.
			
			\begin{description}
				\item[Text]
			\end{description}
			
				the word of the Lord came unto me saying my Servant Fredrick [Frederick G. Williams] shall come and shall have wisdom given him to deal prudently and my handmaden shall be delivered of a living child \& be spared
				
		% -----
		\subsubsection{Revelation, 1 November 1835}
		% -----
			\label{EMS-1835-JS-1-NOV-1835}
			% \label{EMS-1830-JS-DAY-3LETTERMONTH-YEAR}
			
			\begin{description}
				\item[Print Sources]
			\end{description}

				\begin{tabular}{p{0.5in}p{1in}p{0.5in}p{1in}}
				JSP	 	& D5:29--30; J1:81		& JSR 	& 273 \\
				DC	 	& --- 			& HC 	& 2:299 \\
				HDDC 	& ---			& TCHC 	& 2:291 \\
				PJS 	& 1:118--119			& 	 		&  \\
				\end{tabular}
				% HDDC = woodford's DC diss; JSR = Marquardt's JS Revelations; TCHC = Vogel HC
			
			\begin{description}
				\item[Online Access] \href{https://www.josephsmithpapers.org/paper-summary/revelation-1-november-1835/1}{Here}
				% \item[Analysis] \hyperref[XX]{Here}
			\end{description}
			
			\begin{description}
				\item[Background]
			\end{description}
			
				% It might also be useful to give a two sentence context of the period: e.g., "This speech was given in Nauvoo to a small congregation one month prior to the destruction of the Nauvoo Expositor. These and other tensions may have been on the prophet's mind when he gave the speech."
			
			\begin{description}
				\item[Original Source]
			\end{description}
			
				% brief description of original source material, with reference to JSP or somewhere else for more info
				% Background material: it might be helpful to have a note that says whether a given document was presented as a speech, was written in a journal, etc.
				% Also a comment here about how reliable the source might be, or what shape it is in, if there are large omissions or parts that are hard to understand, and so on.
				
				Note the spaces in between lines and block quotations here.
			
			\begin{description}
				\item[Text]
			\end{description}
			
				Verily thus Saith the Lord unto me, his servant Joseph Smith jun mine anger is kindle[d] against my servant Reynolds Cahoon because of his iniquities his covetous and dishonest principles in himself and family and he doth not purge them away and set his house in order, therefore if he repent not chastisment awaiteth him even as it seemeth good in my sight therefore go and declare unto him \sout{this} <these> word<s>
				
		% -----
		\subsubsection{Revelation, 2 November 1835}
		% -----
			\label{EMS-1835-JS-2-NOV-1835}
			% \label{EMS-1830-JS-DAY-3LETTERMONTH-YEAR}
			
			\begin{description}
				\item[Print Sources]
			\end{description}

				\begin{tabular}{p{0.5in}p{1in}p{0.5in}p{1in}}
				JSP	 	& D5:30--32; J1:81		& JSR 		& 274 \\
				DC	 	& --- 			& HC 		& 2:300 \\
				HDDC 	& ---			& TCHC 		& 2:292 \\
				PJS 	& 1:120			& 	 		&  \\
				\end{tabular}
				% HDDC = woodford's DC diss; JSR = Marquardt's JS Revelations; TCHC = Vogel HC
			
			\begin{description}
				\item[Online Access] \href{https://www.josephsmithpapers.org/paper-summary/revelation-2-november-1835/1}{Here}
				% \item[Analysis] \hyperref[XX]{Here}
			\end{description}
			
			\begin{description}
				\item[Background]
			\end{description}
			
				% It might also be useful to give a two sentence context of the period: e.g., "This speech was given in Nauvoo to a small congregation one month prior to the destruction of the Nauvoo Expositor. These and other tensions may have been on the prophet's mind when he gave the speech."
			
			\begin{description}
				\item[Original Source]
			\end{description}
			
				% brief description of original source material, with reference to JSP or somewhere else for more info
				% Background material: it might be helpful to have a note that says whether a given document was presented as a speech, was written in a journal, etc.
				% Also a comment here about how reliable the source might be, or what shape it is in, if there are large omissions or parts that are hard to understand, and so on.
				
				Note the spaces in between lines and block quotations here.
			
			\begin{description}
				\item[Text]
			\end{description}
			
				thus came the word of the Lord unto me saying it is not my will that my servant Frederick [G. Williams] should go to New York, but inasmuch as he wishes to go and visit his relatives that he may warn them to flee the wrath to come let him go and see them, for that purpose and let that be his only business, and behold in this thing he shall be blessed with power \sout{while} to overcome their prejudices, Verily thus saith the Lord Amen.
				
		% -----
		\subsubsection{Revelation, 3 November 1835}
		% -----
			\label{EMS-1835-JS-3-NOV-1835}
			% \label{EMS-1830-JS-DAY-3LETTERMONTH-YEAR}
			
			\begin{description}
				\item[Print Sources]
			\end{description}

				\begin{tabular}{p{0.5in}p{1in}p{0.5in}p{1in}}
				JSP	 	& D5:32--36; J1:83		& JSR 		& 274--275 \\
				DC	 	& --- 			& HC 		& 2:300--301 \\
				HDDC 	& ---			& TCHC 		& 2:292--293 \\
				PJS 	& 1:120--121			& 	 		&  \\
				\end{tabular}
				% HDDC = woodford's DC diss; JSR = Marquardt's JS Revelations; TCHC = Vogel HC
			
			\begin{description}
				\item[Online Access] \href{https://www.josephsmithpapers.org/paper-summary/revelation-3-november-1835/1}{Here}
				% \item[Analysis] \hyperref[XX]{Here}
			\end{description}
			
			\begin{description}
				\item[Background]
			\end{description}
			
				% It might also be useful to give a two sentence context of the period: e.g., "This speech was given in Nauvoo to a small congregation one month prior to the destruction of the Nauvoo Expositor. These and other tensions may have been on the prophet's mind when he gave the speech."
			
			\begin{description}
				\item[Original Source]
			\end{description}
			
				% brief description of original source material, with reference to JSP or somewhere else for more info
				% Background material: it might be helpful to have a note that says whether a given document was presented as a speech, was written in a journal, etc.
				% Also a comment here about how reliable the source might be, or what shape it is in, if there are large omissions or parts that are hard to understand, and so on.
			
			\begin{description}
				\item[Text]
			\end{description}
			
					\hypertarget{EMS-1835-JS-3-NOV-1835-a}%
				Thus
					\marginnote{a}%
				came the word of the Lord unto me \sout{saying} concerning the, Twelve <saying>
				
				behold they are under condemnation, because they have not been sufficiently humble in my sight, and in consequence of their covetous desires, in that they have not dealt equally with each other in the division of the moneys which came into their hands, nevertheless some of them, dealt equally therefore they shall be rewarded, but Verily I say unto you they must all humble themselves before Me, before they will be accounted worthy to receive an endowment to go forth in my name unto all nations, as for my Servant William [Smith] let the Eleven humble themselves in prayer and in faith and wait on me in patience and my servant William shall return, and I will yet make him a polished shaft in my quiver, in bringing down the wickedness and abominations of men and their shall be none mightier than he in his day and generation, nevertheless if he repent not spedily he shall be brought low and shall be chastened sorely for all his iniquities he has commited against me, nevertheless the sin which he hath sin[n]ed against me is not even now more grevious than the sin with which my servant David W. Patten and my servant Orson Hyde and my servant W\supers{m} E. McLellen [McLellin] have sinded [sinned] against me, and the residue are not sufficiently humble before me,
					\hypertarget{EMS-1835-JS-3-NOV-1835-b}%
				behold
					\marginnote{b}%
				the parable which I spake concerning a man having twelve Sons, for what man amon[g] you having twelve Sons and is no respecter to them and they serve him obediantly and he saith unto the one be thou clothed in robes and sit thou here, and to the other be thou clothed in rages [rags] and sit thou there, and looketh upon his sons and saith I am just, ye will answer and say no man, and ye answer truly, therefore Verily thus saith the Lord your God I appointed these twelve that they should be equal in their ministry and in their portion and in their evangelical rights, wherefore they have sined a verry grevious sin, in asmuch as they have made themselves unequal and have not hearkned unto my voice therfor let them repent speedily and prepare their hearts for the solem assembly and for the great day which is to come Verely thus saith the Lord Amen.
				
		% -----
		\subsubsection{Record, 3 November 1835}
		% -----
			\label{EMS-1835-JS-3-NOV-1835-REC}
			% \label{EMS-1830-JS-DAY-3LETTERMONTH-YEAR}
			
			\begin{description}
				\item[Print Sources]
			\end{description}

				\begin{tabular}{p{0.5in}p{1in}p{0.5in}p{1in}}
				JSP	 	& J1:84			& JSR 		& ---  \\
				DC	 	& --- 			& HC 		& ---  \\
				HDDC 	& ---			& TCHC 		& ---  \\
				PJS 	& 1:121			& 	 		&  \\
				\end{tabular}
				% HDDC = woodford's DC diss; JSR = Marquardt's JS Revelations; TCHC = Vogel HC
			
			\begin{description}
				\item[Online Access] \href{https://www.josephsmithpapers.org/paper-summary/journal-1835-1836/20}{Here}
				% \item[Analysis] \hyperref[XX]{Here}
			\end{description}
			
			\begin{description}
				\item[Background]
			\end{description}
			
				% It might also be useful to give a two sentence context of the period: e.g., "This speech was given in Nauvoo to a small congregation one month prior to the destruction of the Nauvoo Expositor. These and other tensions may have been on the prophet's mind when he gave the speech."
			
			\begin{description}
				\item[Original Source]
			\end{description}
			
				% brief description of original source material, with reference to JSP or somewhere else for more info
				% Background material: it might be helpful to have a note that says whether a given document was presented as a speech, was written in a journal, etc.
				% Also a comment here about how reliable the source might be, or what shape it is in, if there are large omissions or parts that are hard to understand, and so on.
			
			\begin{description}
				\item[Text]
			\end{description}
			
				I then went to assist in organizing the Elders School called to order and I made some remarks upon the object of this School, and the great necessity there is \sout{in} <of> our rightly improving our time and reigning up our minds to \sout{the} a sense of the great object that lies before us, viz, that glorious endowment that God has in store for the faithful I then dedicated the School in the name of the Lord Jesus Christ. after the School was dismissed I attended a patriarchal meeting at Br Samuel Smiths, his wifeses parents were blessed also his child \& named Susan[n]ah, at evening I preachd at the School-house to a crowded congregation
				
		% -----
		\subsubsection{Record, 6 November 1835}
		% -----
			\label{EMS-1835-JS-6-NOV-1835}
			% \label{EMS-1830-JS-DAY-3LETTERMONTH-YEAR}
			
			\begin{description}
				\item[Print Sources]
			\end{description}

				\begin{tabular}{p{0.5in}p{1in}p{0.5in}p{1in}}
				JSP	 	& J1:85			& JSR 		& ---  \\
				DC	 	& --- 			& HC 		& ---  \\
				HDDC 	& ---			& TCHC 		& ---  \\
				PJS 	& 1:122--123	& TPJS 		& 89   \\
				\end{tabular}
				% HDDC = woodford's DC diss; JSR = Marquardt's JS Revelations; TCHC = Vogel HC
			
			\begin{description}
				\item[Online Access] \href{https://www.josephsmithpapers.org/paper-summary/journal-1835-1836/21}{Here}
				% \item[Analysis] \hyperref[XX]{Here}
			\end{description}
			
			\begin{description}
				\item[Background]
			\end{description}
			
				% It might also be useful to give a two sentence context of the period: e.g., "This speech was given in Nauvoo to a small congregation one month prior to the destruction of the Nauvoo Expositor. These and other tensions may have been on the prophet's mind when he gave the speech."
			
			\begin{description}
				\item[Original Source]
			\end{description}
			
				% brief description of original source material, with reference to JSP or somewhere else for more info
				% Background material: it might be helpful to have a note that says whether a given document was presented as a speech, was written in a journal, etc.
				% Also a comment here about how reliable the source might be, or what shape it is in, if there are large omissions or parts that are hard to understand, and so on.
			
			\begin{description}
				\item[Text]
			\end{description}
				
				Friday morning 6\supers{th} at home. attended School during the school hours returned and spent the evening at home I was this morning introduced to a man from the east, after hearing my name he \sout{replied} remarked that I was nothing but a man: indicating by this expression that he had supposed that a person, <to> who<m> the Lord should see fit to reveal his will, must be something more than a man, he seems to have forgotten the saying that fell from the lips of St. James, that Elias was a man of like passions like unto us, yet he had such power with God that He in answer to his prayer, shut the heavens that they gave no rain for the space of three years and six months, and again in answer to his prayer the heavens gave forth rain and the earth brought forth fruit; and indeed such is the darkness \& ignorance of this generation that they look upon it as incredible that a man should have any intercourse with his Maker.
				
		% -----
		\subsubsection{Revelation, 7 November 1835}
		% -----
			\label{EMS-1835-JS-7-NOV-1835}
			% \label{EMS-1830-JS-DAY-3LETTERMONTH-YEAR}
			
			\begin{description}
				\item[Print Sources]
			\end{description}

				\begin{tabular}{p{0.5in}p{1in}p{0.5in}p{1in}}
				JSP	 	& D5:36--37; J1:85		& JSR 		& 275 \\
				DC	 	& --- 			& HC 		& 2:302--303 \\
				HDDC 	& ---			& TCHC 		& 2:294 \\
				PJS 	& 1:123			& 	 		&  \\
				\end{tabular}
				% HDDC = woodford's DC diss; JSR = Marquardt's JS Revelations; TCHC = Vogel HC
			
			\begin{description}
				\item[Online Access] \href{https://www.josephsmithpapers.org/paper-summary/revelation-7-november-1835/1}{Here}
				% \item[Analysis] \hyperref[XX]{Here}
			\end{description}
			
			\begin{description}
				\item[Background]
			\end{description}
			
				% It might also be useful to give a two sentence context of the period: e.g., "This speech was given in Nauvoo to a small congregation one month prior to the destruction of the Nauvoo Expositor. These and other tensions may have been on the prophet's mind when he gave the speech."
			
			\begin{description}
				\item[Original Source]
			\end{description}
			
				% brief description of original source material, with reference to JSP or somewhere else for more info
				% Background material: it might be helpful to have a note that says whether a given document was presented as a speech, was written in a journal, etc.
				% Also a comment here about how reliable the source might be, or what shape it is in, if there are large omissions or parts that are hard to understand, and so on.
			
			\begin{description}
				\item[Text]
			\end{description}
			
				The word <of the Lord> came to me saying, behold I am well pleased with my servant Isaac Morley and my servant Edward Partridge, because of the integrity of their harts in laboring in my vinyard for the salvation of the souls of men, Verily I say unto you their sins are forgiven them, therefor say unto them in my name that it is my will that they should tarry for a little season and attend the school, and also the solem assembly for a wise purpose in me, even so amen
				
		% -----
		\subsubsection{Revelation, 8 November 1835}
		% -----
			\label{EMS-1835-JS-8-NOV-1835}
			% \label{EMS-1830-JS-DAY-3LETTERMONTH-YEAR}
			
			\begin{description}
				\item[Print Sources]
			\end{description}

				\begin{tabular}{p{0.5in}p{1in}p{0.5in}p{1in}}
				JSP	 	& D5:38--39; J1:86	& JSR 		& 275 \\
				DC	 	& --- 			& HC 		& 2:304 \\
				HDDC 	& ---			& TCHC 		& 2:295 \\
				PJS 	& 1:124			& 	 		&  \\
				\end{tabular}
				% HDDC = woodford's DC diss; JSR = Marquardt's JS Revelations; TCHC = Vogel HC
			
			\begin{description}
				\item[Online Access] \href{https://www.josephsmithpapers.org/paper-summary/revelation-8-november-1835/1}{Here}
				% \item[Analysis] \hyperref[XX]{Here}
			\end{description}
			
			\begin{description}
				\item[Background]
			\end{description}
			
				% It might also be useful to give a two sentence context of the period: e.g., "This speech was given in Nauvoo to a small congregation one month prior to the destruction of the Nauvoo Expositor. These and other tensions may have been on the prophet's mind when he gave the speech."
			
			\begin{description}
				\item[Original Source]
			\end{description}
			
				% brief description of original source material, with reference to JSP or somewhere else for more info
				% Background material: it might be helpful to have a note that says whether a given document was presented as a speech, was written in a journal, etc.
				% Also a comment here about how reliable the source might be, or what shape it is in, if there are large omissions or parts that are hard to understand, and so on.
			
			\begin{description}
				\item[Text]
			\end{description}
			
				The word of the Lord came unto me saying that President [William W.] Phelps \& President J[ohn] Whitmer are under condemnation before the Lord, for their iniquities
				
		% -----
		\subsubsection{Conversations with Robert Matthews, 9--11 November 1835}
		% -----
			\label{EMS-1835-JS-9-11-NOV-1835}
			% \label{EMS-1830-JS-DAY-3LETTERMONTH-YEAR}
			
			\begin{description}
				\item[Print Sources]
			\end{description}

				\begin{tabular}{p{0.5in}p{1in}p{0.5in}p{1in}}
				JSP	 	& D5:39--47; J1:87--95		& JSR 		& --- \\
				DC	 	& --- 			& HC 		& 2:304--307 \\
				HDDC 	& ---			& TCHC 		& 2:295--299 \\
				TPJS 	& 103--105		& 	 		&  \\
				PJS 	& 1:124--132			& 	 		&  \\
				\end{tabular}
				% HDDC = woodford's DC diss; JSR = Marquardt's JS Revelations; TCHC = Vogel HC
			
			\begin{description}
				\item[Online Access] \href{https://www.josephsmithpapers.org/paper-summary/conversations-with-robert-matthews-9-11-november-1835/1}{Here}
				% \item[Analysis] \hyperref[XX]{Here}
			\end{description}
			
			\begin{description}
				\item[Background]
			\end{description}
			
				% It might also be useful to give a two sentence context of the period: e.g., "This speech was given in Nauvoo to a small congregation one month prior to the destruction of the Nauvoo Expositor. These and other tensions may have been on the prophet's mind when he gave the speech."
			
			\begin{description}
				\item[Original Source]
			\end{description}
			
				% brief description of original source material, with reference to JSP or somewhere else for more info
				% Background material: it might be helpful to have a note that says whether a given document was presented as a speech, was written in a journal, etc.
				% Also a comment here about how reliable the source might be, or what shape it is in, if there are large omissions or parts that are hard to understand, and so on.
				
				Note that the J1 account may have more information about their conversation, because it seems to include some of what Matthews said; compare them.
			
			\begin{description}
				\item[Text]
			\end{description}
			
					\hypertarget{EMS-1835-JS-9-11-NOV-1835-a}%
				while
					\marginnote{a}%
				setting in my house between the hours of \sout{nine} <ten> \& \sout{10} 11 this morning a man came in, and introduced himself to me, calling <himself> \sout{self} <by the name of> Joshua the Jewish minister, his appearance was some \sout{what} <thing> singular, having a beard about 3 inches in length which is quite grey, also his hair is long and considerably silvered with age I should think he is about 50 or 55 years old, tall and strait slender built of thin visage blue eyes, and fair complexion, he wears a sea green frock coat, \& pantaloons of the same, black fur hat with narrow brim, and while speaking frequently shuts his eyes, with a scowl on his countinance; I made some enquiry after his name but received no definite answer; we soon commenced talking upon the subject of religion and after I had made some remarks concerning the bible I commenced giving him a relation of the circumstances connected with the coming forth of the book of Mormon, as follows--- being wrought up in my mind, respecting the subject of religion and looking \sout{upon} <at> the different systems taught the children of men, I knew not who was right or who was wrong and concidering it of the first importance that I should be right, in matters that involve\sout{d} eternal consequences; being thus perplexed in mind I retired to the silent grove and bowd down before the Lord, under a realising sense that he had said (if the bible be true) ask and you shall recieve knock and it shall be opened seek and you shall find and again, if any man lack wisdom let him ask of God who giveth to all men liberally and upbra[i]deth not; information was what I most desired at this time, and with a fixed determination \sout{I} to obtain it,
					\hypertarget{EMS-1835-JS-9-11-NOV-1835-b}%
				I
					\marginnote{b}%
				called upon the Lord for the first time, in the place above stated or in other words I made a fruitless attempt to pray, my toung seemed to be swolen in my mouth, so that I could not utter, I heard a noise behind me like some person walking towards me, <I> strove again to pray, but could not, the noise of walking seemed to draw nearer, I sprung up on my feet, \sout{and} and looked around, but saw no person or thing that was calculated to produce the noise of walking, I kneeled again my mouth was opened and my toung liberated, and I called on the Lord in mighty prayer, a pillar of fire appeared above my head, it presently rested down upon \sout{my} <me> \sout{head}, and filled me with joy unspeakable, a personage appeard in the midst, of this pillar of flame which was spread all around, and yet nothing consumed, another personage soon appeard like unto the first, he said unto me thy sins are forgiven thee, he testifyed unto me that Jesus Christ is the son of God; <and I saw many angels in this vision> I was about 14. years old when I recieved this first communication;
					\hypertarget{EMS-1835-JS-9-11-NOV-1835-c}%
				When
					\marginnote{c}%
				I was about 17 years old I saw another vision of angels, in the night season after I had retired to bed I had not been a sleep, \sout{when} but was meditating upon my past life and experiance, I was verry concious that I had not kept the commandments, and I repented hartily for all my sins and transgression, and humbled myself before Him; <whose eyes are over all things>, all at once the room was iluminated above the brightness of the sun an angel appeared before me, his hands and feet were naked pure and white, and he stood between the floors of the room, clothed <with> \sout{in} purity inexpressible, he said unto me I am a messenger sent from God, be faithful and keep his commandments in all things, he told me of a sacred record which was written on plates of gold, I saw in the vision the place where they were deposited, he said the indians, were the literal descendants of Abraham he explained many \sout{things} of the prophesies to me, one I will mention which is \sout{this} in Malachi 4 behold the day of the Lord cometh \&c; also that the Urim and Thumim, was hid up with the record, and that God would give me power to translate it, with the assistance of this instrument he then gradually vanished out of my sight, or the vision closed, while meditating on what I had seen, the Angel appeard to me again and related the same things and much more, also the third time bearing the same tidings, and departed; during the time I was in this vision I did not realize any thing \sout{else} around me except what was shown me in this communication: after the vision had all passed, I found that it was nearly day-light, the family soon arose, I got up also:---
					\hypertarget{EMS-1835-JS-9-11-NOV-1835-d}%
				on
					\marginnote{d}%
				that day while in the field at work with my Father he asked me if I was sick I replyed, I had but little strenght, he told me to go to the house, I started and went part way and was finally deprived \sout{deprived} of my strength and fell, but how long I remained I do not know; the Angel came to me again and commanded me to go and tell my Father, what I had seen and heard, I did so, he wept and told me that it was a vision from God to attend to it I went and found the place, where the plates were, according to the direction of the Angel, also saw them, and the angel as before; the powers of darkness strove hard against me, I called on God, the Angel told me that the reason why I could not obtain the plates at this time was because I was under transgression, but to come again in one year from that time, I did so, but did not obtain them also the third and the fourth year, at which time I obtained them, and translated them into the english language; by the gift and power of God and have been preaching it ever since.

					\hypertarget{EMS-1835-JS-9-11-NOV-1835-e}%
				While
					\marginnote{e}%
				I was relating this brief history of the establishment of the Church of Christ in these last days, Joshua seemed to be highly entertained after I had got through I observed that, the hour of worship \& time to dine had now arived and invited him to tarry, which he concented to,

				After dinner the conversation was resumed and Joshua proceded to make some remarks on the prophesies, as follows:

				He observed that he was aware that I could bear stronger meat than many others, therefore he should open his mind the more freely;--- Daniel has told us that he is to stand in his proper lot, in the latter days according to his vision he had a right to shut it up and also to open it again after many days, or in the latter times; Daniels Image whose head was gold, and body, armes, legs and feet was composed of the different materials described in his vision represents different governments, the golden head was <to represent> Nebuchodnazer King of Babylon, the other parts other kings \& forms of government, which I shall not now mention in detail, but confine my remarks, more particularly to the feet of the Image; The policy of the wicked spirit, is to separate what God has joined togather and unite what He has separated, which he has succeded in doing to admiration, in the present state of society, which is like unto Iron and clay, there is confusion in all things, both both Political and religious, and notwithstanding all the efforts that are made to bring about a union, society \sout{is} remains disunited, and all attempts to <unite her> are as fruitless, as to attemp to unite Iron \& Clay.

					\hypertarget{EMS-1835-JS-9-11-NOV-1835-f}%
				The
					\marginnote{f}%
				feet of the Image, is the government of these united States, other Nations \& kingdoms are looking up to her for an example, of union fredom and equal rights, and therefore worship her, like as Daniel saw in the vision, although they are begining to loose confidence in her, seeing the broils and discord that distract, her political \& religious horizon this Image is characteristic of all governments and institutions or most of them; as they begin with a head of gold and terminate in the contemp[t]ible feet of Iron \& clay: making a splendid appearance at first, proposing to do much more than the[y] can perform, and finally end in degradation and sink, in infamy; we should not only start to come out of Babylon but leav it entirely lest we are overthrown in her ruins, we should keep improving and reforming, twenty-fours hours for improvement now is worth as much as a year a hundred years ago; the spirit of the Fathers that was cut down, or those that were under the altar, are now rising this is the first resurection the Elder that fall's first will rise last; we should not form any opinion only for the present, and leave the result of futurity with God: I have risen up out of obscurity, but was look\supers{d.} up to when but a youth, in temporal things: It is not necessary that God should give us all things at first or in his first commission to us, but in his second. John saw the angel deliver the gospel in the last days, which would not be necessary if it was already in the world this expression would be inconsistent, the small lights that God has given, is sufficient to lead us out of babylon, when we get out we shall have the greater light. I told Jo[s]hua that I did not understand him concerning the resurection and wishd him to be more explanitory on the subject; he replied that he did not feell impressed by the spirit to unfold it further at present, but perhaps he might at some other time.

					\hypertarget{EMS-1835-JS-9-11-NOV-1835-g}%
				I
					\marginnote{g}%
				then withdrew to do some buisness with another gentleman that called to see me.

				He [Robert Matthews] informed my Scribe that he was born in Washington County Town of Cambridge New York. he says that all the railroads canals and other improvements are performed by spirits of the resurection.

				The silence spoken of by John the Revelator which is to be in heaven for the space of half an hour, is between 1830 \& 1851, during which time the judgments of God will be poured out after that time there will be peace.

				Curiosity to see a man that was reputed to be a jew caused many to call during the day and more particularly at evening suspicions were entertained that said Joshua was the noted Mathias of New York, spoken so much of in the public prints on account of the trials he underwent in that place before a court of justice, for murder manslaughter comtempt of court whip[p]ing his Daughter \&c for the two last crimes he was imprisoned, and came out about 4, months since, after some, equivocating he confessed that he was realy Mathias: after supper I proposed that he should deliver a lecture to us, he did so sitting in his chair; he commenced by saying God said let there be light and there was light, which he dwelt upon through his discource, he made some verry exelent remarks but his mind was evidently filled with darkness, after he dismissed his meeting, and the congregation disperced, he conversed freely upon the circumstances that transpired in New York,

					\hypertarget{EMS-1835-JS-9-11-NOV-1835-h}%
				His
					\marginnote{h}%
				name is Robert Mathias, he say[s] that Joshua, is his priestly name.

				during all this time I did not contradict his sentiments, wishing to draw out all that I could concerning his faith; the next morning Tuesday 10\supers{th} I resumed the conversation and desired him to enlighten my mind more on his views respecting the resurection, he says that he poss[ess]es the spirit of his fathers, that he is a litteral decendant of Mathias the Apostle that was chosen in the place of Judas that fell and that his spirit is resurected in him, and that this is the way or scheme of eternal life, this transmigration of soul or spirit from Father to Son: I told him that his doctrine was of the Devil that he was in reality in possession of wicked and depraved spirit, although he professed to be the spirit of truth, it self, also that he possesses the soul of Christ; he tarried until Wednesday 11.\supers{th}, after breckfast I told him, that my God told me that his God is the Devil, and I could not keep him any longer, and he must depart, and so I for once cast out the Devil in bodily shape, \& I believe a murderer
				
		% -----
		\subsubsection{Discourse, 12 November 1835}
		% -----
			\label{EMS-1835-JS-12-NOV-1835}
			% \label{EMS-1830-JS-DAY-3LETTERMONTH-YEAR}
			
			\begin{description}
				\item[Print Sources]
			\end{description}

				\begin{tabular}{p{0.5in}p{1in}p{0.5in}p{1in}}
				JSP	 	& D5:47--51; J1:96--99		& JSR 		& --- \\
				DC	 	& --- 			& HC 		& 2:308--310 \\
				HDDC 	& ---			& TCHC 		& 2:300--301 \\
				TPJS 	& 89--92			& 	 		&  \\
				PJS 	& 1:132--135			& 	 		&  \\
				\end{tabular}
				% HDDC = woodford's DC diss; JSR = Marquardt's JS Revelations; TCHC = Vogel HC
			
			\begin{description}
				\item[Online Access] \href{https://www.josephsmithpapers.org/paper-summary/discourse-12-november-1835/1}{Here}
				% \item[Analysis] \hyperref[XX]{Here}
			\end{description}
			
			\begin{description}
				\item[Background]
			\end{description}
			
				% It might also be useful to give a two sentence context of the period: e.g., "This speech was given in Nauvoo to a small congregation one month prior to the destruction of the Nauvoo Expositor. These and other tensions may have been on the prophet's mind when he gave the speech."
			
			\begin{description}
				\item[Original Source]
			\end{description}
			
				% brief description of original source material, with reference to JSP or somewhere else for more info
				% Background material: it might be helpful to have a note that says whether a given document was presented as a speech, was written in a journal, etc.
				% Also a comment here about how reliable the source might be, or what shape it is in, if there are large omissions or parts that are hard to understand, and so on.
			
			\begin{description}
				\item[Text]
			\end{description}
			
					\hypertarget{EMS-1835-JS-12-NOV-1835-a}%
				This
					\marginnote{a}%
				evening viz the 12\supers{th} at 6 oclock meet with the council of 12. by their request, 9 of them were present council opened by singing \& prayer, and I made some remarks as follows;--- I am happy in the enjoyment of this opportunity of meeting with this council on this occasion, I am satisfyed that the spirit of the Lord is here, and I am satisfied with all the breth[r]en present, and I need not say that you have my utmost confidence, and that I intend to uphold, you to the uttermost, for I am well aware that you \sout{do} \sout{and} \sout{delight} \sout{in} \sout{so} \sout{doing} have to sustain my character \sout{my} \sout{charcter} against the vile calumnies and reproaches of this ungodly generation and that you delight in so doing:--- darkness prevails, at this time as it was, at the time Jesus Christ was about to be crucified, the powers of darkness strove to obscure the glorious sun of righteousness that began to dawn upon the world, and was soon to burst in great blessings upon the heads of the faithful, and let me tell you brethren that great blessings awate us at this time and will soon be poured out upon us if we are faithful in all things, for we are even entitled to greater blessings than they were, because the[y] had the person of Christ with them, to instruct them in the great plan of salvation, his personal presence we have not, therefore we need great faith on account of our peculiar circumstances and I am determined to do all that I can to uphold you, although I may do many things <invertaintly [inadvertently]> that are not right in the sight of God;
					\hypertarget{EMS-1835-JS-12-NOV-1835-b}%
				you
					\marginnote{b}%
				want to know many things that are before you, that you may know how \sout{how} to prepare your selves for the great things that God is about to bring to pass; but there is on[e] great deficiency or obstruction, in the way that deprives us of the greater blessings, and in order to make the foundation of this church complete and permanent, we must remove this obstruction, which is to attend to certain duties that we have not as yet attended to; I supposed I had established this church on a permanent foundation when I went to the Missourie and indeed I did so, for if I had been taken away it would have been enough, but I yet live, and therefore God requires more at my hands:--- The item to which I wish the more particularly to call your attention to night is the ordinance of washing of feet, this we have not done as yet but it is necessary now as much as it was in the days of the Saviour, and we must have a place prepared, that we may attend to this ordinance, aside from the world; we have not desired much from the hand of the Lord, with that faith and obediance that we ought, yet we have enjoyed great blessings, and we are not so sensible of this as we should be; When or wher has God suffered one of the witnesses or first Elders of this church <to> fall? never nor nowhere amidst all the calamities and judgments that have befallen the inhabitants of the earth his almighty arm has sustained us,
					\hypertarget{EMS-1835-JS-12-NOV-1835-c}%
				men
					\marginnote{c}%
				and Devils have raged and spent the malice in vain. we must have all things prepared and call our solem assembly as the Lord has commanded us, that we may be able to accomplish his great work: and it must be done in Gods own way, the house of the Lord must be prepared, and the solem assembly called and organized in it according to the order of the house of God and in it we must attend to the ordinance of washing of feet; it was never intended for any but official members, it is calculated to unite our hearts, that we may be one in feeling and sentiment and that our faith may be strong, so that satan cannot over throw us, nor have any power over us,--- the endowment you are so anxious about you cannot comprehend now, nor could Gabriel explain it to the understanding of your dark minds, but strive to be prepared in your hearts, be faithful in all things that when we meet in the solem assembly that is such as God shall name out of all the official members, will meet, and we must be clean evry whit, let us be faithful and silent brethren, <and> if God gives you a manifestation, keep it to yourselves, be watchful and prayerful, and you shall have a prelude of those joys that God will pour out on that day, do not watch for iniquity in each other if you do you will not get an endowment for God will not bestow it on such;
					\hypertarget{EMS-1835-JS-12-NOV-1835-d}%
				but
					\marginnote{d}%
				if we are faithful and live by every word that procedes forth from the mouth of God I will venture to prophesy that we shall get a blessing that will be worth remembering if we should live as long as John the Revelator, our blessings will be such as we have not realized before, nor in this generation. The order of the house of God has and ever will be the same, even after Christ comes, and after the termination of the thousand years it will be the same, and we shall finally roll into the celestial kingdom of God and enjoy it forever;--- you need an endowment brethren in order that you may be prepared and able to overcome all things, and those that reject your testimony will be damned the sick will be healed the lame made to walk the deaf to hear and the blind to see through your instrumentality;

					\hypertarget{EMS-1835-JS-12-NOV-1835-e}%
				But
					\marginnote{e}%
				let me tell you that you will not have power after the endowment to heal those who have not faith, nor to benifit them, for you might as well expect to benefit a devil in hell as such a<n> one, who is possessed of his spirit and are willing to keep it for they are habitations for devils and only fit for his society but when you are endowed and prepared to preach the gospel to all nations kindred and toungs in there own languages you must faithfully warn all and bind up the testimony and seal up the law and the destroying angel will follow close at your heels and execute his tremendeous mission upon the children of disobediance, and destroy the workers of iniquity, while the saints will be gathered out from among them and stand in holy places ready to meet the bride groom when he comes.---

				I feel disposed to speak a few words more to you my brethren concerning the endowment, all who are prepared and are sufficiently pure to abide the presence of the Saviour will see him in the solem assembly.
				
		% -----
		\subsubsection{Revelation, 14 November 1835}
		% -----
			\label{EMS-1835-JS-14-NOV-1835}
			% \label{EMS-1830-JS-DAY-3LETTERMONTH-YEAR}
			
			\begin{description}
				\item[Print Sources]
			\end{description}

				\begin{tabular}{p{0.5in}p{1in}p{0.5in}p{1in}}
				JSP	 	& D5:51--53; J1:99--100		& JSR 		& 276 \\
				DC	 	& --- 			& HC 		& 2:311--312 \\
				HDDC 	& ---			& TCHC 		& 2:303--304 \\
				PJS 	& 1:136			& 	 		&  \\
				\end{tabular}
				% HDDC = woodford's DC diss; JSR = Marquardt's JS Revelations; TCHC = Vogel HC
			
			\begin{description}
				\item[Online Access] \href{https://www.josephsmithpapers.org/paper-summary/revelation-14-november-1835/1}{Here}
				% \item[Analysis] \hyperref[XX]{Here}
			\end{description}
			
			\begin{description}
				\item[Background]
			\end{description}
			
				% It might also be useful to give a two sentence context of the period: e.g., "This speech was given in Nauvoo to a small congregation one month prior to the destruction of the Nauvoo Expositor. These and other tensions may have been on the prophet's mind when he gave the speech."
			
			\begin{description}
				\item[Original Source]
			\end{description}
			
				% brief description of original source material, with reference to JSP or somewhere else for more info
				% Background material: it might be helpful to have a note that says whether a given document was presented as a speech, was written in a journal, etc.
				% Also a comment here about how reliable the source might be, or what shape it is in, if there are large omissions or parts that are hard to understand, and so on.
			
			\begin{description}
				\item[Text]
			\end{description}
			
				Thus came the word of the Lord unto me saying:
				
				verily thus saith the \sout{the} Lord unto my servant Joseph concerning my servant Warren [Parrish], behold his sins are forgiven him because of his desires to do the works of righteousness therefore in as much as he will continue to hearken unto my voice he shall be blessed with wisdom and with a sound mind even above his fellows, behold it shall come to pass in his day that he shall <see> great things shew forth themselves unto my people, he shall see much of my ancient records, and shall know of hid[d]en things, and shall be endowed with a knowledge of hiden languages, and if he desires and shall seek it at my hand, he shall be privileged with writing much of my word, as a scribe unto me for the benefit of my people, therefore this shall be his calling until I shall order it otherwise in my wisdom and it shall be said of him in a time to come, behold Warren the Lords Scribe, for the Lords Seer whom he hath appointed in Israel; Therefore <if he will> keep my commandments he shall be lifted up at the last day, even so Amen
				
		% -----
		\subsubsection{Record, 14 November 1835}
		% -----
			\label{EMS-1835-JS-14-NOV-1835-REC}
			% \label{EMS-1830-JS-DAY-3LETTERMONTH-YEAR}
			
			\begin{description}
				\item[Print Sources]
			\end{description}

				\begin{tabular}{p{0.5in}p{1in}p{0.5in}p{1in}}
				JSP	 	& J1:100		& JSR 		& --- \\
				DC	 	& --- 			& HC 		& --- \\
				HDDC 	& ---			& TCHC 		& --- \\
				PJS 	& 1:136--137			& 	 		&  \\
				\end{tabular}
				% HDDC = woodford's DC diss; JSR = Marquardt's JS Revelations; TCHC = Vogel HC
			
			\begin{description}
				\item[Online Access] \href{https://www.josephsmithpapers.org/paper-summary/journal-1835-1836/37}{Here}
				% \item[Analysis] \hyperref[XX]{Here}
			\end{description}
			
			\begin{description}
				\item[Background]
			\end{description}
			
				% It might also be useful to give a two sentence context of the period: e.g., "This speech was given in Nauvoo to a small congregation one month prior to the destruction of the Nauvoo Expositor. These and other tensions may have been on the prophet's mind when he gave the speech."
			
			\begin{description}
				\item[Original Source]
			\end{description}
			
				% brief description of original source material, with reference to JSP or somewhere else for more info
				% Background material: it might be helpful to have a note that says whether a given document was presented as a speech, was written in a journal, etc.
				% Also a comment here about how reliable the source might be, or what shape it is in, if there are large omissions or parts that are hard to understand, and so on.
			
			\begin{description}
				\item[Text]
			\end{description}
			
				A Gentleman called this after noon by the name of Erastus Holmes of Newbury Clemon [Clermont] Co. Ohio, he called to make enquiry about the establishment of the Church of the latter-day Saints and to be instructed more perfectly in our doctrine \&c I commenced and gave him a brief relation of my experience while in my juvenile years, say from 6, years old up to the time I received the first visitation of Angels which was when I was about 14, years old and also the the visitations that I received afterward, concerning the book of Mormon, and a short account of the rise and progress of the church, up to this, date he listened verry attentively and seemed highly gratified, and intends to unite with the Church he is a verry candid man indeed and I am much pleased with him.
				
		% -----
		\subsubsection{Letter to the Elders of the Church, 16 November 1835}
		% -----
			\label{EMS-1835-JS-16-NOV-1835}
			% \label{EMS-1830-JS-DAY-3LETTERMONTH-YEAR}
			
			\begin{description}
				\item[Print Sources]
			\end{description}

				\begin{tabular}{p{0.5in}p{1in}p{0.5in}p{1in}}
				JSP	 	& D5:53--60		& JSR 		& --- \\
				DC	 	& --- 			& HC 		& 2:259--264 \\
				HDDC 	& ---			& TCHC 		& 2:277--282 \\
				TPJS 	& 83--89			& 	 		&  \\
				PJS 	& 1:137			& 	 		&  \\
				\end{tabular}
				% HDDC = woodford's DC diss; JSR = Marquardt's JS Revelations; TCHC = Vogel HC
			
			\begin{description}
				\item[Online Access] \href{https://www.josephsmithpapers.org/paper-summary/letter-to-the-elders-of-the-church-16-november-1835/1}{Here}
				% \item[Analysis] \hyperref[XX]{Here}
			\end{description}
			
			\begin{description}
				\item[Background]
			\end{description}
			
				% It might also be useful to give a two sentence context of the period: e.g., "This speech was given in Nauvoo to a small congregation one month prior to the destruction of the Nauvoo Expositor. These and other tensions may have been on the prophet's mind when he gave the speech."
			
			\begin{description}
				\item[Original Source]
			\end{description}
			
				% brief description of original source material, with reference to JSP or somewhere else for more info
				% Background material: it might be helpful to have a note that says whether a given document was presented as a speech, was written in a journal, etc.
				% Also a comment here about how reliable the source might be, or what shape it is in, if there are large omissions or parts that are hard to understand, and so on.
			
			\begin{description}
				\item[Text]
			\end{description}
			
					\hypertarget{EMS-1835-JS-16-NOV-1835-a}%
				\phantom{ }\textit{To the elders of the church of the Latter Day Saints}.
					\marginnote{a}

				At the close of my letter in the September No. of the ``Messenger and Advocate,'' I promised to continue the subject there commenced: I do so with a hope that it may be a benefit and a means of assistance to the elders in their labors, while they are combatting the prejudices of a crooked and perverse generation, by having in their possession, the facts of my religious principles, which are misrepresented by almost all those whose crafts are in danger by the same; and also to aid those who are anxiously inquiring, and have been excited to do so from rumor, in accertaining correctly, what my principles are.

				I have been drawn into this course of proceeding, by persecution, that is brought upon us from false rumor, and misrepresentations concerning my sentiments.

					\hypertarget{EMS-1835-JS-16-NOV-1835-b}%
				But
					\marginnote{b}%
				to proceed, in the letter alluded to. The principles of repentance and baptism for the remission of sins, are not only set forth, but many passages of scripture, were quoted, clearly illucidating the subject; let me add, that I do positively rely upon the truth and veracity of those principles inculcated in the new testament; and then pass from the above named items, on to the item or subject of the gathering, and show my views upon this point: which is an item which I esteem to be of the greatest importance to those who are looking for salvation in this generation, or in these what may be called ``the latter times,'' as all the prophets that have written, from the days of righteous Abel down to the last man, that has left any testimony on record, for our consideration, in speaking of the salvation of Israel in the last days, goes directly to show, that it consists in the work of the gathering.

					\hypertarget{EMS-1835-JS-16-NOV-1835-c}%
				Firstly,
					\marginnote{c}%
				I shall begin by quoting from the prophecy of Enoch, speaking of the last days: ``Righteousness will I send down out of heaven, and truth will I send forth out of the earth, to bear testimony of mine Only Begotten, his resurrection from the dead, -[this resurrection I understand to be the corporeal body]- yea, and also the resurrection of all men, righteousness and truth will I cause to sweep the earth as with a flood, to gather out mine own elect from the four quarters of the earth, unto a place which I shall prepare; a holy city, that my people may gird up their loins, and be looking forth for the time of my coming: for there shall be my tabernacle; and it shall be called Zion, a New Jerusalem.''

					\hypertarget{EMS-1835-JS-16-NOV-1835-d}%
				Now
					\marginnote{d}%
				I understand by this quotation, that God clearly manifested to Enoch, the redemption which he prepared, by offering the Messiah as a Lamb slain from before the foundation of the world: by virtue of the same, the glorious resurrection of the Savior, and the resurrection of all the human family,---even a resurrection of their corporeal bodies: and also righteousness and truth to sweep the earth as with a flood. Now I ask how righteousness and truth are agoing to sweep the earth as with a flood? I will answer:---Men and angels are to be co-workers in bringing to pass this great work: and a Zion is to be prepared; even a New Jerusalem, for the elect that are to be gathered from the four quarters of the earth, and to be established an holy city: for the tabernacle of the Lord shall be with them.

				Now Enoch was in good company in his views upon this subject. See Revelations, 23:3.---``And I heard a great voice out of heaven saying, Behold the tabernacle of God is with men, and he will dwell with them, and they shall be his people, and God himself shall be with them, and be their God.'' I discover by this quotation, that John upon the isle of Patmos, saw the same things concerning the last days, which Enoch saw. But before the tabernacle can be with men, the elect must be gathered from the four quarters of the earth.

					\hypertarget{EMS-1835-JS-16-NOV-1835-e}%
				And
					\marginnote{e}%
				to show further upon this subject of the gathering: Moses, after having pronounced the blessing and the cursing upon the children of Israel, for their obedience or disobedience, says thus:---``And it shall come to pass, when all these things are come upon thee, the blessing and the curse which I have set before thee; and thou shalt call them to mind, among all the nations whither the Lord thy God hath driven thee, and shalt return unto the Lord thy God, and shalt obey his voice, according to all that I command thee, this day, thou and thy children, with all thine heart, and with all thy soul, that then the Lord thy God, will turn thy captivity, and have compassion upon thee, and will return and gather thee from all the nations whither the Lord thy God hath scattered thee; and if any of thine be driven out unto the utmost parts of heaven; from thence will the Lord thy God gather thee; and from thence will he fetch thee.''

				It has been said by many of the learned, and wise men, or historians, that the Indians, or aboriginees of this continent, are of the scattered tribes of Israel. It has been conjectured by many others, that the aboriginees of this continent, are not of the tribes of Israel; but the ten tribes have been led away into some unknown regions of the north. Let this be as it may, the prophesy I have just quoted, ``will fetch them'' in the last days, and place them, in the land which their fathers possessed: and you will find in the 7th verse of the 30th chapt. quoted:---``And the Lord thy God will put all these curses upon thine enemies and on them that hate thee, which persecuted thee.''

					\hypertarget{EMS-1835-JS-16-NOV-1835-f}%
				Many
					\marginnote{f}%
				may say that this scripture is fulfilled, but let them mark carefully what the prophet says: ``If any are driven out unto the utmost parts of heaven;'' (which must mean the breadths of the earth.) Now this promise is good to any, if there should be such, that are driven out, even in the last days: therefore, the children of the fathers have claim unto this day: and if these curses are to be laid over on the heads of their enemies, wo be unto the Gentiles: See book of Mormon, page 487, Wo unto the unbelieving of the Gentiles, saith the Father. Again see book of Mormon, page 497, which says: ``Behold this people will I establish in this land, unto the fulfilling of the covenant which I made with your father Jacob: and it shall be a New Jerusalem.'' Now we learn from the book of Mormon, the very identical continent and spot of land upon which the New Jerusalem is to stand, and it must be caught up according to the vision of John upon the isle of Patmos. Now many will be disposed to say, that this New Jerusalem spoken of, is the Jerusalem that was built by the Jews on the eastern continent: but you will see from Revelations, 21:2, there was a New Jerusalem coming down from God out of heaven, adorned as a bride for her husband. That after this the Revelator was caught away in the Spirit to a great and high mountain, and saw the great and holy city descending out of heaven from God. Now there are two cities spoken of here, and as every thing cannot be had in so narrow a compass as a letter, I shall say with brevity, that there is a New Jerusalem to be established on this continent.--- And also the Jerusalem shall be rebuilt on the eastern continent. See book of Mormon, page 566. Behold, Ether saw the days of Christ, and he spake also concerning the house of Israel, and the Jerusalem from whence Lehi should come: after it should be destroyed it should be built up again, a holy city unto the Lord: wherefore, it could not be a New Jerusalem, for it had been in a time of old. This may suffice upon the subject of gathering until my next.

					\hypertarget{EMS-1835-JS-16-NOV-1835-g}%
				I
					\marginnote{g}%
				now proceed, at the close of my letter, to make a few remarks on the duty of elders with regard to their teaching parents and children, husbands and wives, masters and slaves, or servants, \&c. as I said I would in my former letter. And firstly, it becomes an elder when he is travelling through the world, warning the inhabitants of the earth to gather together, that they may be built up an holy city unto the Lord, instead of commencing with children, or those who look up to parents or guardians, to influence their minds, thereby drawing them from their duties, which they rightfully owe to such, they should commence their labors with parents, or guardians, and their teachings should be such as are calculated to turn the hearts of the fathers to the children, and the hearts of the children to the fathers. And no influence should be used, with children contrary to the consent of their parents or guardians.--- But all such as can be persuaded in a lawful and righteous manner, and with common consent, we should feel it our duty to influence them to gather with the people of God. But otherwise let the responsibility rest upon the heads of parents or guardians, and all condemnation or consequences, be upon their heads, according to the dispensation which he hath committed unto us: for God has so ordained, that his work shall be cut short in righteousness, in the last days: therefore, first teach the parents, and then, with their consent, let him persuade the children to embrace the gospel also. And if children embrace the gospel, and their parents or guardians are unbelievers, teach them to stay at home and be obedient to their parents or guardians, if they require it; but if they consent to let them gather with the people of God let them do so and there shall be no wrong and let all things be done carefully, and righteously, and God will extend his guardian care to all such.

					\hypertarget{EMS-1835-JS-16-NOV-1835-h}%
				And
					\marginnote{h}%
				secondly, it should be the duty of elders, when they enter into any house, to let their labors and warning voice, be unto the master of that house: and if he receive the gospel, then he may extend his influence to his wife also, with consent, that peradventure she may receive the gospel; but if a man receive not the gospel, but gives his consent that his wife may receive it, and she believes, then let her receive it. But if the man forbid his wife, or his children before they are of age, to receive the gospel, then it should be the duty of the elder to go his way and use no influence against him: and let the responsibility be upon his head---shake off the dust of thy feet as a testimony against him, and thy skirts shall then be clear of their souls. Their sins are not to be answered upon such as God hath sent to warn them to flee the wrath to come, and save themselves from this untoward generation. The servants of God will not have gone over the nations of the Gentiles, with a warning voice, until the destroying angel will commence to waste the inhabitants of the earth; and as the prophet hath said, ``It shall be a vexation to hear the report.'' I speak because I feel for my fellow-men: I do it in the name of the Lord, being moved upon by the Holy Spirit. O that I could snatch them from the vortex of misery, into which I behold them plunging themselves, by their sins, that I may be enabled, by the warning voice, to be an instrument of bringing them to unfeigned repentance, that they may have faith to stand in the evil day.

					\hypertarget{EMS-1835-JS-16-NOV-1835-i}%
				Thirdly,
					\marginnote{i}%
				it should be the duty of an elder, when he enters into a house to salute the master of that house, and if he gain his consent, then he may preach to all that are in that house, but if he gain not his consent, let him go not unto his slaves or servants, but let the responsibility be upon the head of the master of that house, and the consequences thereof; and the guilt of that house is no longer upon thy skirts: Thou art free; therefore, shake off the dust of thy feet, and go thy way. But if the master of that house give consent, that thou mayest preach to his family, his wife, his children, and his servants, his man-servants, or his maid-servants, or his slaves, then it should be the duty of the elder to stand up boldly for the cause of Christ, and warn that people with one accord, to repent and be baptized for the remission of sins, and for the Holy Ghost, always commanding them in the name of the Lord, in the spirit of meekness to be kindly affected one towards another; that the fathers should be kind to their children, husbands to their wives; masters to their slaves or servants; children obedient to their parents, wives to their husbands, and slaves or servants to their masters:

					\hypertarget{EMS-1835-JS-16-NOV-1835-j}%
				``Wives
					\marginnote{j}%
				submit youselves unto your own husbands, as unto the Lord. For the husband is the head of the wife, even as Christ is the head of the church: and he is the Savior of the body. Therefore as the church is subject unto Christ, so let the wives be to their own husbands in every thing. Husbands, love your wives even as Christ also loved the church and gave himself for it; that he might sanctify and cleanse it with the washing of water by the word, that he might present it to himself a glorious church, not having spot, or wrinkle, or any such thing; but that it should be holy and without blemish. So ought men to love their wives as their own bodies. He that loveth his wife loveth himself. For no man ever yet hated his own flesh; but nourisheth and cherisheth it, even as the Lord the church: for we are members of his body, of his flesh, and of his bones.--- For this cause shall a man leave his father and mother, and shall be joined unto his wife, and they two shall be one flesh.''---Ephesians, Chapt. V. from the 22d to the end of the 21st [31st] verse.

					\hypertarget{EMS-1835-JS-16-NOV-1835-k}%
				``Wives
					\marginnote{k}%
				submit yourselves unto your own husbands, as it is fit in the Lord. Husbands, love your wives, and be not bitter against them. Children, obey your parents in all things: for this is well pleasing unto the Lord. Fathers, provoke not your children to anger, lest they be discouraged. Servants, obey in all things your masters according to the flesh: not with eye service as menpleasers; but in singleness of heart, fearing God.''---Colocians, Chapt. III. from the 18th to the end of the 22d verse.

				But I must close this letter and resume the subject in another number.

				In the bonds of the new and everlasting covenant
				
				\begin{tightright}
				JOSEPH SMITH, jr.
				\end{tightright}

				To J[ohn] \textsc{Whitmer, Esq}.
				
		% -----
		\subsubsection{Letter and Revelation to Harvey Whitlock, 16 November 1835}
		% -----
			\label{EMS-1835-JS-16-NOV-1835-HW}
			% \label{EMS-1830-JS-DAY-3LETTERMONTH-YEAR}
			
			\begin{description}
				\item[Print Sources]
			\end{description}

				\begin{tabular}{p{0.5in}p{1in}p{0.5in}p{1in}}
				JSP	 	& D5:60--62		& JSR 		& 276--277 \\
				DC	 	& --- 			& HC 		& 2:314--316 \\
				HDDC 	& ---			& TCHC 		& 2:313--314 \\
				PJS 	& 1:140--142			& 	 		&  \\
				\end{tabular}
				% HDDC = woodford's DC diss; JSR = Marquardt's JS Revelations; TCHC = Vogel HC
			
			\begin{description}
				\item[Online Access] \href{https://www.josephsmithpapers.org/paper-summary/letter-and-revelation-to-harvey-whitlock-16-november-1835/1}{Here}
				% \item[Analysis] \hyperref[XX]{Here}
			\end{description}
			
			\begin{description}
				\item[Background]
			\end{description}
			
				% It might also be useful to give a two sentence context of the period: e.g., "This speech was given in Nauvoo to a small congregation one month prior to the destruction of the Nauvoo Expositor. These and other tensions may have been on the prophet's mind when he gave the speech."
			
			\begin{description}
				\item[Original Source]
			\end{description}
			
				% brief description of original source material, with reference to JSP or somewhere else for more info
				% Background material: it might be helpful to have a note that says whether a given document was presented as a speech, was written in a journal, etc.
				% Also a comment here about how reliable the source might be, or what shape it is in, if there are large omissions or parts that are hard to understand, and so on.
			
			\begin{description}
				\item[Text]
			\end{description}
			
				\begin{tightright}
						\hypertarget{EMS-1835-JS-16-NOV-1835-HW-a}%
					Kirtland Nov. 16\supers{th} 1835
						\marginnote{a}
				\end{tightright}

				Bro Harvey Whitlock

				I have received your letter of the 28\supers{th} Sept. 1835, and I have read it twice, and it gave me sensations that are better imagined than discribed; let it suffice, that I say the verry flood-gates of my heart were broken up: I could not refrain from weeping, I thank God, that it has entered into your heart, to try to return to the Lord, and to his people; if it so be, that he will have mercy upon you.

					\hypertarget{EMS-1835-JS-16-NOV-1835-HW-b}%
				I
					\marginnote{b}%
				have inquired of the Lord concerning your case, these words came to me

				Verily thus saith the Lord unto you; let him who was my servant Harvey, return unto me;--- and unto the bosom of my Church, and forsake all the sins wherewith he has offended against me and persue from hence forth a virtuous and upright life, and remain under the direction of those whom I have appointed to be pillars, and heads of my Church, and behold, saith the Lord, your God; his sins shall be blotted out from under heaven, and shall be forgotten from among men, and shall not come up in mine ears, nor be recorded as <a> memorial against him, but I will lift him up as out of deep mire, and he shall be exalted upon the high places, and shall be counted worthy to stand ammong princes, and shall yet be made a polished shaft in my quiver, of bringing down the strong holds of wickedness, among those who set themselves up on high, that they may take council against me, and against annointed ones in the last days.

				Therefore let him prepare himself speedily and come unto you; even to Kirtland and inasmuch as he shall harken unto all your council from henceforth he shall be restored unto his former state, and shall be saved unto the uttermost, even as the Lord your God liveth Amen.

					\hypertarget{EMS-1835-JS-16-NOV-1835-HW-c}%
				Thus
					\marginnote{c}%
				you see my dear Brother the willingness of our heavenly Father to forgive sins and restore to favour all those who are willing to humble themselves before him, and confess their sins and forsake them, and return to him with full purpose of heart (acting no hypocrisy) to serve him to the end.

				Marvle not that the Lord has condescended to speak from the heavens and give you instructions whereby you may learn your duty, he has heard your prayers, and witnessed your humility; and holds forth the hand of paternal affection, for your return; the angels rejoice over you, while the saints are willing to recieve you again into fellowship.

				I hope on the rec[e]ipt of this, you will \sout{not} loose \sout{any} no time in coming to Kirtland: for if you get here in season, you will have the privelege of attending the School of the prophets, which has already commenced and, also receive\sout{d} instruction in doctrine, and principle, from those whom God has appointed whereby you may be qualified to go forth, and declare the true doctrines of the kingdom according to the \sout{true} \sout{doctrines} \sout{of} \sout{the} mind and, will of God. and when you come to Kirtland, it will be explained to you why God has condescended to give you a revelation according to your request.

				please give my respects to you[r] family, and bee assured I am yours in the bonds of the new and everlasting covenant

				\begin{tightright}
				Joseph Smith Jun
				\end{tightright}
				
		% -----
		\subsubsection{Revelation, 16 November 1835}
		% -----
			\label{EMS-1835-JS-16-NOV-1835-REV}
			% \label{EMS-1830-JS-DAY-3LETTERMONTH-YEAR}
			
			\begin{description}
				\item[Print Sources]
			\end{description}

				\begin{tabular}{p{0.5in}p{1in}p{0.5in}p{1in}}
				JSP	 	& D5:63--64		& JSR 		& 277 \\
				DC	 	& --- 			& HC 		& 2:316 \\
				HDDC 	& ---			& TCHC 		& 2:314 \\
				PJS 	& 1:142			& 	 		&  \\
				\end{tabular}
				% HDDC = woodford's DC diss; JSR = Marquardt's JS Revelations; TCHC = Vogel HC
			
			\begin{description}
				\item[Online Access] \href{https://www.josephsmithpapers.org/paper-summary/revelation-16-november-1835/1}{Here}
				% \item[Analysis] \hyperref[XX]{Here}
			\end{description}
			
			\begin{description}
				\item[Background]
			\end{description}
			
				% It might also be useful to give a two sentence context of the period: e.g., "This speech was given in Nauvoo to a small congregation one month prior to the destruction of the Nauvoo Expositor. These and other tensions may have been on the prophet's mind when he gave the speech."
			
			\begin{description}
				\item[Original Source]
			\end{description}
			
				% brief description of original source material, with reference to JSP or somewhere else for more info
				% Background material: it might be helpful to have a note that says whether a given document was presented as a speech, was written in a journal, etc.
				% Also a comment here about how reliable the source might be, or what shape it is in, if there are large omissions or parts that are hard to understand, and so on.
			
			\begin{description}
				\item[Text]
			\end{description}
			
				The same night \sout{that} I received the word of the Lord on Mr. Hlmes [Erastus Holmes's] case, he had, desired that I would inquire at the hand of the Lord whether it was his duty to be baptised here, or wait until he returned home;--- The word of the Lord came \sout{to} unto me saying, that Mr. Holmes had better not be baptised here, and that he had better not return by water, also that there were three men that were seeking his destruction, to be ware of his eneys [enemies]
				
		% -----
		\subsubsection{Book of Abraham Manuscript, circa July--circa November 1835--A [Abraham 1:4--2:6]}
		% -----
			\label{EMS-1835-JS-CIR-JUL-CIR-NOV-1835}
			% \label{EMS-1830-JS-DAY-3LETTERMONTH-YEAR}
			
			\begin{description}
				\item[Print Sources]
			\end{description}

				\begin{tabular}{p{0.5in}p{1in}p{0.5in}p{1in}}
				JSP	 	& D5:69--80		& JSR 		& --- \\
				DC	 	& --- 		& HC 		& --- \\
				HDDC 	& ---		& TCHC 		& --- \\
				\end{tabular}
				% HDDC = woodford's DC diss; JSR = Marquardt's JS Revelations; TCHC = Vogel HC
			
			\begin{description}
				\item[Online Access] \href{https://www.josephsmithpapers.org/paper-summary/book-of-abraham-manuscript-circa-july-circa-november-1835-a-abraham-14-26/1}{Here}
				% \item[Analysis] \hyperref[XX]{Here}
			\end{description}
			
			\begin{description}
				\item[Background]
			\end{description}
			
				% It might also be useful to give a two sentence context of the period: e.g., "This speech was given in Nauvoo to a small congregation one month prior to the destruction of the Nauvoo Expositor. These and other tensions may have been on the prophet's mind when he gave the speech."
			
			\begin{description}
				\item[Original Source]
			\end{description}
			
				% brief description of original source material, with reference to JSP or somewhere else for more info
				% Background material: it might be helpful to have a note that says whether a given document was presented as a speech, was written in a journal, etc.
				% Also a comment here about how reliable the source might be, or what shape it is in, if there are large omissions or parts that are hard to understand, and so on.
				
				No attempt here to represent characters in the margins. Also, a discussion of how we aren't including this: https://www.josephsmithpapers.org/paper-summary/egyptian-alphabet-circa-early-july-circa-november-1835-a/1 .
			
			\begin{description}
				\item[Text]
			\end{description}
				
				\begin{tightright}
				<J>
				\end{tightright}

				$\blacklozenge$ [1]	sign of the fifth degree of the \sout{first} <Sec\sout{c}ond> part

				$\blacklozenge$ [2]	I sought for <mine> \sout{the} appointment \sout{whereunto} unto the priesthood according to the appointment of God unto the fathers concerning the seed

				$\blacklozenge$ [3]	my fathers having turned from their righteousness and from unto them unto the worshiping of the Gods of the hethens

				$\blacklozenge$ [4]	utterly refused to harken to my voice for their hearts were set to do evil and were wholly turned to the God of Elk<=>Kener and the God of Zibnah and the God of Mah-mackrah and the God of Pharoah King of Egypt therefore they turned their hearts to the sacrafice of the heathens in offering up their children unto these dumb Idols and harkened not unto my voice but indeovered [endeavored] to take away my life by the hand of the priest of Elk=Kener

				$\blacklozenge$ [5]	The priest of Elk=Keenah was also the priest of Pharoah, now at this time it was the custom of the priest of Pharaoh the King of Egypt to offer up upon the Alter which was built in the land of Chaldea for the offering unto \sout{ther} these strange gods both men, women, and children--- and it came to pass that the priest made an offering unto the god of Pharaoh and also unto the god of Shag=reel even after the manner of the Egyptians now the god of Shag-reel was the Sun--- even a thank offering of a child did the priest of Pharaoh offer upon the Alter which stood by the hill called Potipher<s> hill at the head of the plain\sout{s} of Olishem

				$\blacklozenge$ [6]	Now this priest had offered \sout{off} upon this alter three virgins at one time who were the daughters of \sout{Onitus} Onitah---one of the \sout{regular} royal discent directly from the loins of Ham these virgins were offered up because of their virtue they would not bow down to worship Gods of wood, or of stone therefore they were Killed upon this alter

				$\blacklozenge$ [7]	And it was done after the manner of the Egyptians and it came to pass that the priests laid violence upon me that they might slay me also, as they did those virgins upon this alter, and that you might have a knowledge of this alter <I will refer you to the representation that is at the (commencement of this record>

				$\blacklozenge$ [8]	It was made after, the form of a bedsted such as was had among the Chaldeans and it stood before the Gods of Elk-keenah Zibnah Mah-Mach-rah---and als[o] a God like unto that of pharaoh King of Egypt

				[p. 1]

				<That you may have an understanding of their gods I have given you the fashion of them in the figures at the begining which manner of figures is called by the Chaldians, \sout{Ca$\diamond$ak-lee=} Kah-lee=-nos.------>

				<\uline{K}>

				$\blacklozenge$ [9]	And as they lifted up their hands upon me that they might offer me up \sout{to} \sout{and} <\sout{an} and> take away my life behold I lifted up my voice unto the Lord my \sout{good} God; and the lord harkened, and heard and he filled me with a vision of the almighty and the angel of his presence stood by my feet and immediately loosed my bands

				$\blacklozenge$ [10]	And his voice was unto me. Abram Abram Behold my name is Jehovah. and I have heard thee and have come down to deliver thee. and to take thee away from thy fathers house, and from all thy Kinsfolks, in to a strange land which thou knowest not of, and this because \sout{their} \sout{hearts} \sout{are} \sout{turned} they have turned their hearts away from me to worship the god of Elk Kee-nah and the god of Zibnah- and of Mah-Mach-rah--- and the god of pharaoh King of Egypt. Therefore I have come down to visit them. and to distroy him, who hath lifted up his hand against thee \sout{Abraham} <Abram> my son to \sout{distroy thy} take away thy life, Behold I will lead thee by my hand and I will take thee, to put upon thee my name even the priesthood of thy father, and my power shall be over thee; as it was with Noah so shall it be with thee, that through thy ministry, my name shall be known, in the earth forever, for I am thy God

				$\blacklozenge$ [11]	Behold Potiphers hill was in the land of Ur of Chaldea and the Lord broke down the alter of Elk-Keenah and of the gods of the land, and utterly distroyed them \sout{gods of the land} and smote the priest\sout{s} that he died and there was great morning in Chaldeea and also in the court of Pharaoh which Pharaoh signifies King by royal blood. Now this King of Egypt was a discendent from the loins of Ham and was a partaker of the blood of the Cananitess by birth: From this decent sprang all the Egyptians and thus the blood of the cannites was preserved in the land

				\begin{tightright}
				$\blacklozenge$

				<\uline{L}>
				\end{tightright}

				$\blacklozenge$ [12]	The land of Egypt being first discovered by a woman, who was the daughter of Ham; and the daughter of Zep-tah. which in the Chaldea signifies Egypt, which sign[i]fies that which is forbidden. Whin this woman discovered the land it was under water, who after settled her sons in it: And thus from Ham sprang \sout{the} that race which preserved the curse in the land.

				$\blacklozenge$ [13]	Now the <first> government of Egypt, was established by Pharaoh the eldest \sout{sun} <son> of Egyptes the daughter of Ham; and it was after the manner of the government of Ham, which was Patriarchal. Pharaoh being a righteous man established his kingdom, and Judged his people wisely and Justly all his days, seeking earnestly to imitate that order established by the fathers in the first generation in the days of the first Patriarchal reign, even in the reign of Adam. And also Noah his father\sout{. For} \sout{in} \sout{his} \sout{days} \sout{Who} <who> blessed him with the blessings of the earth, and \sout{of} with the blessings of wisdom, but cursed him as pertaining to the priesthood.

				$\blacklozenge$ [14]	Now Pharaoh being of that leniage [lineage] by which he could not have the right of priesthood; notwithstanding the Pharaohs would fain claim it from Noah through Ham: Therefore, my father was led away by their---idolatry, <;> but I shall indeaver [endeavor] hereafter to dilliniate the chronology run[n]ing back from myself to the begining of <the> creation, for the reccords, have \sout{came} <come> into my hands which I hold unto this present time

				$\blacklozenge$ [15]	Now after the priest of Elkkeenah was smitten that he died, there came a fulfilment of those things which were spoken unto me concerning the land of Chaldea, that there should be a famine in the land; and accordingly a famine prevailed throughout all the land of Chaldea: \sout{and} And my father was sorely tormented because of the famine, and he repented of the evil which he had determined against me, to take away my life: But the reccords of the fathers even the patriarchs concerning the right of priesthood, the lord my god preserved in mine own hand<s:>

				<\uline{M}> Therefore a knowledge of the begining of creation \sout{an} and also of the planets, and of the stars, as it was made known unto the fathers, have I kept even unto this day.------

				$\blacklozenge$ [16]	And I shall endeaver to write some of these things, upon this reccord, for the benefit of my posterity, that shall come after me

				$\blacklozenge$ [17]	Now the Lord God caused the famine to wax soar in the land of \sout{E$\diamond$} Ur insomuch that Haran my brother died: but Terah my father yet lived in the land of Ur of the chaldees. And it came to pass; that I Abram took sarai to wife, and Nahor my brother took Milcah to wife

				$\blacklozenge$ [18]	Who was the daughter of Haron

				$\blacklozenge$[19]	Now the Lord had said unto me Abram get the[e] out of thy country, and from thy Kindred and from thy fathers home, unto a land that I will shew thee: Therefore I left the land of Ur of the chaldees to go into the land of canaan; and I took Lot my brothers son, and his wife, and Sarai my wife; and also my father followed after me unto the land which we denominated Haran. And the famine abated, and my father tarried in Haran and dwelt there, as there were many flocks in Haran; And my father turned again unto his idolitry: Therefore he continued in Haran

				Now the Lord had said unto Abram <me> get thee out of thy country, and from thy kindr[ed] and from thy fathers home unto a land that I will shew thee. Therefore I left the land of Ur of the chaldees to go into the land of canaan, and I took Lot my bro son and his wife and sarah my wife and also my father follo[we]d me unto the land which we denominated Haran and the famine abated, and my father tarried in Haran and dwelt there as \sout{they} <there> were <many> flock in Haran. and my father turned again unto his idolitry Therefore he continued in Haran but <I> Abram and and Lot my brothers son prayed unto the Lord, and the Lord appeared
				
		% -----
		\subsubsection{Record, 18 November 1835}
		% -----
			\label{EMS-1835-JS-18-NOV-1835}
			% \label{EMS-1830-JS-DAY-3LETTERMONTH-YEAR}
			
			\begin{description}
				\item[Print Sources]
			\end{description}

				\begin{tabular}{p{0.5in}p{1in}p{0.5in}p{1in}}
				JSP	 	& J1:106		& JSR 		& --- \\
				DC	 	& --- 				& HC 		& --- \\
				HDDC 	& ---				& TCHC 		& --- \\
				PJS 	& 1:143--144			& 	 		&  \\
				\end{tabular}
				% HDDC = woodford's DC diss; JSR = Marquardt's JS Revelations; TCHC = Vogel HC
			
			\begin{description}
				\item[Online Access] \href{https://www.josephsmithpapers.org/paper-summary/journal-1835-1836/47}{Here}
				% \item[Analysis] \hyperref[XX]{Here}
			\end{description}
			
			\begin{description}
				\item[Background]
			\end{description}
			
				% It might also be useful to give a two sentence context of the period: e.g., "This speech was given in Nauvoo to a small congregation one month prior to the destruction of the Nauvoo Expositor. These and other tensions may have been on the prophet's mind when he gave the speech."
			
			\begin{description}
				\item[Original Source]
			\end{description}
			
				% brief description of original source material, with reference to JSP or somewhere else for more info
				% Background material: it might be helpful to have a note that says whether a given document was presented as a speech, was written in a journal, etc.
				% Also a comment here about how reliable the source might be, or what shape it is in, if there are large omissions or parts that are hard to understand, and so on.
			
			\begin{description}
				\item[Text]
			\end{description}
			
				...at evening Bishop [Newel K.] Whitney his wife Father and Mother, and \sout{wife} Sister in law, came in and invited me and my wife to go with them \& visit Father Smith \& family my wife was unwell and could not go; however I and my Scribe went, when we got there, we found that some of the young Elders, were about engaging in a debate, upon the subject of miracles, the question was this; was or was it not the design of Christ to establish his gospel by miracles,
				
				After an interesting debate of three hours or more, during which time much talent was displayed, it was decided by the presidents of the debate in the negative; which was a righteous descision I discovered in this debate, much warmth displayed, to[o] much zeal for mastery, to[o] much of that enthusiasm that characterises a lawyer at the bar, who is determined to defend his cause right or wrong. I therefore availed myself of this favorable opportunity, to drop a few words upon this subject by way of advise, that they might improve their minds and cultivate their powers of intellect in a proper manner, that they might not incur the displeasure of heaven, that they should handle sacred things verry sacredly, and with a due deference to the opinions of others and with an eye single to the glory of God
				
		% -----
		\subsubsection{Record, 24 November 1835}
		% -----
			\label{EMS-1835-JS-24-NOV-1835}
			% \label{EMS-1830-JS-DAY-3LETTERMONTH-YEAR}
			
			\begin{description}
				\item[Print Sources]
			\end{description}

				\begin{tabular}{p{0.5in}p{1in}p{0.5in}p{1in}}
				JSP	 	& J1:109--110		& JSR 		& --- \\
				DC	 	& --- 				& HC 		& --- \\
				HDDC 	& ---				& TCHC 		& --- \\
				PJS 	& 1:145--146			& 	 		&  \\
				\end{tabular}
				% HDDC = woodford's DC diss; JSR = Marquardt's JS Revelations; TCHC = Vogel HC
			
			\begin{description}
				\item[Online Access] \href{https://www.josephsmithpapers.org/paper-summary/journal-1835-1836/50}{Here}
				% \item[Analysis] \hyperref[XX]{Here}
			\end{description}
			
			\begin{description}
				\item[Background]
			\end{description}
			
				% It might also be useful to give a two sentence context of the period: e.g., "This speech was given in Nauvoo to a small congregation one month prior to the destruction of the Nauvoo Expositor. These and other tensions may have been on the prophet's mind when he gave the speech."
				
				Useful footnotes in the J1 account; compare to 3 December 1835 (which is not in here, see J1:114); compare to 14 December 1836 (which is in here)
			
			\begin{description}
				\item[Original Source]
			\end{description}
			
				% brief description of original source material, with reference to JSP or somewhere else for more info
				% Background material: it might be helpful to have a note that says whether a given document was presented as a speech, was written in a journal, etc.
				% Also a comment here about how reliable the source might be, or what shape it is in, if there are large omissions or parts that are hard to understand, and so on.
			
			\begin{description}
				\item[Text]
			\end{description}
			
				...I had an invitation, to attend a wedding at Br. Hiram [Hyrum] Smith's in the evening also to solemnize the matrimonial ceremony, <between Newell Knights [Newel Knight] \& Lydia Goldthwaite [Bailey]> I and my wife, went, when we arrived a conciderable company, had collected, the bridegroom \& bride came in, and took their seats, which gave me to understand that they were ready, I requesteded them to arise and join hands, I then remarked that marriage was an institution of h[e]aven institude [instituted] in the garden of Eden, that it was necessary that it should be Solemnized by the authority of the everlasting priesthood, before joining hands however, we attended prayers. I then made the remarks above stated; The ceremony was original <with me> it was in substance as follows, You covenant to be each others companions through life, and discharge the duties of husband \& wife in every respect to which they assented, I then pronounced them husband \& Wife in the name of God and also \sout{pronounced} the blessings that the Lord confered upon adam \& Eve in the garden of Eden; that is to multiply and replenish the earth, with the addition of long life and prosperity; dismissed them and returned home...
				
		% -----
		\subsubsection{Letter to the Elders of the Church, 30 November--1 December 1835}
		% -----
			\label{EMS-1835-JS-30-NOV-1-DEC-1835}
			% \label{EMS-1830-JS-DAY-3LETTERMONTH-YEAR}
			
			\begin{description}
				\item[Print Sources]
			\end{description}

				\begin{tabular}{p{0.5in}p{1in}p{0.5in}p{1in}}
				JSP	 	& D5:89--100		& JSR 		& --- \\
				DC	 	& --- 				& HC 		& 2:264--272 \\
				HDDC 	& ---				& TCHC 		& 2:304--311 \\
				TPJS 	& 94--102			& 	 		&  \\
				PJS 	& 1:149			& 	 		&  \\
				\end{tabular}
				% HDDC = woodford's DC diss; JSR = Marquardt's JS Revelations; TCHC = Vogel HC
			
			\begin{description}
				\item[Online Access] \href{https://www.josephsmithpapers.org/paper-summary/letter-to-the-elders-of-the-church-30-november-1-december-1835/1}{Here}
				% \item[Analysis] \hyperref[XX]{Here}
			\end{description}
			
			\begin{description}
				\item[Background]
			\end{description}
			
				% It might also be useful to give a two sentence context of the period: e.g., "This speech was given in Nauvoo to a small congregation one month prior to the destruction of the Nauvoo Expositor. These and other tensions may have been on the prophet's mind when he gave the speech."
			
			\begin{description}
				\item[Original Source]
			\end{description}
			
				% brief description of original source material, with reference to JSP or somewhere else for more info
				% Background material: it might be helpful to have a note that says whether a given document was presented as a speech, was written in a journal, etc.
				% Also a comment here about how reliable the source might be, or what shape it is in, if there are large omissions or parts that are hard to understand, and so on.
			
			\begin{description}
				\item[Text]
			\end{description}
			
					\hypertarget{EMS-1835-JS-30-NOV-1-DEC-1835-a}%
				\phantom{ }\textit{To the Elders of the Church of the Latter Day Saints}.
					\marginnote{a}

				I have shown unto you, in my last, that there are two Jerusalems spoken of in holy writ, in a manner I think satisfactorily to your minds: At any rate I have given my views upon the subject. I shall now proceed to make some remarks from the sayings of the Savior, recorded in the 13th chapter of his gospel according to St Matthew, which in my mind affords us as clear an understanding, upon the important subject of the gathering, as any thing recorded in the bible. At the time the Savior spoke these beautiful sayings and parables, contained in the chapter above quoted, we find him seated in a ship, on the account of the multitude that pressed upon him to hear his words, and he commenced teaching them by saying: ``Behold a sower went forth to sow, and when he sowed, some seeds fell by the way side, and the fowls came and devoured them up; some fell upon stony places, where they had not much earth, and forthwith they sprang up because they had no deepness of earth, and when the sun was up, they were scorched, and because they had not root they withered away; and some fell among thorns and the thorns sprang up and choked them; but other, fell into good ground and brought forth fruit, some an hundred fold, some sixty fold, some thirty fold: who hath ears to hear let him hear. And the disciples came and said unto him, why speakest thou unto them in parables, (I would remark here, that the ``\textit{them},'' made use of, in this interrogation, is a personal pronoun and refers to the multitude,) he answered and said unto them, (that is the disciples,) it is given unto \textit{you} to know the mysteries of the kingdom of heaven, but unto \textit{them} (that is unbelievers) it is not given, for whosoever hath, to him shall be given, and he shall have more abundance; but whosoever hath not, shall be taken away, even that he hath.''

					\hypertarget{EMS-1835-JS-30-NOV-1-DEC-1835-b}%
				We
					\marginnote{b}%
				understand from this saying, that those who had previously been looking for a Messiah to come, according to the testimony of the Prophets, and were then, at that time, looking for a Messiah, but had not sufficient light on the account of their unbelief, to discern him to be their Savior; and he being the true Messiah, consequently they must be disappointed and lose even all the knowledge, or have taken away from them, all the light, understanding and faith, which they had upon this subject: therefore he that will not receive the greater light, must have taken away from him, all the light which he hath. And if the light which is in you, become darkness, behold how great is that darkness? Therefore says the Savior, speak I unto them in parables, because they, seeing, see not; and hearing, they hear not; neither do they understand: and in them is fulfilled the prophecy of Esaias, which saith: by hearing ye shall hear and shall not understand; and seeing ye shall see and not perceive.

				Now we discover, that the very reasons assigned by this prophet, why they would not receive the Messiah, was, because they did or would not understand; and seeing they did not perceive: for this people's heart is waxed gross; their ears are dull of hearing; their eyes they have closed, lest at any time, they should see with their eyes, and hear with their ears, and understand with their hearts, and should be converted and I should heal them.

					\hypertarget{EMS-1835-JS-30-NOV-1-DEC-1835-c}%
				But
					\marginnote{c}%
				what saith he to his disciples: Blessed are your eyes, for they see, and your ears, for they hear; for verily I say unto you, that many prophets and righteous men have desired to see those things which ye see, and have not seen them; and to hear those things which ye hear, and have not heard them.

				We again make a remark here, for we find that the very principles upon which the disciples were accounted blessed, was because they were permitted to see with their eyes, and hear with their ears, and the condemnation which rested upon the multitude, which received not his saying, was because they were not willing to see with their eyes and hear with their ears; not because they could not and were not privileged to see, and hear, but because their hearts were full of iniquity and abomination: as your fathers did so do ye.--- The prophet foreseeing that they would thus harden their hearts plainly declared it; and herein is the condemnation of the world, that light hath come into the world, and men choose darkness rather than light because their deeds are evil: This is so plainly taught by the Savior, that a wayfaring man need not mistake it.

					\hypertarget{EMS-1835-JS-30-NOV-1-DEC-1835-d}%
				And
					\marginnote{d}%
				again hear ye the parable of the sower: Men are in the habit, when the truth is exhibited by the servants of God, of saying, all is mystery, they are spoken in parables, and, therefore, are not to be understood, it is true they have eyes to see, and see not; but none are so blind as those who will not see: And although the Savior spoke this parable to such characters, yet unto his disciples he expounded it plainly; and we have reason to be truly humble before the God of our fathers, that he hath left these things on record for us, so plain, that, notwithstanding the exertions and combined influence of the priests of Baal, they have not power to blind our eyes and darken our understanding, if we will but open our eyes and read with candor, for a moment. But listen to the explanation of the parable: when any one heareth the word of the kingdom, and understandeth it not, then cometh the wicked one and catcheth away that which was sown in his heart. Now mark the expression; that which was before sown in his heart; this is he which received seed by the way side; men who have no principle of righteousness in themselves, and whose hearts are full of iniquity, and who have no desire for the principles of truth, do not understand the word of truth, when they hear it.--- The devil taketh away the word of truth out of their hearts, because there is no desire for righteousness in them. But he that received the seed into stony places the same is he that heareth the word and, anon, with joy receiveth it, yet hath he not root in himself, but dureth for awhile; for when tribulation or persecution ariseth because of the word, by and by he is offended. He also that received seed among the thorns is he that receiveth the word, and the cares of this world, and the deceitfulness of riches choke the word, and he becometh unfruitful: but he that received seed into the good ground, is he that heareth the word and understandeth it which also beareth fruit and bringeth forth some an hundred fold, some sixty, some thirty. Thus the Savior himself explains unto his disciples the parable, which he put forth and left no mystery or darkness upon the minds of those who firmly believe on his words.

					\hypertarget{EMS-1835-JS-30-NOV-1-DEC-1835-e}%
				We
					\marginnote{e}%
				draw the conclusion then, that the very reason why the multitude, or the world, as they were designated by the Savior, did not receive an explanation upon his parables, was, because of unbelief. To you, he says, (speaking to his disciples) it is given to know the mysteries of the kingdom of God: and why? because of the faith and confidence which they had in him. This parable was spoken to demonstrate the effects that are produced by the preaching of the word; and we believe that it has an allusion directly, to the commencement, or the setting up of the kingdom in that age: therefore, we shall continue to trace his sayings concerning this kingdom from that time forth, even unto the end of the world.

					\hypertarget{EMS-1835-JS-30-NOV-1-DEC-1835-f}%
				Another
					\marginnote{f}%
				parable put he forth unto them, saying, (which parable has an allusion to the setting up of the kingdom, in that age of the world also) the kingdom of Heaven is likened unto a man which sowed good seed in his field, but while men slept an enemy came and sowed tares among the wheat and went his way; but when the blade was sprung up, and brought forth fruit, then appeared the tares also; so the servants of the householder came and said unto him, sir, didst not thou sow good seed in thy field? from whence then hath it tares? He said unto them, an enemy hath done this. The servants said unto him wilt thou then that we go and gather them up; but he said nay, lest while ye gather up the tares, ye root up also the wheat with them.--- Let both grow together until the harvest, and in the time of the harvest, I will say to the reapers, gather ye together first the tares, and bind them in bundles, to burn them; but gather the wheat into my barn.

					\hypertarget{EMS-1835-JS-30-NOV-1-DEC-1835-g}%
				Now
					\marginnote{g}%
				we learn by this parable, not only the setting up of the kingdom in the days of the Savior, which is represented by the good seed, which produced fruit, but also the corruptions of the church, which is represented by the tares, which were sown by the enemy, which his disciples would fain have plucked up, or cleansed the church of, if their views had been favored by the Savior; but he, knowing all things, says not so; as much as to say, your views are not correct, the church is in its infancy, and if you take this rash step, you will destroy the wheat or the church with the tares: therefore it is better to let them grow together until the harvest, or the end of the world, which means the destruction of the wicked; which is not yet fulfilled; as we shall show hereafter, in the Savior's explanation of the parable, which is so plain, that there is no room left for dubiety upon the mind, notwithstanding the cry of the priests, parables, parables! figures, figures! mystery, mystery! all is mystery! but we find no room for doubt here, as the parables were all plainly elucidated.

				And again, another parable put he forth unto them, having an allusion to the kingdom which should be set up, just previous or at the time of harvest, which reads as follows:---The kingdom of heaven is like to a grain of mustard seed, which a man took and sowed in his field, which indeed is the least of all seeds, but when it is grown it is the greatest among herbs, and becometh a tree, so that the birds of the air come and lodge in the branches thereof. Now we can discover plainly, that this figure is given to represent the church as it shall come forth in the last days. Behold the kingdom of heaven is likened unto it. Now what is like unto it?

					\hypertarget{EMS-1835-JS-30-NOV-1-DEC-1835-h}%
				Let
					\marginnote{h}%
				us take the book of Mormon, which a man took and hid in his field; securing it by his faith, to spring up in the last days, or in due time: let us behold it coming forth out of the ground, which is indeed accounted the least of all seeds, but behold it branching forth; yea, even towering, with lofty branches, and God-like majesty, until it becomes the greatest of all herbs: and it is truth, and it has sprouted and come forth out of the earth; and righteousness begins to look down from heaven; and God is sending down his powers, gifts and angels, to lodge in the branches thereof: The kingdom of heaven is like unto a mustard seed. Behold, then, is not this the kingdom of heaven that is raising its head in the last days, in the majesty of its God; even the church of the Latter day saints,---like an impenetrable, immovable rock in the midst of the mighty deep, exposed to storms and tempests of satan, but has, thus far, remained steadfast and is still braving the mountain waves of opposition, which are driven by the tempestuous winds of sinking crafts, have and are still dashing with tremendous foam, across its triumphing brow, urged onward with redoubled fury by the enemy of righteousness, with his pitchfork of lies, as you will see fairly represented in a cut, contained in Mr. [Eber D.] Howe's ``Mormonism Unveiled?''

					\hypertarget{EMS-1835-JS-30-NOV-1-DEC-1835-i}%
				And
					\marginnote{i}%
				we hope that this adversary of truth will continue to stir up the sink of iniquity, that people may the more readily discern between the righteous and wicked. We also would notice one of the modern sons of Sceva, who would fain have made people believe that he could cast out devils, by a certain pamphlet (viz. the ``Millenial Harbinger,'') that went the rounds through our country, who felt so fully authorized to brand Jo Smith, with the appellation of Elymus the sorcerer, and to say with Paul, O full of all subtilty and all mischief, thou child of the devil, thou enemy of all righteousness, wilt thou not cease to pervert the right ways of the Lord! We would reply to this gentleman---Paul we know, and Christ we know, but who are ye? And with the best of feelings, we would say to him, in the language of Paul to those who said they were John's disciples, but had not so much as heard there was a Holy Ghost, to repent and be baptised for the remission of sins by those who have legal authority, and under their hands you shall receive the Holy Ghost, according to the scriptures.

				Then laid they \textit{their} hands on them, and they received the Holy Ghost.---Acts: ch. 8, v. 17.

				And, when Paul had laid \textit{his} hands upon them, the Holy Ghost came on them; and they spake with tongues, and prophesied.---Acts: ch. 19, v. 6.

				Of the doctrine of baptisms, and of laying on of hands, and of resurrection of the dead, and of eternal judgment.---Heb. ch. 6, v.2.

					\hypertarget{EMS-1835-JS-30-NOV-1-DEC-1835-j}%
				How
					\marginnote{j}%
				then shall they call on him in whom they have not believed? and how shall they believe in him of whom they have not heard? and how shall they hear without a preacher? And how shall they preach except they be sent? as it is written, How beautiful are the feet of them that preach the gospel of peace, and bring glad tidings of good things!---Rom. ch. 10, v. 14--15.

				But if this man will not take our admonition, but will persist in his wicked course, we hope that he will continue trying to cast out devils, that we may have the clearer proof that the kingdom of satan is divided against itself, and consequently cannot stand: for a kingdom divided against itself, speedily hath an end. If we were disposed to take this gentleman upon his own ground and justly heap upon him that which he so readily and unjustly heaps upon others, we might go farther; we might say that he has wickedly and maliciously lied about, vilified and traduced the characters of innocent men. We might invite the gentleman to a public investigation of these matters; yea, and we do challenge him to an investigation upon any or all principles wherein he feels opposed to us, in public or in private.

					\hypertarget{EMS-1835-JS-30-NOV-1-DEC-1835-k}%
				We
					\marginnote{k}%
				might farther say that, we could introduce him to ``Mormonism Unveiled.'' Also to the right honorable Doct. P. Hurlburt [Doctor Philastus Hurlbut], who is the legitimate author of the same, who is not so much a doctor of physic, as of falsehood, or by name We could also give him an introduction to the reverend Mr. Howe, the illegitimate author of ``Mormonism Unveiled,'' in order to give currency to the publication, as Mr. Hurlburt, about this time, was bound over to court, for threatening life. He is also an associate of the celebrated Mr. [Orris] Clapp, who has of late immortalised his name by swearing that he would not believe a Mormon under oath; and by his polite introduction to said Hurlburt's wife, which cost him (as we have been informed) a round sum. Also his son Mathew [Matthew Clapp] testified that, the book of Mormon had been proved false an hundred times, by How's book: and also, that he would not believe a Mormon under oath. And also we could mention the reverend Mr. [Adamson] Bentley, who, we believe, has been actively engaged in injuring the character of his brother-in-law, viz: Elder S[idney] Rigdon.

				Now, the above statements are according to our best information: and we believe them to be true; and this is as fair a sample of the doctrine of Campbellism, as we ask, taking the statements of these gentlemen, and judging them by their fruits. And we might add many more to the black catalogue; even the ringleaders, not of the Nazarenes, for how can any good thing come out of Nazareth, but of the far-famed Mentor mob: all sons and legitimate heirs to the same spirit of Alexander Campbell, and ``Mormonism Unveiled,'' according to the representation in the cut spoken of above.

					\hypertarget{EMS-1835-JS-30-NOV-1-DEC-1835-l}%
				The
					\marginnote{l}%
				above cloud of darkness has long been beating with mountain waves upon the immovable rock of the church of the Latter Day Saints, and notwithstanding all this, the mustard seed is still towering its lofty branches, higher and higher, and extending itself wider and wider, and the charriot wheels of the kingdom are still rolling on, impelled by the mighty arm of Jehovah; and in spite of all opposition will still roll on until his words are all fulfilled.

				Our readers will excuse us for deviating from the subject, when they take into consideration the abuses, that have been heaped upon us heretofore, which we have tamely submitted to, until forbearance is no longer required at our hands, having frequently turned both the right and left cheek, we believe it our duty now to stand up in our own defence. With these remarks we shall proceed with the subject of the gathering.

				And another parable spake he unto them: The kingdom of heaven is like unto leaven which a woman took and hid in three measures of meal, until the whole was leavened. It may be understood that the church of the Latter Day Saints, has taken its rise from a little leaven that was put into three witnesses. Behold, how much this is like the parable: it is fast leavening the lump, and will soon leaven the whole. But let us pass on.

					\hypertarget{EMS-1835-JS-30-NOV-1-DEC-1835-m}%
				All
					\marginnote{m}%
				these things spake Jesus unto the multitudes, in parables, and without a parable spake he not unto them, that it might be fulfilled which was spoken by the prophet, saying: I will open my mouth in parables: I will utter things which have been kept secret from the foundation of the world: Then Jesus sent the multitude away and went into the house, and his disciples came unto him, saying, declare unto us the parable of the tares of the field. He answered and said unto them, he that soweth the good seed is the son of man; the field is the world; the good seed are the children of the kingdom, but the tares are the children of the wicked one. Now let our readers mark the expression, the field is the world; the tares are the children of the wicked one: the enemy that sowed them is the devil;
					\hypertarget{EMS-1835-JS-30-NOV-1-DEC-1835-n}%
				the
					\marginnote{n}%
				harvest is the end of the world. Let them carefully mark this expression also, \textit{the end of the world}, and the reapers are the angels. Now men cannot have any possible grounds to say that this is figurative, or that it does not mean what it says; for he is now explaining what he had previously spoken in parables; and according to this language, the end of the world is the destruction of the wicked; the harvest and the end of the world have an allusion directly to the human family in the last days, instead of the earth, as many have imagined, and that which shall precede the coming of the Son of man, and the restitution of all things spoken of by the mouth of all the holy prophets since the world began; and the angels are to have something to do in this great work, for they are the reapers: as therefore the tares are gathered and burned in the fire, so shall it be in the end of this world; that is, as the servants of God go forth warning the nations, both priests and people, and as they harden their hearts and reject the light of the truth, these first being delivered over unto the buffetings of satan, and the law and the testimony being closed up, as it was with the Jews, they are left in darkness, and delivered over unto the day of burning: thus being bound up by their creeds and their bands made strong by their \textit{priests}, are prepared for the fulfilment of the saying of the Savior: the Son of man shall send forth his angels, and gather out of his kingdom all things that offend, and them which do iniquity, and shall cast them into a furnace of fire and there shall be wailing and gnashing of teeth.

					\hypertarget{EMS-1835-JS-30-NOV-1-DEC-1835-o}%
				We
					\marginnote{o}%
				understand, that the work of the gathering together of the wheat into barns, or garners, is to take place while the tares are being bound over, and preparing for the day of burning: that after the day of burnings, the righteous shall shine forth like the sun, in the kingdom of their Father: who hath ears to hear let him hear.

				But to illustrate more clearly upon this gathering, we have another parable. Again the kingdom of heaven is like a treasure hid in a field, the which when a man hath found, he hideth and for joy thereof, goeth and selleth all that he hath and buyeth that field: for the work after this pattern, see the church of the Latter Day Saints, selling all that they have and gathering themselves together unto a place that they may purchase for an inheritance, that they may be together and bear each other's afflictions in the day of calamity.

				Again the kingdom of heaven is like unto a merchant man seeking goodly pearls, who when he had found one pearl of great price, went and sold all that he had, and bought it. For the work of this example, see men travelling to find places for Zion, and her stakes or remnants, who when they find the place for Zion, or the pearl of great price; straitway sell all that they have and buy it.

					\hypertarget{EMS-1835-JS-30-NOV-1-DEC-1835-p}%
				Again
					\marginnote{p}%
				the kingdom of heaven is like unto a net that was cast into the sea, and gathered of every kind, which when it was full they drew to shore, and sat down and gathered the good into vessels, and cast the bad away.--- For the work of this pattern, behold the seed of Joseph, spreading forth the gospel net, upon the face of the earth, gathering of every kind, that the good may be saved in vessels prepared for that purpose, and the angels will take care of the bad: so shall it be at the end of the world, the angels shall come forth, and sever the wicked from among the just, and cast them into the furnace of fire, and there shall be wailing and gnashing of teeth.

				Jesus saith unto them, have you understood all these things? they say unto him yea Lord: and we say yea Lord, and well might they say yea Lord, for these things are so plain and so glorious, that every Saint in the last days must respond with a hearty \textit{amen} to them.

					\hypertarget{EMS-1835-JS-30-NOV-1-DEC-1835-q}%
				Then
					\marginnote{q}%
				said he unto them, therefore every scribe which is instructed into the kingdom of heaven, is like unto a man that is an house holder; which bringeth forth out of his treasure things that are new and old.

				For the work of this example, see the book of Mormon, coming forth out of the treasure of the heart; also the covenants given to the Latter Day Saints: also the translation of the bible: thus bringing forth out of the heart, things new and old: thus answering to three measures of meal, undergoing the purifying touch by a revelation of Jesus Christ, and the ministering of angels, who have already commenced this work in the last days, which will answer to the leaven which leavened the whole lump. Amen.

				So I close but shall continue the subject in another number.

				In the bonds of the new and everlasting covenant.

				\begin{tightright}
				JOSEPH SMITH, jr.
				\end{tightright}

				To J[ohn] \textsc{Whitmer} Esq.
				
		% -----
		\subsubsection{Record, 2 December 1835}
		% -----
			\label{EMS-1835-JS-2-DEC-1835}
			% \label{EMS-1830-JS-DAY-3LETTERMONTH-YEAR}
			
			\begin{description}
				\item[Print Sources]
			\end{description}

				\begin{tabular}{p{0.5in}p{1in}p{0.5in}p{1in}}
				JSP	 	& J1:113--114		& JSR 		& --- \\
				DC	 	& --- 				& HC 		& --- \\
				HDDC 	& ---				& TCHC 		& --- \\
				PJS 	& 1:149--150			& 	 		&  \\
				\end{tabular}
				% HDDC = woodford's DC diss; JSR = Marquardt's JS Revelations; TCHC = Vogel HC
			
			\begin{description}
				\item[Online Access] \href{https://www.josephsmithpapers.org/paper-summary/journal-1835-1836/54}{Here}
				% \item[Analysis] \hyperref[XX]{Here}
			\end{description}
			
			\begin{description}
				\item[Background]
			\end{description}
			
				% It might also be useful to give a two sentence context of the period: e.g., "This speech was given in Nauvoo to a small congregation one month prior to the destruction of the Nauvoo Expositor. These and other tensions may have been on the prophet's mind when he gave the speech."
			
			\begin{description}
				\item[Original Source]
			\end{description}
			
				% brief description of original source material, with reference to JSP or somewhere else for more info
				% Background material: it might be helpful to have a note that says whether a given document was presented as a speech, was written in a journal, etc.
				% Also a comment here about how reliable the source might be, or what shape it is in, if there are large omissions or parts that are hard to understand, and so on.
			
			\begin{description}
				\item[Text]
			\end{description}
			
				...we overtook a team with two men on the sleigh. I politely asked them to let me pass, they granted my request, and as we passed them, they bawled out, do you get any revelation lately, with an adition of blackguard that I did not understand, this is a fair sample of the character of Mentor Street inhabitants, who are ready to abuse and scandalize, men who never laid a straw in their way, and infact those whos faces they never saw, and cannot, bring an acusation, against, either of a temporal or spirtual nature; except our firm belief in the fulness of the gospel and I was led to marvle \sout{that} \sout{God} at the long suffering and condescention of our heavenly Father, in permitting, these ungodly wretches, to possess, this goodly land, which is \sout{the} indeed as beautifully situated, and its soil as fertile, as any in this region of country, and its inhabitance, \sout{as} wealthy even blessed, above measure, in temporal things, and fain, would God bless, them with, \sout{with} spiritual blessings, even eternal life, were it not for their evil hearts of unbelief, and we are led to \sout{cry} \sout{in} \sout{our} \sout{hearts} mingle our prayers with those saints that have suffered the like treatment before us, whose souls are under the altar crying to the Lord for vengance upon those that dwell upon the earth and we rejoice that the time is at hand when, the wicked who will not repent will be swept <from the earth> with the besom of destruction and the earth become an inheritance for the poor and the meek...
				
		% -----
		\subsubsection{Record, 9 December 1835}
		% -----
			\label{EMS-1835-JS-9-DEC-1835}
			% \label{EMS-1830-JS-DAY-3LETTERMONTH-YEAR}
			
			\begin{description}
				\item[Print Sources]
			\end{description}

				\begin{tabular}{p{0.5in}p{1in}p{0.5in}p{1in}}
				JSP	 	& J1:118		& JSR 		& --- \\
				DC	 	& --- 				& HC 		& --- \\
				HDDC 	& ---				& TCHC 		& --- \\
				PJS 	& 1:155--156			& 	 		&  \\
				\end{tabular}
				% HDDC = woodford's DC diss; JSR = Marquardt's JS Revelations; TCHC = Vogel HC
			
			\begin{description}
				\item[Online Access] \href{https://www.josephsmithpapers.org/paper-summary/journal-1835-1836/60}{Here}
				% \item[Analysis] \hyperref[XX]{Here}
			\end{description}
			
			\begin{description}
				\item[Background]
			\end{description}
			
				% It might also be useful to give a two sentence context of the period: e.g., "This speech was given in Nauvoo to a small congregation one month prior to the destruction of the Nauvoo Expositor. These and other tensions may have been on the prophet's mind when he gave the speech."
			
			\begin{description}
				\item[Original Source]
			\end{description}
			
				% brief description of original source material, with reference to JSP or somewhere else for more info
				% Background material: it might be helpful to have a note that says whether a given document was presented as a speech, was written in a journal, etc.
				% Also a comment here about how reliable the source might be, or what shape it is in, if there are large omissions or parts that are hard to understand, and so on.
			
			\begin{description}
				\item[Text]
			\end{description}
			
				...My heart swells with gratitude inexpressible when I realize the great condescention of my heavenly Fathers, in opening the hearts of these, my beloved brethren to administer so liberally, to my wants and I ask God in the name of Jesus Christ, to multiply, blessings, without number upon their heads, and bless me with much wisdom and understanding, and dispose of me, to the best advantage, for my brethren, and the advancement, of thy cause and Kingdom, and whether my days are many or few whether in life or in death I say in my heart, O Lord let me enjoy the society of such brethren...
				
		% -----
		\subsubsection{Record, 10 December 1835}
		% -----
			\label{EMS-1835-JS-10-DEC-1835}
			% \label{EMS-1830-JS-DAY-3LETTERMONTH-YEAR}
			
			\begin{description}
				\item[Print Sources]
			\end{description}

				\begin{tabular}{p{0.5in}p{1in}p{0.5in}p{1in}}
				JSP	 	& J1:118--119		& JSR 		& --- \\
				DC	 	& --- 				& HC 		& --- \\
				HDDC 	& ---				& TCHC 		& --- \\
				PJS 	& 1:156--157			& 	 		&  \\
				\end{tabular}
				% HDDC = woodford's DC diss; JSR = Marquardt's JS Revelations; TCHC = Vogel HC
			
			\begin{description}
				\item[Online Access] \href{https://www.josephsmithpapers.org/paper-summary/journal-1835-1836/61}{Here}
				% \item[Analysis] \hyperref[XX]{Here}
			\end{description}
			
			\begin{description}
				\item[Background]
			\end{description}
			
				% It might also be useful to give a two sentence context of the period: e.g., "This speech was given in Nauvoo to a small congregation one month prior to the destruction of the Nauvoo Expositor. These and other tensions may have been on the prophet's mind when he gave the speech."
			
			\begin{description}
				\item[Original Source]
			\end{description}
			
				% brief description of original source material, with reference to JSP or somewhere else for more info
				% Background material: it might be helpful to have a note that says whether a given document was presented as a speech, was written in a journal, etc.
				% Also a comment here about how reliable the source might be, or what shape it is in, if there are large omissions or parts that are hard to understand, and so on.
			
			\begin{description}
				\item[Text]
			\end{description}
				
					\hypertarget{EMS-1835-JS-10-DEC-1835-a}%
				...a
					\marginnote{a}%
				beautiful morning, indeed, and fine sleighing, this day my brethren, meet according, to previous arangement, to chop and haul wood for me, and they have been verry industrious, and I think they have supplyed me, with my winters wood, for which I am sincerely grateful to each and every, one of them, for this expression of their, goodness towards me
				
				And in the name of Jesus Christ I envoke the rich benediction of heavn to rest upon them, \sout{even} \sout{all} and their families, and I ask my heavenly Father to preserve their health's and those of their wives and children, that they may have strength. of body to perform, their, labours, in their several ocupations in life, and the use and activity of their limbs, also powers of intellect and understanding hearts, that they may treasure up, wisdom, \sout{and} understanding, \sout{until} and inteligence, above measure, and be preserved from plagues pestilence, and famine, and from the power of the adversary, and the hands of evil designing, men and have power over all their enemys; and the way be prepared before them, that they may journey to the land of Zion and be established, on their inheritances, to enjoy undisturbe[d], peace and happiness for ever, and ultimately, to be crowned with everlasting life in the celestial kingdom of God, which blessings I ask in the name of Jesus of Nazareth. Amen
				
					\hypertarget{EMS-1835-JS-10-DEC-1835-b}%
				I
					\marginnote{b}%
				would remember Elder Leonard Rich who was the first one that proposed to the brethren, to assist me, in, obtaining wood for the use of my family, for which I pray my heavenly Father, to bless <him> with all the blessings, named above, and I shall ever remember him with much gratitude, for this testimony, of benevolence and respect, and thank the great I am, for puting into his heart, to do me this kindness, and I say in my heart, I will trust in thy goodness, and mercy, forever, for thy wisdom and benevolence <O Lord> is unbounded and beyond the comprehension; of men and all of thy ways cannot be found out...
			
		% -----
		\subsubsection{Record, 12 December 1835}
		% -----
			\label{EMS-1835-JS-12-DEC-1835}
			% \label{EMS-1830-JS-DAY-3LETTERMONTH-YEAR}
			
			\begin{description}
				\item[Print Sources]
			\end{description}

				\begin{tabular}{p{0.5in}p{1in}p{0.5in}p{1in}}
				JSP	 	& J1:120--121		& JSR 		& --- \\
				DC	 	& --- 				& HC 		& --- \\
				HDDC 	& ---				& TCHC 		& --- \\
				PJS 	& 1:158--159			& 	 		&  \\
				\end{tabular}
				% HDDC = woodford's DC diss; JSR = Marquardt's JS Revelations; TCHC = Vogel HC
			
			\begin{description}
				\item[Online Access] \href{https://www.josephsmithpapers.org/paper-summary/journal-1835-1836/64}{Here}
				% \item[Analysis] \hyperref[XX]{Here}
			\end{description}
			
			\begin{description}
				\item[Background]
			\end{description}
			
				% It might also be useful to give a two sentence context of the period: e.g., "This speech was given in Nauvoo to a small congregation one month prior to the destruction of the Nauvoo Expositor. These and other tensions may have been on the prophet's mind when he gave the speech."
				
				See the J1 entry for 16 December.
			
			\begin{description}
				\item[Original Source]
			\end{description}
			
				% brief description of original source material, with reference to JSP or somewhere else for more info
				% Background material: it might be helpful to have a note that says whether a given document was presented as a speech, was written in a journal, etc.
				% Also a comment here about how reliable the source might be, or what shape it is in, if there are large omissions or parts that are hard to understand, and so on.
			
			\begin{description}
				\item[Text]
			\end{description}
			
					\hypertarget{EMS-1835-JS-12-DEC-1835-a}%
				...at
					\marginnote{a}%
				about 12 oclock a number of young person[s] called to see the \sout{records} Egyptian records I requested my Scribe to exibit them, he did so, one of the young ladies, who had been examining them, was asked if they had the appearance of Antiquity, she observed with an air of contempt that they did not, on hearing this I was surprised at the ignorance she displayed, and I observed to her that she was an anomaly in creation for all the wise and learned that had ever examined them, without hesitation pronounced them antient, I further remarked that, it was downright wickedness ignorance bigotry and superstition that caused her to make the remark, and that I would put it on record, and I have done so because it is a fair sample of the prevailing spirit of the times showing that the victims of priestcraft and superstition, would not believe though one should rise from the dead.
				
					\hypertarget{EMS-1835-JS-12-DEC-1835-b}%
				At
					\marginnote{b}%
				evening attended a debate, at Br. W\supersu{m}\supers{.} Smiths, the question proposed to debate upon was, as follows.--- was it necessary for God to reveal himself to man, in order for their happiness,--- I was on the affirmative and the last One to speak on that side of the question,--- but while listning, with interest to the, ingenuity displayed, on both sides of the qu[e]stion, I was called, away to visit, Sister Angeline Work[s], who was suposed to be dangerously sick...
				
		% -----
		\subsubsection{Record, 13 December 1835}
		% -----
			\label{EMS-1835-JS-13-DEC-1835}
			% \label{EMS-1830-JS-DAY-3LETTERMONTH-YEAR}
			
			\begin{description}
				\item[Print Sources]
			\end{description}

				\begin{tabular}{p{0.5in}p{1in}p{0.5in}p{1in}}
				JSP	 	& J1:121--122		& JSR 		& --- \\
				DC	 	& --- 				& HC 		& --- \\
				HDDC 	& ---				& TCHC 		& --- \\
				PJS 	& 1:160			& 	 		&  \\
				\end{tabular}
				% HDDC = woodford's DC diss; JSR = Marquardt's JS Revelations; TCHC = Vogel HC
			
			\begin{description}
				\item[Online Access] \href{https://www.josephsmithpapers.org/paper-summary/journal-1835-1836/66}{Here}
				% \item[Analysis] \hyperref[XX]{Here}
			\end{description}
			
			\begin{description}
				\item[Background]
			\end{description}
			
				% It might also be useful to give a two sentence context of the period: e.g., "This speech was given in Nauvoo to a small congregation one month prior to the destruction of the Nauvoo Expositor. These and other tensions may have been on the prophet's mind when he gave the speech."
			
			\begin{description}
				\item[Original Source]
			\end{description}
			
				% brief description of original source material, with reference to JSP or somewhere else for more info
				% Background material: it might be helpful to have a note that says whether a given document was presented as a speech, was written in a journal, etc.
				% Also a comment here about how reliable the source might be, or what shape it is in, if there are large omissions or parts that are hard to understand, and so on.
			
			\begin{description}
				\item[Text]
			\end{description}
			
				...We then rode to Mr. [Isaac] McWithy's a distance of about 3, miles from Town, where I had been Solicited, to attend, and solemnize, the matrimonial covenant betwe[e]n Mr. E[dwin] Webb \& Miss E. A. McWithy [Eliza Ann McWethy], the parents and many of the connections of both parties were present, with a large and respectable company of friends, who were invited as guests; and after making the necessary arangements the company come to order, and the Groom \& bride, with the attendants politely came forward, and took their seats, and having been requested, to make some preliminary remarks upon the subject of matrimony, touching the design of the All Mighty in this institution, also the duties, of husbands \& wives towards eac[h]other, and after opening our interview, with singing and prayer, I delivered a lecture of about 40, minuits in length, during this time all seemed to be interested, except\sout{one} one or two individuals, who manifested, a spirit of grovling contempt, which I was constrained to reprove and rebuke sharply, after I had \sout{been} closed my remarks, I sealed the matrimonial ceremony in the name of God, and pronounced the blessings of heaven. upon the heads of the young married couple we then closed by returning thanks...
				
		% -----
		\subsubsection{Record, 14 December 1835}
		% -----
			\label{EMS-1835-JS-14-DEC-1835}
			% \label{EMS-1830-JS-DAY-3LETTERMONTH-YEAR}
			
			\begin{description}
				\item[Print Sources]
			\end{description}

				\begin{tabular}{p{0.5in}p{1in}p{0.5in}p{1in}}
				JSP	 	& J1:122		& JSR 		& --- \\
				DC	 	& --- 				& HC 		& --- \\
				HDDC 	& ---				& TCHC 		& --- \\
				PJS 	& 1:161			& 	 		&  \\
				\end{tabular}
				% HDDC = woodford's DC diss; JSR = Marquardt's JS Revelations; TCHC = Vogel HC
			
			\begin{description}
				\item[Online Access] \href{https://www.josephsmithpapers.org/paper-summary/journal-1835-1836/67}{Here}
				% \item[Analysis] \hyperref[XX]{Here}
			\end{description}
			
			\begin{description}
				\item[Background]
			\end{description}
			
				% It might also be useful to give a two sentence context of the period: e.g., "This speech was given in Nauvoo to a small congregation one month prior to the destruction of the Nauvoo Expositor. These and other tensions may have been on the prophet's mind when he gave the speech."
			
			\begin{description}
				\item[Original Source]
			\end{description}
			
				% brief description of original source material, with reference to JSP or somewhere else for more info
				% Background material: it might be helpful to have a note that says whether a given document was presented as a speech, was written in a journal, etc.
				% Also a comment here about how reliable the source might be, or what shape it is in, if there are large omissions or parts that are hard to understand, and so on.
			
			\begin{description}
				\item[Text]
			\end{description}
			
				...To day Samuel Branum [Brannan] came to my house, much afflicted with a swelling on his left arm, which was occasioned by a bruise on his elbow, we had been called to pray for him and anoint him with oil, but his faith was not sufficient to effect, a cure, and my wife prepared a poultice of herbs and applyed to it and he tarryed with me over night...
				
		% -----
		\subsubsection{Record, 15 December 1835}
		% -----
			\label{EMS-1835-JS-15-DEC-1835}
			% \label{EMS-1830-JS-DAY-3LETTERMONTH-YEAR}
			
			\begin{description}
				\item[Print Sources]
			\end{description}

				\begin{tabular}{p{0.5in}p{1in}p{0.5in}p{1in}}
				JSP	 	& J1:123		& JSR 		& --- \\
				DC	 	& --- 				& HC 		& --- \\
				HDDC 	& ---				& TCHC 		& --- \\
				PJS 	& 1:161--162			& 	 		&  \\
				\end{tabular}
				% HDDC = woodford's DC diss; JSR = Marquardt's JS Revelations; TCHC = Vogel HC
			
			\begin{description}
				\item[Online Access] \href{https://www.josephsmithpapers.org/paper-summary/journal-1835-1836/68}{Here}
				% \item[Analysis] \hyperref[XX]{Here}
			\end{description}
			
			\begin{description}
				\item[Background]
			\end{description}
			
				% It might also be useful to give a two sentence context of the period: e.g., "This speech was given in Nauvoo to a small congregation one month prior to the destruction of the Nauvoo Expositor. These and other tensions may have been on the prophet's mind when he gave the speech."
			
			\begin{description}
				\item[Original Source]
			\end{description}
			
				% brief description of original source material, with reference to JSP or somewhere else for more info
				% Background material: it might be helpful to have a note that says whether a given document was presented as a speech, was written in a journal, etc.
				% Also a comment here about how reliable the source might be, or what shape it is in, if there are large omissions or parts that are hard to understand, and so on.
			
			\begin{description}
				\item[Text]
			\end{description}
			
					\hypertarget{EMS-1835-JS-15-DEC-1835-a}%
				...This
					\marginnote{a}%
				afternon Elder Orson Hyde, handed me a Letter, the purport of which is that he is dissatisfyed with the committee, in their dealings, with him in temporal affairs, that is that they do not deal as liberally \sout{in} <with> him as they do with Elder William Smith, also requested me to reconcile the revelation, given to the 12, since their return from the East,
				
				That unless these things and others named in the letter, could be reconciled to his mind his honour would not stand united with them,--- this I believe is the amount of the contents of the letter although much was written, my feelings on this occasion, were much laserated, knowing that I had dealt in righteousness with him in all things and endeavoured to promote his happiness and well being, as much as lay in my power,
					\hypertarget{EMS-1835-JS-15-DEC-1835-b}%
				and
					\marginnote{b}%
				I feel that these reflections are ungrateful and founded in jealousy and that the adversary is striving with all his subtle devises and influence to destroy him by causing a division amon[g] the twelve that God has chosen to open the gospel kingdom in all nations, but I pray my Heavenly Father in the name of Jesus of Nazeareth that he may be delivered from the power of the destroyer, \sout{and} that his faith fail not in this hour of temptation, and prepare him and all the Elders to receive an endument, in thy house, even according to \sout{thy} <thine> own order from time to time as thou seeest them worthy to be called into thy Solemn Assembly. 
				
		% -----
		\subsubsection{Record, 16 December 1835}
		% -----
			\label{EMS-1835-JS-16-DEC-1835}
			% \label{EMS-1830-JS-DAY-3LETTERMONTH-YEAR}
			
			\begin{description}
				\item[Print Sources]
			\end{description}

				\begin{tabular}{p{0.5in}p{1in}p{0.5in}p{1in}}
				JSP	 	& J1:123--124		& JSR 		& --- \\
				DC	 	& --- 				& HC 		& --- \\
				HDDC 	& ---				& TCHC 		& --- \\
				PJS 	& 1:163			& 	 		&  \\
				\end{tabular}
				% HDDC = woodford's DC diss; JSR = Marquardt's JS Revelations; TCHC = Vogel HC
			
			\begin{description}
				\item[Online Access] \href{https://www.josephsmithpapers.org/paper-summary/journal-1835-1836/70}{Here}
				% \item[Analysis] \hyperref[XX]{Here}
			\end{description}
			
			\begin{description}
				\item[Background]
			\end{description}
			
				% It might also be useful to give a two sentence context of the period: e.g., "This speech was given in Nauvoo to a small congregation one month prior to the destruction of the Nauvoo Expositor. These and other tensions may have been on the prophet's mind when he gave the speech."
			
			\begin{description}
				\item[Original Source]
			\end{description}
			
				% brief description of original source material, with reference to JSP or somewhere else for more info
				% Background material: it might be helpful to have a note that says whether a given document was presented as a speech, was written in a journal, etc.
				% Also a comment here about how reliable the source might be, or what shape it is in, if there are large omissions or parts that are hard to understand, and so on.
			
			\begin{description}
				\item[Text]
			\end{description}
			
				...Returned home Elder McLellen [William E. McLellin] Elder B[righam] Young and Elder J[ared] Carter called and paid me a visit, with which I was much gratified, I exibited and explaind the Egyptian Records to them, and explained many things to them concerning the dealings of God with the ancient<s> and the formation of the planetary System, they seemed much pleased with the interview...

		% -----
		\subsubsection{Letter to William Smith, circa 18 December 1835}
		% -----
			\label{EMS-1835-JS-CIR-18-DEC-1835}
			% \label{EMS-1830-JS-DAY-3LETTERMONTH-YEAR}
			
			\begin{description}
				\item[Print Sources]
			\end{description}

				\begin{tabular}{p{0.5in}p{1in}p{0.5in}p{1in}}
				JSP	 	& D5:115--121; J1:131--134		& JSR 		& --- \\
				DC	 	& --- 				& HC 		& 2:340--343 \\
				HDDC 	& ---				& TCHC 		& 2:334--337 \\
				PJS 	& 1:170--175			& 	 		&  \\ 
				\end{tabular}
				% HDDC = woodford's DC diss; JSR = Marquardt's JS Revelations; TCHC = Vogel HC
			
			\begin{description}
				\item[Online Access] \href{https://www.josephsmithpapers.org/paper-summary/letter-to-william-smith-circa-18-december-1835/1}{Here}
				% \item[Analysis] \hyperref[XX]{Here}
			\end{description}
			
			\begin{description}
				\item[Background]
			\end{description}
			
				% It might also be useful to give a two sentence context of the period: e.g., "This speech was given in Nauvoo to a small congregation one month prior to the destruction of the Nauvoo Expositor. These and other tensions may have been on the prophet's mind when he gave the speech."
			
			\begin{description}
				\item[Original Source]
			\end{description}
			
				% brief description of original source material, with reference to JSP or somewhere else for more info
				% Background material: it might be helpful to have a note that says whether a given document was presented as a speech, was written in a journal, etc.
				% Also a comment here about how reliable the source might be, or what shape it is in, if there are large omissions or parts that are hard to understand, and so on.
			
			\begin{description}
				\item[Text]
			\end{description}
			
				\begin{tightcenter}
						\hypertarget{EMS-1835-JS-CIR-18-DEC-1835-a}%
					Kirtland Friday Dec \sout{17\supers{th}} <18\supers{th}> 1835
						\marginnote{a}
				\end{tightcenter}

				Answer to the foregoing Letter from Br. William Smith a Copy

				Br. William.

				having received your letter I now procede to answer it, and shall first procede, to give a brief naration of my feelings and motives, since the night I first came to the knowledge, of your having a debating school, which was at the time I happened, in with, Bishop [Newel K.] Whitney his Father and Mother \&c--- which was the first that I knew any thing about it, and from that time I took an interest in them, and was delighted with it, and formed a determination, to attend the school for the purpose of obtaining information, and with the idea of imparting the same, through the assistance of the spirit of the Lord, if by any means I should have faith to do so; and with this intent, I went to the school on <last> Wedensday night, not with the idea of braking up the school, neither did it enter into my heart, that there was any wrangling or jealousy's in your heart, against me;

					\hypertarget{EMS-1835-JS-CIR-18-DEC-1835-b}%
				Notwithstanding
					\marginnote{b}%
				previous to my leaving home there were feelings of solemnity, rolling across my breast, which were unaccountable to me, and also these feelings continued by spells to depress my \sout{feelings} <spirit> and seemed to manifest that all was not right, even after the \sout{debate} school commenced, and during the debate, yet I strove to believe that all would work together for good; I was pleased with the power of the arguments, that were aduced, and did not feel to cast any reflections, upon any one that had spoken; but I felt that it was \sout{my} <the> duty of old men that set as presidents to be as grave, at least as young men, and that it was our duty to smile at solid arguments, and sound reasoning, and be impreesed, with solemnity, which should be manifest in our countanance, when folly and that which militates against truth and righteousness, rears its head

				Therefore in the spirit of my calling and in view of the authority of the priesthood that has been confered upon me, it would be my duty to reprove whatever I esteemed to be wrong fondly hoping in my heart that all parties, would concider it right, and therefore humble themselves, that satan might not take the advantage of us, and hinder the progress of our School.

					\hypertarget{EMS-1835-JS-CIR-18-DEC-1835-c}%
				Now
					\marginnote{c}%
				Br. William I want you should bear with me, notwithstanding my plainness---

				I would say to you that my feelings, were grieved at the interuption you made upon Elder McLellen [William E. McLellin], I thought, you should have concidered your relation, with him, in your Apostle ship, and not manifest any division of sentiment, between you, and him, for a surrounding multitude to take the advantage of you:--- Therefore by way of entreaty, on the account of the anxiety I had for your influence and wellfare, I said, unto you, do not have any feelings, or something to that amount, why I am thus particular, is that if You, have misconstrued, my feelings, toward you, you may be corrected.---

				But to procede--- after the school was closed Br. Hyrum [Smith], requested, the privilege, of speaking, you objected, however you said if he would not abuse the school, he might speak, and that you would not allow any man to abuse the school in your house,---

				Now you had no reason to suspect that Hyrum, would abuse the school, therefore my feelings were mortifyed, at those unnecessa[r]y observations, I undertook to reason, with you but you manifisted, an inconciderate and stubourn spirit, I then dispared, of benefiting you, on the account of the spirit you manifested, which drew from, me the expression that you was as ugly as the Devil.

					\hypertarget{EMS-1835-JS-CIR-18-DEC-1835-d}%
				Father
					\marginnote{d}%
				then commanded silence and I formed a determination, to obey his mandate, and was about to leave the house, with the impression, that You was under the influence of a wicked spirit, you replyed that you, would say what you pleased in your own house, Father replyed, say what you please, but let the rest hold their, toungs, then a reflection, rushed through my mind, of the, anxiety, and care I \sout{had} <hav> had for you and your family, in doing what I did, in finishing your house and providin flour for your family \&c and also father had possession in the house, as well, as your self; and when at any time have I transgressed, the commandments of my father? or sold my birthright, that I should not have the privilege of speaking in my fathers house, or in other words in my fathers family, or in your house, (for so we will call it, and so it shall be,) that I should not have the privilege, of reproving a younger brother, therefore I said I will speak, for I built the house, and it is as much mine as yours, or something, to that effect,
					\hypertarget{EMS-1835-JS-CIR-18-DEC-1835-e}%
				(I
					\marginnote{e}%
				should have said that. I helped finish the house,) I said it merely to show that it could not be, the right spirit, that would rise up for trifling matters, and undertake to put me to silence, I saw that your indignation was kindled against me, and you made towards me, I was not then to be moved, and I thought, to pull off my loose coat, least it should tangle me, and you be left to hurt me, but not with the intention, of hurting You, but you was to[o] soon for me, and having once fallen into the hands of a mob, and \sout{now} been wounded in my side, and now into the hands of a brother, my side gave way, and after having been rescued, from your grasp, I left your house, with, feelings that were indiscribale [indescribable], the scenery had changed, and all those expectations, that I had cherished, when going to your house, of brotherly kindness, charity forbearance and natural, affection, that in duty binds us not to make eachothers offenders for a word. \sout{but}

				But alass! abuse, anger, malice, hatred, and rage <with a lame side> with marks, of violence <heaped> upon \sout{my} \sout{body} me by a brother, were the reflections of my disapointment, and with these I returned home, not able to sit down, or rise up, without help, but through the blessings of God I am now better.---

					\hypertarget{EMS-1835-JS-CIR-18-DEC-1835-f}%
				I
					\marginnote{f}%
				have received your letter and purused it with care, I have not entertained a feeling of malice, against you, I am, older than you\sout{r} and have endured, more suffering, having been mar[r]ed by mobs, the labours of my calling, a series of persecution, and inguries, continually heaped upon me, all serve to debilitate, my body, and it may <be> that I cannot boast of being stronger, than you, if I could, or could not, would this be an honor, or dishonor to me,--- if I could boast like David of slaying a Goliath, who defied the armies of the living God, or like Paul, of contending with Peter face to face, with sound arguments, it might be an honor, But to mangle the flesh or seek revenge upon one who never done you any wrong, can not be a source of sweet reflection, to you, nor to me, neither to an honorable father \& mother, brothers, and sisters, and when we reflect, with what care \sout{our} \sout{parents} and with what unremiting diligence our parents, have strove to watch over us, and how many hours, of sorrow, and anxiety, they have spent over our cradles and bedsides, in times of sickness, how careful we ought to be of their feelings in their old age, it cannot be a source of sweet reflection to us to say or do any thing that will bring their grey hairs down with sorrow to the grave,

					\hypertarget{EMS-1835-JS-CIR-18-DEC-1835-g}%
				In
					\marginnote{g}%
				your letter you asked my forgivness, which I readily grant, but it seems to me, that you still retain an idea, that I have given you reasons to be angry or disaffected with me,

				Grant me the privilege of saying then, that however hasty, or harsh, I may have spoken, at any time to you, it has been done for the express purpose of endeavouring, to warn exhort, admonish, and rescue you, from falling into difficulties, and sorrows which I foresaw you plunging into, by giving way to that wicked spirit, which you call your passions, which you should curbe and break down, and put under your feet, which if you do not you, never can be saved, in my view, in the kingdom of God.

				God requires the will of his creatures, to be swallowed up in his will.

				You desire to remain in the church, but forsake your apostleship, this is a stratigem of the evil one, when he has gained one advantage, \sout{your} he lays a plan for another, \sout{by} <but> by maintaining your apostleship in rising up, and making one tremendeous effort, you may overcome your passions, and please God and by forsaking your apostleship, is not to be willing, to make that sacrafice that God requires at your hands and is to incur his displeasure, and without pleasing God do not think, that it will be any better for you, when a man falls one step he must regain that step again, or fall another, he has still more to gain, or eventually all is lost.

					\hypertarget{EMS-1835-JS-CIR-18-DEC-1835-h}%
				I
					\marginnote{h}%
				desire brother William that you will humble yourself, I freely forgive you and you know, my unshaken and \sout{unshaken} unchangable disposition I \sout{think} know in whom I trust, I stand upon the rock, the floods cannot, no they shall not overthrow me, you know the doctrine I teach is true, and you know that God has blessed me, I brought salvation to my fathers house, as an instrument in the hand of God, when they were in a miserable situation, You know that it is my duty to admonish you when you do wrong this liberty I shall always take, and you shall have the same privelege, I take the privelege, to admonish you because of my birthright, and I grant you the privilege because it is my duty, to be humble and to receive rebuke, and instruction, from a brother or a friend.

					\hypertarget{EMS-1835-JS-CIR-18-DEC-1835-i}%
				As
					\marginnote{i}%
				it regards, what course you shall persue hereafter, I do not pretend to say, I leave you in the hands of God and his church. Make your own desision, I will do you good altho you mar me, or slay me, by so doing my garments, shall be clear of your sins, and if at any time you should concider me to be an imposter, for heavens sake leave me in the hands of God, and not think to take vengance on me your self.

				Tyrany ursurpation, and to take mens rights ever has and ever shall be banished from my heart.

				David sought not to kill Saul, although he was guilty of crimes that never entered my heart.

				And now may God have mercy upon my fathers house, may God take away enmity, from betwe[e]n me and thee, and may all blessings be restored, and the past be forgotten forever, may humble repentance bring us both to thee <O God> and to thy power and protection, and a crown, to enjoy the society of father mother Alvin Hyrum Sophron[i]a Samuel Catharine [Katharine] Carloss [Don Carlos] Lucy the Saints and all the sanctif[ie]d in peace forever<, is the prayer of>

				\begin{tightright}
				\sout{This} \sout{from} Your brother

				Joseph Smith Jun
				\end{tightright}

				To William Smith
				
		% -----
		\subsubsection{Record, 19 December 1835}
		% -----
			\label{EMS-1835-JS-19-DEC-1835}
			% \label{EMS-1830-JS-DAY-3LETTERMONTH-YEAR}
			
			\begin{description}
				\item[Print Sources]
			\end{description}

				\begin{tabular}{p{0.5in}p{1in}p{0.5in}p{1in}}
				JSP	 	& J1:135		& JSR 		& --- \\
				DC	 	& --- 				& HC 		& --- \\
				HDDC 	& ---				& TCHC 		& --- \\
				PJS 	& 1:176				& 	 		&  \\
				\end{tabular}
				% HDDC = woodford's DC diss; JSR = Marquardt's JS Revelations; TCHC = Vogel HC
			
			\begin{description}
				\item[Online Access] \href{https://www.josephsmithpapers.org/paper-summary/journal-1835-1836/88}{Here}
				% \item[Analysis] \hyperref[XX]{Here}
			\end{description}
			
			\begin{description}
				\item[Background]
			\end{description}
			
				% It might also be useful to give a two sentence context of the period: e.g., "This speech was given in Nauvoo to a small congregation one month prior to the destruction of the Nauvoo Expositor. These and other tensions may have been on the prophet's mind when he gave the speech."
			
			\begin{description}
				\item[Original Source]
			\end{description}
			
				% brief description of original source material, with reference to JSP or somewhere else for more info
				% Background material: it might be helpful to have a note that says whether a given document was presented as a speech, was written in a journal, etc.
				% Also a comment here about how reliable the source might be, or what shape it is in, if there are large omissions or parts that are hard to understand, and so on.
			
			\begin{description}
				\item[Text]
			\end{description}
			
				...\textbf{I have had many solemn feelings this day Concerning my Brothe[r] William and have prayed in my heart to fervently that the Lord will not \sout{him} <cast him> off but <he> may return to the God of Jacob and magnify his apostleship and calling may this be his happy lot for the Lord of Glorys Sake Amen}
				
		% -----
		\subsubsection{Record, 21 December 1835}
		% -----
			\label{EMS-1835-JS-21-DEC-1835}
			% \label{EMS-1830-JS-DAY-3LETTERMONTH-YEAR}
			
			\begin{description}
				\item[Print Sources]
			\end{description}

				\begin{tabular}{p{0.5in}p{1in}p{0.5in}p{1in}}
				JSP	 	& J1:135		& JSR 		& --- \\
				DC	 	& --- 				& HC 		& --- \\
				HDDC 	& ---				& TCHC 		& --- \\
				PJS 	& 1:176				& 	 		&  \\
				\end{tabular}
				% HDDC = woodford's DC diss; JSR = Marquardt's JS Revelations; TCHC = Vogel HC
			
			\begin{description}
				\item[Online Access] \href{https://www.josephsmithpapers.org/paper-summary/journal-1835-1836/89}{Here}
				% \item[Analysis] \hyperref[XX]{Here}
			\end{description}
			
			\begin{description}
				\item[Background]
			\end{description}
			
				% It might also be useful to give a two sentence context of the period: e.g., "This speech was given in Nauvoo to a small congregation one month prior to the destruction of the Nauvoo Expositor. These and other tensions may have been on the prophet's mind when he gave the speech."
			
			\begin{description}
				\item[Original Source]
			\end{description}
			
				% brief description of original source material, with reference to JSP or somewhere else for more info
				% Background material: it might be helpful to have a note that says whether a given document was presented as a speech, was written in a journal, etc.
				% Also a comment here about how reliable the source might be, or what shape it is in, if there are large omissions or parts that are hard to understand, and so on.
			
			\begin{description}
				\item[Text]
			\end{description}
			
				...\textbf{At home Spent this in indeavering to treasure up know[l]edge for the be[n]ifit of my Calling the} <day> \textbf{pas[s]ed of[f] very pleasantly for which I thank the Lord for his blessings to my soul his great mercy over my Family in sparing\sout{ly} our lives O Continue thy Care over me and mine for Christ sake}
				
		% -----
		\subsubsection{Record, 22 December 1835}
		% -----
			\label{EMS-1835-JS-22-DEC-1835}
			% \label{EMS-1830-JS-DAY-3LETTERMONTH-YEAR}
			
			\begin{description}
				\item[Print Sources]
			\end{description}

				\begin{tabular}{p{0.5in}p{1in}p{0.5in}p{1in}}
				JSP	 	& J1:135		& JSR 		& --- \\
				DC	 	& --- 				& HC 		& --- \\
				HDDC 	& ---				& TCHC 		& --- \\
				PJS 	& 1:176--177				& 	 		&  \\
				\end{tabular}
				% HDDC = woodford's DC diss; JSR = Marquardt's JS Revelations; TCHC = Vogel HC
			
			\begin{description}
				\item[Online Access] \href{https://www.josephsmithpapers.org/paper-summary/journal-1835-1836/89}{Here}
				% \item[Analysis] \hyperref[XX]{Here}
			\end{description}
			
			\begin{description}
				\item[Background]
			\end{description}
			
				% It might also be useful to give a two sentence context of the period: e.g., "This speech was given in Nauvoo to a small congregation one month prior to the destruction of the Nauvoo Expositor. These and other tensions may have been on the prophet's mind when he gave the speech."
			
			\begin{description}
				\item[Original Source]
			\end{description}
			
				% brief description of original source material, with reference to JSP or somewhere else for more info
				% Background material: it might be helpful to have a note that says whether a given document was presented as a speech, was written in a journal, etc.
				% Also a comment here about how reliable the source might be, or what shape it is in, if there are large omissions or parts that are hard to understand, and so on.
			
			\begin{description}
				\item[Text]
			\end{description}
			
				\textbf{At home \sout{this} Continued my studys O may God give me learning even Language and indo [endue] me with qualifycations to magnify his name while I live I also deliv[er]ed an address to the Church this Evening the Lord blessed my Soul, my scribe also is unwell O my God heal him and for his kindness to me O my soul be thou greatful to him and bless him and he shall be blessed \sout{of} \sout{for} \sout{ever} of God forever I believe him to be a faithful friend to me therefore my soul delighteth in him Amen}
				
		% -----
		\subsubsection{Revelation, 26 December 1835 [D\&C 108]}
		% -----
			\label{EMS-1835-JS-26-DEC-1835}
			% \label{EMS-1830-JS-DAY-3LETTERMONTH-YEAR}
			
			\begin{description}
				\item[Print Sources]
			\end{description}

				\begin{tabular}{p{0.5in}p{1in}p{0.5in}p{1in}}
				JSP	 	& D5:122--123; J1:137--138		& JSR 		& 277--278 \\
				DC	 	& Section 108 		& HC 		& 2:345 \\
				HDDC 	& 1433--1439		& TCHC 		& 2:338--339 \\
				PJS 	& 1:178			& 	 		&  \\
				\end{tabular}
				% HDDC = woodford's DC diss; JSR = Marquardt's JS Revelations; TCHC = Vogel HC
			
			\begin{description}
				\item[Online Access] \href{https://www.josephsmithpapers.org/paper-summary/revelation-26-december-1835-dc-108/1}{Here}
				% \item[Analysis] \hyperref[DC108]{Here}
			\end{description}
			
			\begin{description}
				\item[Background]
			\end{description}
			
				% It might also be useful to give a two sentence context of the period: e.g., "This speech was given in Nauvoo to a small congregation one month prior to the destruction of the Nauvoo Expositor. These and other tensions may have been on the prophet's mind when he gave the speech."
			
			\begin{description}
				\item[Original Source]
			\end{description}
			
				% brief description of original source material, with reference to JSP or somewhere else for more info
				% Background material: it might be helpful to have a note that says whether a given document was presented as a speech, was written in a journal, etc.
				% Also a comment here about how reliable the source might be, or what shape it is in, if there are large omissions or parts that are hard to understand, and so on.
			
			\begin{description}
				\item[Text]
			\end{description}
			
				The following is a revelation given to Lyman Sherman this day 26 Dec 1835
				
				Verily thus saith the Lord unto you my servant Lyman your sins are forgiven you because you have obeyed my voice in coming up hither this morning to receive councel of him whom I have appointed
				
				Therefore let your soul be at rest concerning your spiritual standing, and resist no more my voice, and arise up, and be more careful henceforth in observing your vows which you have made and do make, and you shall be blessed with exceding great blessings. Wait patiently untill the time when the solemn assembly shall be called of my servants then you shall be numbered with the first of mine elders and receive right by ordination with the rest of mine elders whom I have chosen
				
				Behold this is the promise of the father unto you if you continue faithful---
				
				and it shall be fulfilled upon you in that day that you shall have right to preach my gospel wheresoever I shall send you from henceforth from that time, Therefore strengthen your brethren in all your conversation in all your prayers, and in all your exhortations, and in all your doings, and behold and lo I am with you to bless you and deliver you forever Amen

% -----------------------------------------------------------
\pagebreak{}
% -----------------------------------------------------------